%%%%%%%%%%%%%%%%%%%%%%%%%%%%%%%%%%%%%%%%%%%%%%%%%%%%%%%%%%%%%%%%%%%%%%
% Lớp 11 | Giới hạn. Hàm số liên tục | Giới hạn của hàm số
% Dạng 1D3?2-3: PP đặt thừa số chung, kết quả hữu hạn
% Số câu: 244  (sinh tự động bởi tach_cau.py; từ nay thầy sửa trực tiếp file này)
%%%%%%%%%%%%%%%%%%%%%%%%%%%%%%%%%%%%%%%%%%%%%%%%%%%%%%%%%%%%%%%%%%%%%%

%%%%%%%%%%%%%%%% TRẮC NGHIỆM 4 PHƯƠNG ÁN %%%%%%%%%%%%%%%%

% nguồn: Ok_Toan11_TracNghiem_Lan2
\begin{ex}%[1D3N2-3]%[ma:0898c45a]
	Hãy xác  định xem kết quả nào \textbf{sai}?
	 \choice
		{$\lim\limits_{x\to  1} \dfrac{x+1}{x}=2$}
		{$\lim\limits_{x\to+\infty} \dfrac{x+2}{x-4}=1$} 
		{$\lim\limits_{x\to  1} \dfrac{x^2-3x+2}{x-1}=-1$}
		{\True$\lim\limits_{x\to 4}  \dfrac{x^2-16}{x^2+x-20}=\dfrac{9}{8}$}
	 \loigiai{
	Ta có $\lim\limits_{x\to 4} \dfrac{x^2-16}{x^2+x-20}=\lim\limits_{x\to 4} \dfrac{\left(x-4\right)\left(x+4\right)}{\left(x-4\right)\left(x+5\right)}=\lim\limits_{x\to  4} \dfrac{x+4}{x+5}=\dfrac{8}{9}$.}
\end{ex}

% nguồn: Ok_Toan11_TracNghiem_Lan2
\begin{ex}%[1D3N2-3]%[ma:fa5cf29c]
	Tìm giới  hạn $L=\lim\limits_{x\to \dfrac{\pi}{2}} \dfrac{\cos x}{x-\frac{\pi}{2}}$.
	\choice 
		{$L=1$} 
		{\True$L=-1$} 
		 {$L=0$}
		{$L=\dfrac{\pi}{2}$} 
	\loigiai{ 
		Đặt  $t=x-\dfrac{\pi}{2}$.\\
		Khi $x\to  \dfrac{\pi}{2}$ thì $t\to 0$.\\
		Vậy $L=\lim\limits_{t\to 0} \dfrac{\cos \left(t+\dfrac{\pi}{2} \right)}{t}=\lim\limits_{t\to 0} \dfrac{-\sin{t}}{t}=-1$. 
	 }
\end{ex}

% nguồn: Ok_Toan11_TracNghiem_Lan2
\begin{ex}%[1D3N2-3]%[ma:578a4562]
	Giới hạn $\lim\limits_{x\to -2}\dfrac{2x^2+3x-2}{x^2-4}$  bằng bao nhiêu?
	\choice 
		 {\True $\dfrac{5}{4}$}
		{$-\dfrac{5}{4}$} 
		 {$\dfrac{1}{4}$}
		 {$2$}
	\loigiai{ 
	 Ta có $\lim\limits_{x\to -2}\dfrac{2x^2+3x-2}{x^2-4}=\lim\limits_{x\to -2}\dfrac{\left( 2x-1 \right)\left( x+2 \right)}{\left( x-2 \right)\left( x+2 \right)}=\lim\limits_{x\to -2}\dfrac{2x-1}{x-2}=\dfrac{5}{4}$.}
\end{ex}

% nguồn: Ok_Toan11_TracNghiem_Lan2
\begin{ex}%[1D3N2-3]%[ma:e4bd573b]
	$\lim\limits_{x\to 1}\dfrac{x^2-1}{x-1}$ có giá trị  bằng
	\choice 
		 {$1$}
		{$+\infty  $}
		 {$0$}
		{\True $2$} 
	 \loigiai{
	Ta có $\lim\limits_{x\to 1}\dfrac{x^2-1}{x-1}=\lim\limits_{x\to 1}\dfrac{\left( x-1 \right)\left( x+1 \right)}{x-1}=\lim\limits_{x\to  1}\left( x+1 \right)=2$.}
\end{ex}

% nguồn: Ok_Toan11_TracNghiem_Lan2
\begin{ex}%[1D3N2-3]%[ma:7ecf060e]
	Giá trị của  $I=\lim\limits_{x\to -\sqrt{2}}\dfrac{x+\sqrt{2}}{x^2-2}$ bằng
	\choice 
		 {$2$}
		 {\True $\dfrac{-1}{2\sqrt{2}}$}
		{$1$} 
		 {$\sqrt{2}$}
	\loigiai{ 
	$I=\lim\limits_{x\to -\sqrt{2}}\dfrac{x+\sqrt{2}}{x^2-2}=\lim\limits_{x\to  -\sqrt{2}}\dfrac{x+\sqrt{2}}{\left( x+\sqrt{2} \right)\left( x-\sqrt{2} \right)}=\lim\limits_{x\to -\sqrt{2}}\dfrac{1}{x-\sqrt{2}}=\dfrac{-1}{2\sqrt{2}}$.}
\end{ex}

% nguồn: Ok_Toan11_TracNghiem_Lan2
\begin{ex}%[1D3N2-3]%[ma:2b315529]
	Giới hạn $\lim\limits_{x\to-\infty} \dfrac{2x-1}{x+2}$  bằng
	\choice 
		 {$1$}
		 {$-\dfrac{1}{2}$}
		{\True $2$} 
		{$-\infty$} 
	 \loigiai{
		Ta có $\lim\limits_{x\to-\infty} \dfrac{2x-1}{x+2}=\lim\limits_{x\to-\infty}  \dfrac{2-\dfrac{1}{x}}{1+\dfrac{2}{x}}=2$.
	} 
\end{ex}

% nguồn: Ok_Toan11_TracNghiem_Lan2
\begin{ex}%[1D3N2-3]%[ma:b572f22d]
	Giá trị của $\lim\limits_{x \to 0} \dfrac{\sqrt{x^3+x^2+1}-1}{x^2}$ bằng 
	 \choice
		{$1$} 
		{\True  $\dfrac{1}{2}$}
		 {$-1$}
		 {$0$}
	 \loigiai{
	$\lim\limits_{x \to  0} \dfrac{\sqrt{x^3+x^2+1}-1}{x^2}=\lim\limits_{x \to 0} \dfrac{x^3+x^2+1-1}{x^2\left(\sqrt{x^3+x^2+1}+1\right)}=\lim\limits_{x \to 0} \dfrac{x+1}{\left(\sqrt{x^3+x^2+1}+1\right)}=\dfrac{1}{2}$.}
\end{ex}

% nguồn: Ok_Toan11_TracNghiem_Lan2
\begin{ex}%[1D3N2-3]%[ma:c3b64c4e]
	Giới hạn $\lim\limits_{x\to  -1}\dfrac{3x^2-2x-5}{x^2-1}$ bằng
	\choice 
		 {$-3$}
		 {$+\infty $}
		 {$0$}
		{\True  $4$}
	 \loigiai{
	Ta có $\lim\limits_{x\to -1}\dfrac{3x^2-2x-5}{x^2-1}=\lim\limits_{x\to -1}\dfrac{\left( 3x-5  \right)\left( x+1 \right)}{\left( x-1 \right)\left( x+1 \right)}=\lim\limits_{x\to -1}\dfrac{3x-5}{x-1}=\dfrac{3\cdot\left( -1 \right)-5}{\left( -1 \right)-1}=\dfrac{-8}{-2}=4$.}
\end{ex}

% nguồn: Ok_Toan11_TracNghiem_Lan2
\begin{ex}%[1D3N2-3]%[ma:117909bb]
	Tìm  $\lim\limits_{x \to 2} \dfrac{x^2-5x+6}{\sqrt{4x+1}-3}$ là
	\choice 
		 {$\dfrac{3}{2}$}
		 {$-\dfrac{2}{3}$}
		 {\True$-\dfrac{3}{2}$}
		 {$\dfrac{1}{2}$}
	\loigiai{ 
	$\lim\limits_{x \to 2}  \dfrac{x^2-5x+6}{\sqrt{4x+1}-3}=\lim\limits_{x \to 2} \dfrac{\left(x-2\right)\left(x-3\right)\left(\sqrt{4x+1}+3\right)}{4\left(x-2\right)}=\lim\limits_{x \to 2} \dfrac{\left(x-3\right)\left(\sqrt{4x+1}+3\right)}{4}=-\dfrac{3}{2}$.}
\end{ex}

% nguồn: Ok_Toan11_TracNghiem_Lan2
\begin{ex}%[1D3N2-3]%[ma:cf3f6e19]
	Tính $\lim\limits_{x\to 0}  \dfrac{\sqrt[3]{8+x^2}-2}{x^2}$.
	 \choice
		{\True $\dfrac{1}{12}$} 
		{$\dfrac{1}{4}$} 
		 {$\dfrac{1}{3}$}
		{$\dfrac{1}{6}$} 
	 \loigiai{
		Ta  có
		\begin{align*}
			\lim\limits_{x\to 0} \dfrac{\sqrt[3]{8+x^2}-2}{x^2}  =&\lim\limits_{x\to 0} \dfrac{8+x^2-8}{x^2\left(\sqrt[3]{\left(8+x^2\right)^2}+2\sqrt[3]{8+x^2}+4\right)}\\
			=&\lim\limits_{x\to 0} \dfrac{1}{\sqrt[3]{\left(8+x^2\right)^2}+2\sqrt[3]{8+x^2}+4}=\dfrac{1}{12}. 
		\end{align*}
	} 
\end{ex}

% nguồn: Ok_Toan11_TracNghiem_Lan2
\begin{ex}%[1D3N2-3]%[ma:9502dcda]
	Tính giới hạn $L=\lim\limits_{x\to  1}\dfrac{x^2+3x-4}{x-1}$
	 \choice
		 {$L=-5$}
		{$L=0$} 
		{$L=-3$} 
		 {\True $L=5$}
	 \loigiai{
	Ta có $L=\lim\limits_{x\to 1}\dfrac{x^2+3x-4}{x-1}=\lim\limits_{x\to  1}\dfrac{\left( x-1 \right)\left( x+4 \right)}{x-1}=\lim\limits_{x\to 1}\left( x+4 \right)=5$.}
\end{ex}

% nguồn: Ok_Toan11_TracNghiem_Lan2
\begin{ex}%[1D3N2-3]%[ma:13f3a5ac]
	Giới hạn $\lim\limits_{x\to  0} \dfrac{\sqrt{x^2-3x+4}-2}{x}$ bằng
	\choice 
		 {$-\dfrac{1}{2}$}
		{$\dfrac{1}{2}$} 
		{\True  $-\dfrac{3}{4}$}
		 {$-\dfrac{2}{3}$}
	\loigiai{ 
		$\lim\limits_{x\to 0} \dfrac{\sqrt{x^2-3x+4}-2}{x}=\lim\limits_{x\to  0} \dfrac{x^2-3x+4-4}{x\left(\sqrt{x^2-3x+4}+2\right)}=\lim\limits_{x\to 0} \dfrac{x-3}{\sqrt{x^2-3x+4}+2}=-\dfrac{3}{4}.$
	} 
\end{ex}

% nguồn: Ok_Toan11_TracNghiem_Lan2
\begin{ex}%[1D3N2-3]%[ma:86d13dcb]
	Giới hạn $\lim\limits_{x\to  2}\dfrac{x-2}{x^2-4}$ bằng
	\choice 
		{$2$} 
		{$4$} 
		{\True$\dfrac{1}{4}$} 
		 {$0$}
	 \loigiai{
		$\lim\limits_{x\to 2} \dfrac{x-2}{x^2-4}=\lim\limits_{x\to  2} \dfrac{x-2}{(x-2)(x+2)}=\lim\limits_{x\to 2} \dfrac{1}{x+2}=\dfrac{1}{4}.$
	 }
\end{ex}

% nguồn: Ok_Toan11_TracNghiem_Lan2
\begin{ex}%[1D3N2-3]%[ma:4d321f2d]
	Giới  hạn $\lim\limits_{x\to-\infty} \dfrac{4x+1}{-x+1}$ bằng
	\choice 
		{$2$} 
		 {$4$}
		 {$-1$}
		{\True  $-4$}
	\loigiai{ 
		Ta có $\lim\limits_{x\to-\infty} \dfrac{4x+1}{-x+1}=\lim\limits_{x\to-\infty}  \dfrac{4+\dfrac{1}{x}}{-1+\dfrac{1}{x}}=-4$.
	 }
\end{ex}

% nguồn: Ok_Toan11_TracNghiem_Lan2
\begin{ex}%[1D3N2-3]%[ma:bafc6255]
	$\lim\limits_{x\to  4}\dfrac{x^2-3x-4}{x-4}$ bằng
	 \choice
		{Không tồn  tại}
		{$0$} 
		 {\True $5$}
		{$4$} 
	 \loigiai{
	Ta có $\lim\limits_{x\to 4}\dfrac{x^2-3x-4}{x-4}=\lim\limits_{x\to 4}\dfrac{(x+1)(x-4)}{x-4}=\lim\limits_{x\to 4}(x+1)=5$.} 
\end{ex}

% nguồn: Ok_TracNghiemLop11_Lan1
\begin{ex}%[1D3N2-3]%[ma:d66c80e8]
	Giới hạn  của $\displaystyle\lim\limits_{x \rightarrow 3} \dfrac{x^2-9}{x-3}$ là
	\choice 
		 {$3$}
		 {$0$}
		 {\True $6$}
		{$-6$} 
	 \loigiai{
		Ta có  $\displaystyle\lim\limits_{x \rightarrow 3} \dfrac{x^2-9}{x-3} = \displaystyle\lim\limits_{x \rightarrow 3} \dfrac{\left(x-3\right)\left(x+3\right)}{x-3} = \displaystyle\lim\limits_{x \rightarrow 3} \left(x+3\right) = 6$.
	 }
\end{ex}

% nguồn: Ok_TracNghiemLop11_Lan1
\begin{ex}%[1D3N2-3]%[ma:090b7e3d]
	Giá  trị của $\lim\limits_{x\to -2}\dfrac{4x^3 -1}{3x^2 +x +2}$ có kết quả là
	\choice{\True $  -\dfrac{11}{4}$}
		{$-\infty  $}
		{$\dfrac{11}{4} $} 
		{ $ +\infty $} 
	 \loigiai{
	Ta có $\lim\limits_{x\to -2}\dfrac{4x^3 -1}{3x^2 +x +2}=\dfrac{4\cdot(-2)^3 -1}{3\cdot(-2)^2  +-2 +2}=\dfrac{-11}{4}$.}
\end{ex}

% nguồn: Ok_TracNghiemLop11_Lan1
\begin{ex}%[1D3N2-3]%[ma:16ed1c00]
	Giới hạn  của $\displaystyle\lim\limits_{x \rightarrow 3} \dfrac{x-9}{x-2}$ là
	\choice 
		{$-3$} 
		{$0$} 
		 {\True $-6$}
		{$6$} 
	 \loigiai{
		Ta có $\displaystyle\lim\limits_{x \rightarrow 3} \dfrac{x-9}{x-2} =  \dfrac{3-9}{3-2}=-6$.
	 }
\end{ex}

% nguồn: Ok_Toan11_TracNghiem_Lan2
\begin{ex}%[1D3H2-3]%[ma:2278b91e]
	Giới hạn của $\lim\limits_{x  \to 0} \dfrac{x^4-4x^3}{5x^3}$ bằng
	\choice 
		{$0$} 
		 {$-4$}
		{$\dfrac{-3}{5}$} 
		{\True $\dfrac{-4}{5}$} 
	 \loigiai{
		 Ta có $\lim\limits_{x \to 0} \dfrac{x^4-4x^3}{5x^3}=\lim\limits_{x \to 0} \dfrac{x^3(x-4)}{5x^3}=\lim\limits_{x \to 0} \dfrac{x-4}{5}=\dfrac{-4}{5}$.
	} 
\end{ex}

% nguồn: Ok_Toan11_TracNghiem_Lan2
\begin{ex}%[1D3H2-3]%[ma:c86c8384]
	Biết $\lim\limits_{x\to 1}\dfrac{(x+1)(x^2-2)}{x^2+x+3}=\dfrac{a}{b}$,  $a, b\in \mathbb{Z}$; $b<0$. Biết $\dfrac{a}{b}$ tối giản thì $a+b$ bằng bao nhiêu?
	\choice 
		{\True $-3$} 
		{$-6$} 
		 {$3$}
		 {$1$}
	\loigiai{ 
		Ta có $\lim\limits_{x\to  1}\dfrac{(x+1)(x^2-2)}{x^2+x+3}
		 =\dfrac{2}{-5}\Rightarrow \heva{&a=2\\&b=-5}$ nên $a+b=2+(-5)=-3$.
	} 
\end{ex}

% nguồn: Ok_Toan11_TracNghiem_Lan2
\begin{ex}%[1D3H2-3]%[ma:80fe8281]
	Giá trị $\displaystyle\lim\limits _{x \to-\infty}\dfrac{3  x-1}{x+5}$ bằng
	\choice 
		{\True $3$} 
		 {$-3$}
		 {$-\dfrac{1}{5}$}
		 {$5$}
	 \loigiai{
	Ta có $\displaystyle\lim\limits _{x  \to-\infty}\dfrac{3 x-1}{x+5}=\displaystyle\lim\limits _{x \to-\infty}\dfrac{3-\dfrac{1}{x}}{1+\dfrac{5}{x}}=3$.}
\end{ex}

% nguồn: Ok_Toan11_TracNghiem_Lan2
\begin{ex}%[1D3H2-3]%[ma:3fe75230]
	Cho $a$ thỏa mãn $\lim\limits_{x\to 1}\dfrac{x^2-(a^2-1)x+a^2-2}{x-1}=-1$. Khẳng  định nào sau đây đúng?
	 \choice
		{$a  <-2$}
		{\True $a  < 3$}
		{$a  > 1$}
		{$a > 0$} 
	 \loigiai{Ta có
		\begin{align*}
			 &\lim\limits_{x\to 1}\dfrac{x^2-(a^2-1)x+a^2-2}{x-1}=-1\\
			\Leftrightarrow&\lim\limits_{x\to  1}\dfrac{x(x-1)-(a^2-2)(x-1)}{x-1}=-1\\
			\Leftrightarrow&\lim\limits_{x\to  1}(x-a^2+2)=-1\\
			\Leftrightarrow  &1-a^2+2=-1\\
			 \Leftrightarrow &a^2=4\\
			\Leftrightarrow  &a=\pm 2.
		\end{align*}
		 Vậy $a < 3$.
	} 
\end{ex}

% nguồn: Ok_Toan11_TracNghiem_Lan2
\begin{ex}%[1D3H2-3]%[ma:ed0fc898]
	Tính $\lim\limits_{x\to  3} \dfrac{\sqrt{x^2+x}-2\sqrt{3}}{x-3}$?
	\choice 
		{$\dfrac{\sqrt{3}}{12}$} 
		{$-\dfrac{\sqrt{3}}{12}$} 
		{\True  $\dfrac{7\sqrt{3}}{12}$}
		 {$-\dfrac{7\sqrt{3}}{12}$}
	 \loigiai{
		 Ta có
		\begin{eqnarray*}
			 & &\lim\limits_{x\to 3} \dfrac{\sqrt{x^2+x}-2\sqrt{3}}{x-3}\\
			&=&\lim\limits_{x\to 3} \dfrac{x^2+x-12}{\left(x-3\right)\left(\sqrt{x^2+x}+2\sqrt{3}  \right)}\\
			&=&\lim\limits_{x\to  3} \dfrac{\left(x-3\right)\left(x+4\right)}{\left(x-3\right)\left(\sqrt{x^2+x}+2\sqrt{3} \right)}\\
			&=&\lim\limits_{x\to 3} \dfrac{x+4}{\sqrt{x^2+x}+2\sqrt{3}}\\ 
			&=&\dfrac{3+4}{\sqrt{3^2+3}+2\sqrt{3}}\\ 
			&=&\dfrac{7}{4\sqrt{3}}=\dfrac{7\sqrt{3}}{12}. 
		\end{eqnarray*}
	} 
\end{ex}

% nguồn: Ok_Toan11_TracNghiem_Lan2
\begin{ex}%[1D3H2-3]%[ma:93415856]
	 Kết quả của ${\lim\limits_{x \to -2} \dfrac{x^2+3x+2}{x+2}}$ là
	\choice 
		 {\True $-1$}
		 {$3$}
		 {$2$}
		 {$1$}
	\loigiai{ 
		Ta có $\lim\limits_{x \to -2}  \dfrac{x^2+3x+2}{x+2}=\lim\limits_{x \to -2}(x+1)=-1$.
	} 
\end{ex}

% nguồn: Ok_Toan11_TracNghiem_Lan2
\begin{ex}%[1D3H2-3]%[ma:ee3c7670]
	Tính giới hạn $\lim\limits_{x \to  2} \dfrac{x^2-4}{x-2}$ ta được kết quả là
	 \choice
		 {$+\infty$}
		 {$2$}
		{$0$} 
		{\True $4$} 
	 \loigiai{
		Ta có $\lim\limits_{x  \to 2} \dfrac{x^2-4}{x-2}=\lim\limits_{x \to 2}(x+2)=4$.
	} 
\end{ex}

% nguồn: Ok_Toan11_TracNghiem_Lan2
\begin{ex}%[1D3H2-3]%[ma:02e1dc95]
	Giá  trị của $\lim\limits_{x \to -1} \dfrac{x^2-2x-3}{x^2-1}$ là
	 \choice
		{$1$} 
		 {\True $2$}
		{$-\infty$} 
		{$0$} 
	\loigiai{ 
		Ta có $\lim\limits_{x \to -1} \dfrac{x^2-2x-3}{x^2-1}=\lim\limits_{x \to -1} \dfrac{(x+1)(x-3)}{(x+1)(x-1)}=\lim\limits_{x \to -1}  \dfrac{x-3}{x-1}=\dfrac{-1-3}{-1-1}=2$.
	} 
\end{ex}

% nguồn: Ok_Toan11_TracNghiem_Lan2
\begin{ex}%[1D3H2-3]%[ma:e9de2588]
	$\lim\limits_{x\to -\infty}\dfrac{1}{x^2+9}$  bằng
	 \choice
		{$-\infty$} 
		 {1}
		 {\True 0}
		{$+\infty$} 
	 \loigiai{
	Ta có $\heva{&\lim\limits_{x\to -\infty}1=1 \\&\lim\limits_{x\to  -\infty}(x^2+9)=+\infty}\Rightarrow \lim\limits_{x\to -\infty}\dfrac{1}{x^2+9}=0$.}
\end{ex}

% nguồn: Ok_Toan11_TracNghiem_Lan2
\begin{ex}%[1D3H2-3]%[ma:424de8f4]
	Giới hạn $\displaystyle\lim\limits _{x  \to+\infty}\dfrac{c x^{2}+a}{x^{2}+b}$ bằng
	 \choice
		 {$a$}
		{$b$} 
		{\True $c$} 
		{$\dfrac{a+b}{c}$} 
	\loigiai{ 
	Ta có $\displaystyle\lim\limits _{x \to+\infty}\dfrac{c x^{2}+a}{x^{2}+b}=\displaystyle\lim\limits _{x \to+\infty}\dfrac{c+\dfrac{a}{x^{2}}}{1+\dfrac{b}{x^{2}}}=\dfrac{c+0}{1+0}=c$.} 
\end{ex}

% nguồn: Ok_Toan11_TracNghiem_Lan2
\begin{ex}%[1D3H2-3]%[ma:903ad484]
	Giá trị $\displaystyle\lim\limits _{x \to-\infty}\dfrac{(x-1)(x+2)}{x^{2}+9}$ bằng 
	 \choice
		{$\dfrac{2}{9}$} 
		{\True $1$} 
		{$-1$} 
		{$-\dfrac{1}{9}$} 
	\loigiai{ 
	Ta có $\displaystyle\lim\limits  _{x \to-\infty}\dfrac{(x-1)(x+2)}{x^{2}+9}=\displaystyle\lim\limits _{x \to-\infty}\dfrac{\left(1-\dfrac{1}{x}\right)\left(1+\dfrac{2}{x}\right)}{1+\dfrac{9}{x^{2}}}=1$.}
\end{ex}

% nguồn: Ok_Toan11_TracNghiem_Lan2
\begin{ex}%[1D3H2-3]%[ma:007675d8]
	Tính giới  hạn$\lim\limits_{x\to -\infty}\dfrac{\sqrt{4x^2+x+1}-\sqrt{x^2-x+3}}{3x+2}$.
	 \choice
		 {\True $-\dfrac{1}{3}$}
		 {$\dfrac{2}{3}$}
		 {$\dfrac{1}{3}$}
		{$-\dfrac{2}{3}$} 
	\loigiai{ 
		\begin{eqnarray*}
			 &&\lim\limits_{x\to -\infty}\dfrac{\sqrt{4x^2+x+1}-\sqrt{x^2-x+3}}{3x+2}\\
			&=&\lim\limits_{x\to  -\infty}\dfrac{-x\sqrt{4+\dfrac{1}{x}+\dfrac{1}{x^2}}+x\sqrt{1-\dfrac{1}{x}+\dfrac{3}{x^2}}}{3x+2}\\
			 &=&\lim\limits_{x\to -\infty}\dfrac{-\sqrt{4+\dfrac{1}{x}+\dfrac{1}{x^2}}+\sqrt{1-\dfrac{1}{x}+\dfrac{3}{x^2}}}{3+\dfrac{2}{x}}\\
			 &=&-\dfrac{1}{3}.
	\end{eqnarray*}}
\end{ex}

% nguồn: Ok_Toan11_TracNghiem_Lan2
\begin{ex}%[1D3H2-3]%[ma:87f0757a]
	${\lim\limits_{x\to -\infty}\dfrac{x-\sqrt{x^2+x}}{x+1}}$  bằng
	\choice 
		{$1$} 
		{\True  $2$}
		{$0$} 
		 {$-\infty $}
	 \loigiai{
	Ta có  $\lim\limits_{x\to -\infty}\dfrac{x-\sqrt{x^2+x}}{x+1}=\lim\limits_{x\to -\infty}\dfrac{x+x\sqrt{1+\dfrac{1}{x}}}{x+1}=\lim\limits_{x\to -\infty}\dfrac{1+\sqrt{1+\dfrac{1}{x}}}{1+\dfrac{1}{x}}=2$.}
\end{ex}

% nguồn: Ok_Toan11_TracNghiem_Lan2
\begin{ex}%[1D3H2-3]%[ma:9d96c719]
	Biết rằng $\lim\limits_{x \to \sqrt{3}} \dfrac{2x^3+6 \sqrt{3}}{3-x^2}=a\sqrt{3}+b$, với $a$, $b$ là các số nguyên. Tính  $a^2+b^2$.
	 \choice
		 {\True $9$}
		{$25$} 
		 {$5$}
		 {$13$}
	\loigiai{ 
		Ta  có
		\begin{eqnarray*}
			\lim\limits_{x \to -\sqrt{3}}  \dfrac{2x^3+6 \sqrt{3}}{3-x^2}&=&\lim\limits_{x \to -\sqrt{3}} \dfrac{2\left(x+\sqrt{3}\right)\left(x^2-\sqrt{3}x+3\right)}{\left(\sqrt{3}-x\right)\left(\sqrt{3}+x\right)}\\
			&=&\lim\limits_{x \to -\sqrt{3}}  \dfrac{2\left(x^2-\sqrt{3}x+3\right)}{\sqrt{3}-x}\\
			&=&\dfrac{2\left[\left(-\sqrt{3}\right)^2-\sqrt{3}  \cdot\left(-\sqrt{3}\right)+3\right]}{\sqrt{3}-\left(-\sqrt{3}\right)}=\dfrac{18}{2 \sqrt{3}}\\
			&=&3\sqrt{3}. 
		\end{eqnarray*}
		 Suy ra $a=3$; $b=0$.\\
		 Vậy $a^2+b^2=9$.
	 }
\end{ex}

% nguồn: Ok_Toan11_TracNghiem_Lan2
\begin{ex}%[1D3H2-3]%[ma:8160965f]
	Kết quả của $\lim\limits_{x \to 3}  \dfrac{9-x^2}{x^2-4x+3}$ là
	\choice 
		{\True $-3$} 
		{$4$} 
		 {$-4$}
		 {$3$}
	\loigiai{ 
		Ta có  $\lim\limits_{x \to 3} \dfrac{9-x^2}{x^2-4x+3}=\lim\limits_{x \to 3} \dfrac{(3-x)(3+x)}{(x-3)(x-1)}=\lim\limits_{x \to 3} \dfrac{-x-3}{x-1}=-3$.
	 }
\end{ex}

% nguồn: Ok_Toan11_TracNghiem_Lan2
\begin{ex}%[1D3H2-3]%[ma:4482b828]
	Tính  giới hạn $L=\lim\limits_{x \to 1} \dfrac{2x^2-x-1}{x-1}$.
	\choice 
		{$L=\dfrac{3}{2}$} 
		 {\True $L=3$}
		 {$L=1$}
		 {$L=-\dfrac{1}{2}$}
	 \loigiai{
		Ta có $L=\lim\limits_{x \to 1} \dfrac{2x^2-x-1}{x-1}=\lim\limits_{x \to 1}  \dfrac{(x-1)(2x+1)}{x-1}=\lim\limits_{x \to 1}(2x+1)=3$.
	} 
\end{ex}

% nguồn: Ok_Toan11_TracNghiem_Lan2
\begin{ex}%[1D3H2-3]%[ma:dcff8307]
	Tính $\lim\limits_{x\to 1^{+}}  \dfrac{x^2-3x+2}{6\sqrt{x+8}-x-17}$.
	 \choice
		{$-\infty$} 
		{$0$} 
		{\True $+\infty$} 
		 {$\dfrac{1}{6}$}
	 \loigiai{
		Ta có \begin{align*}
		\lim\limits_{x\to 1^{+}} \dfrac{x^2-3x+2}{6\sqrt{x+8}-x-17}=&\lim\limits_{x\to 1^{+}} \dfrac{\left(x-1\right)\left(x-2\right)\left(6\sqrt{x+8}+x+17\right)}{-\left(x-1\right)^2}\\=&\lim\limits_{x\to  1^{+}} \dfrac{\left(x-2\right)\left(6\sqrt{x+8}+x+17\right)}{-\left(x-1\right)}.
	\end{align*}
	 Mà
	\begin{itemize}
		\item $\lim\limits_{x\to  1^{+}} \left(x-2\right)\left(6\sqrt{x+8}+x+17\right)=-36$.
		\item $\lim\limits_{x\to 1^{+}}  \left(-x+1\right)=0$ và $-x+1< 0$.
	\end{itemize}
	$\Rightarrow \lim\limits_{x\to 1^{+}} \dfrac{x^2-3x+2}{6\sqrt{x+8}-x-17}=+\infty$. 
 }
\end{ex}

% nguồn: Ok_Toan11_TracNghiem_Lan2
\begin{ex}%[1D3H2-3]%[ma:723e8db6]
	Tính  giới hạn $\displaystyle\lim\limits _{x \to+\infty}\dfrac{x^{2}-2\,018 x+3}{2 x^{2}+2\,018 x}$ được
	\choice 
		{$2\,018$} 
		 {\True $\dfrac{1}{2}$}
		 {$2$}
		 {$\dfrac{1}{2\,018}$}
	 \loigiai{
	Ta có $\displaystyle\lim\limits _{x \to+\infty}\dfrac{x^{2}-2\,018 x+3}{2  x^{2}+2\,018 x}=\displaystyle\lim\limits _{x \to+\infty}\dfrac{1-\dfrac{2\,018}{x}+\dfrac{3}{x^{2}}}{2+\dfrac{2\,018}{x}}=\dfrac{1}{2}$.}
\end{ex}

% nguồn: Ok_Toan11_TracNghiem_Lan2
\begin{ex}%[1D3H2-3]%[ma:d4ae7b5a]
	 Tính $L=\displaystyle\lim\limits _{x \to-\infty}\dfrac{2 x+1}{x+1}$.
	\choice 
		 {$L=-2$}
		{$L=-1$} 
		{$L=-\dfrac{1}{2}$} 
		{\True  $L=2$}
	 \loigiai{
	Ta có $L=\displaystyle\lim\limits _{x \to-\infty}\dfrac{2 x+1}{x+1}=\displaystyle\lim\limits _{x \to-\infty}\dfrac{x\left(2+\dfrac{1}{x}\right)}{x\left(1+\dfrac{1}{x}\right)}=\displaystyle\lim\limits _{x  \to-\infty}\dfrac{2+\dfrac{1}{x}}{1+\dfrac{1}{x}}=\dfrac{2+0}{1+0}=2$.}
\end{ex}

% nguồn: Ok_Toan11_TracNghiem_Lan2
\begin{ex}%[1D3H2-3]%[ma:b5e8696c]
	Kết quả của $\lim\limits_{x \to  4} \dfrac{x^2-16}{x-4}$ là
	 \choice
		 {$5$}
		{$6$} 
		{$7$} 
		{\True  $8$}
	 \loigiai{
		Ta có $\lim\limits_{x \to 4} \dfrac{x^2-16}{x-4}=\lim\limits_{x \to  4}(x+4)=8$.
	} 
\end{ex}

% nguồn: Ok_Toan11_TracNghiem_Lan2
\begin{ex}%[1D3H2-3]%[ma:34b9796a]
	Tính $\lim\limits_{x\to 2}\left(\dfrac{1}{x^2-3x+2}+\dfrac{1}{x^2-5x+6}\right)$. 
	\choice 
		 {$2$}
		 {$+\infty$}
		{\True  $-2$}
		{$0$} 
	\loigiai{Ta  có
		\begin{align*}
			&\lim\limits_{x\to  2}\left(\dfrac{1}{x^2-3x+2}+\dfrac{1}{x^2-5x+6}\right)\\
			=&\lim\limits_{x\to  2}\dfrac{2x^2-8x+8}{(x^2-3x+2)(x^2-5x+6)}\\
			=&\lim\limits_{x\to  2}\dfrac{2(x-2)^2}{(x-1)(x-2)(x-2)(x-3)}\\
			=&\lim\limits_{x\to  2}\dfrac{2}{(x-1)(x-3)}\\
			 =&\dfrac{2}{(1)\cdot(-1)}=-2.
		\end{align*}
	 }
\end{ex}

% nguồn: Ok_Toan11_TracNghiem_Lan2
\begin{ex}%[1D3H2-3]%[ma:204f98e5]
	$\lim\limits_{x\to +\infty}\dfrac{x-3}{6x^2+x+1}$ có kết quả  là
	\choice 
		{$-2$} 
		{$+\infty$} 
		 {$3$}
		{\True $0$} 
	 \loigiai{
		Ta có $\lim\limits_{x\to +\infty}\dfrac{x-3}{6x^2+x+1} 
		=\lim\limits_{x\to  +\infty}\dfrac{\dfrac{1}{x}-\dfrac{3}{x^2}}{6+\dfrac{1}{x}+\dfrac{1}{x^2}}=0$.
	 }
\end{ex}

% nguồn: Ok_Toan11_TracNghiem_Lan2
\begin{ex}%[1D3H2-3]%[ma:8ffd51c2]
	 Tính $\lim\limits_{x \to 1} \dfrac{x^2+x-2}{x-1}$.
	\choice 
		{$0$} 
		{$1$} 
		{\True  $3$}
		{$+\infty$} 
	\loigiai{ 
		Ta có $\lim\limits_{x \to  1} \dfrac{x^2+x-2}{x-1}=\lim\limits_{x \to 1} \dfrac{(x-1)(x+2)}{x-1}=\lim\limits_{x \to 1}(x+2)=3$.
	 }
\end{ex}

% nguồn: Ok_Toan11_TracNghiem_Lan2
\begin{ex}%[1D3H2-3]%[ma:542c26be]
	Tìm giới hạn $\lim\limits_{x\to  +\infty}\dfrac{2x-3}{1-3x}$.
	 \choice
		 {$\dfrac{2}{3}$}
		{\True  $-\dfrac{2}{3}$}
		{$-\dfrac{3}{2}$} 
		{$2$} 
	\loigiai{ 
	 Ta có $\lim\limits_{x\to +\infty}\dfrac{2x-3}{1-3x}$$=\lim\limits_{x\to +\infty}\dfrac{2-\dfrac{3}{x}}{\dfrac{1}{x}-3}=-\dfrac{2}{3}$.}
\end{ex}

% nguồn: Ok_Toan11_TracNghiem_Lan2
\begin{ex}%[1D3H2-3]%[ma:b58241f2]
	Cho $\lim\limits_{x\to 1} \dfrac{x^2+ax+b}{x^2-1}=\dfrac{-1}{2}$; $\left(a,b\in \mathbb{R}\right)$. Tổng $S=a^2+b^2$ bằng 
	 \choice
		 {$S=4$}
		{$S=1$} 
		{\True $S=13$} 
		 {$S=9$}
	\loigiai{ 
		Do $\lim\limits_{x\to 1} \dfrac{x^2+ax+b}{x^2-1}=\dfrac{-1}{2}; \left(a,b\in \mathbb{R}\right)$ nên phương trình $x^2+ax+b=0$ có một nghiệm $x=1$.\\ 
		Khi đó $a+b+1=0$ nên $x^2+ax-a-1=\left(x-1\right)\cdot  \left(x+a+1\right)$.\\
		Suy ra $\lim\limits_{x\to  1} \dfrac{x^2+ax+b}{x^2-1}=\lim\limits_{x\to 1} \dfrac{x+a+1}{x+1}=\dfrac{a+2}{2}$.\\
		Mà  $\lim\limits_{x\to 1} \dfrac{x^2+ax+b}{x^2-1}=\dfrac{-1}{2}; \left(a,b\in \mathbb{R}\right)$ nên $\dfrac{a+2}{2}=\dfrac{-1}{2} \Leftrightarrow a=-3$.
		\[\Rightarrow b=-a-1=2.\] 
		Vậy $S=a^2+b^2=13$. 
	} 
\end{ex}

% nguồn: Ok_Toan11_TracNghiem_Lan2
\begin{ex}%[1D3H2-3]%[ma:c73d4663]
	Giá trị $\displaystyle\lim\limits _{x  \to-\infty}\dfrac{2 x-1}{3-x}$ bằng
	\choice 
		{\True $-2$} 
		 {$\dfrac{2}{3}$}
		{$1$} 
		{$2$} 
	\loigiai{ 
	Ta có $\displaystyle\lim\limits _{x \to-\infty}\dfrac{2 x-1}{3-x}=\displaystyle\lim\limits  _{x \to-\infty}\dfrac{2-\dfrac{1}{x}}{\dfrac{3}{x}-1}=-2$.}
\end{ex}

% nguồn: Ok_Toan11_TracNghiem_Lan2
\begin{ex}%[1D3H2-3]%[ma:8f939989]
	Giá trị  $\displaystyle\lim\limits _{x \to-\infty}\dfrac{3-4 x}{5 x+2}$ bằng
	 \choice
		 {$\dfrac{5}{4}$}
		 {$-\dfrac{5}{4}$}
		{\True  $-\dfrac{4}{5}$}
		{$\dfrac{4}{5}$} 
	 \loigiai{
	Ta có $\displaystyle\lim\limits _{x \to-\infty}\dfrac{3-4 x}{5  x+2}=\displaystyle\lim\limits _{x \to-\infty}\dfrac{x\left(\dfrac{3}{x}-4\right)}{x\left(5+\dfrac{2}{x}\right)}=\displaystyle\lim\limits _{x \to-\infty}\dfrac{\left(\dfrac{3}{x}-4\right)}{\left(5+\dfrac{2}{x}\right)}=\dfrac{-4}{5}$.}
\end{ex}

% nguồn: Ok_Toan11_TracNghiem_Lan2
\begin{ex}%[1D3H2-3]%[ma:7ec86051]
	Tính  giới hạn $L=\lim\limits_{x\to-2} \dfrac{\sqrt{2x^2+x+3}-3}{4-x^2}$.
	 \choice
		{$L=-\dfrac{2}{7}$} 
		{\True  $L=-\dfrac{7}{24}$}
		 {$L=-\dfrac{9}{31}$}
		 {$L=0$}
	 \loigiai{
		Ta  có
		\begin{align*}
			L&=\lim\limits_{x\to-2} \dfrac{\left(\sqrt{2x^2+x+3}-3\right)\left(\sqrt{2x^2+x+3}+3\right)}{-\left(x^2-4\right)\left(\sqrt{2x^2+x+3}+3\right)}\\ 
			&=\lim\limits_{x\to-2}  \dfrac{2x^2+x-6}{-\left(x^2-4\right)\left(\sqrt{2x^2+x+3}+3\right)}\\
			&=\lim\limits_{x\to-2}  \dfrac{\left(2x-3\right)\left(x+2\right)}{-\left(x-2\right)\left(x+2\right)\left(\sqrt{2x^2+x+3}+3\right)}\\
			&=\lim\limits_{x\to-2}  \dfrac{3-2x}{\left(x-2\right)\left(\sqrt{2x^2+x+3}+3\right)}\\&=-\dfrac{7}{24}.
		\end{align*}
	} 
\end{ex}

% nguồn: Ok_Toan11_TracNghiem_Lan2
\begin{ex}%[1D3H2-3]%[ma:f132a637]
	Tính giới hạn $\lim\limits_{x\to-1}  \dfrac{x^2+2x+1}{2x^3+2}$.
	\choice 
		{$-\infty$} 
		{\True $0$} 
		{$+\infty$} 
		{$\dfrac{1}{2}$} 
	\loigiai{ 
		 Ta có $\lim\limits_{x\to-1} \dfrac{x^2+2x+1}{2x^3+2}$ $=\lim\limits_{x\to-1} \dfrac{\left(x+1\right)^2}{2\left(x+1\right)\left(x^2-x+1\right)}$ $=\lim\limits_{x\to-1} \dfrac{x+1}{2\left(x^2-x+1\right)}$ $=0$.
	 }
\end{ex}

% nguồn: Ok_Toan11_TracNghiem_Lan2
\begin{ex}%[1D3H2-3]%[ma:3be97d8b]
	Tìm  $\lim\limits_{x\to 2} \dfrac{2x^2-5x+2}{x-2}$.
	\choice 
		{$\dfrac{3}{2}$} 
		{\True $3$} 
		 {$1$}
		{$2$} 
	 \loigiai{
		Ta  có $\lim\limits_{x\to 2} \dfrac{2x^2-5x+2}{x-2}=\lim\limits_{x\to 2} \dfrac{\left(2x-1\right)\left(x-2\right)}{x-2}=\lim\limits_{x\to 2} \left(2x-1\right)=3$.
	} 
\end{ex}

% nguồn: Ok_Toan11_TracNghiem_Lan2
\begin{ex}%[1D3H2-3]%[ma:5db07a6b]
	Tính  giới hạn $\lim\limits_{x\to +\infty}\dfrac{2x-1}{4x+2}$.
	 \choice
		{\True  $\dfrac{1}{2}$}
		{$1$} 
		 {$\dfrac{-1}{4}$}
		 {$\dfrac{-1}{2}$}
	 \loigiai{
	 Ta có $\lim\limits_{x\to +\infty}\dfrac{2x-1}{4x+2}=\lim\limits_{x\to +\infty}\dfrac{2-\dfrac{1}{x}}{4+\dfrac{2}{x}}=\dfrac{1}{2}$.}
\end{ex}

% nguồn: Ok_Toan11_TracNghiem_Lan2
\begin{ex}%[1D3H2-3]%[ma:978b0dfa]
	 Giới hạn $\lim\limits_{x\to\pm\infty}\dfrac{2x^2-3x+7}{3-4x^2}$ bằng
	 \choice
		 {$\dfrac{7}{3}$}
		 {$\dfrac{1}{2}$}
		 {$\dfrac{2}{3}$}
		{\True $-\dfrac{1}{2}$} 
	 \loigiai{
	Ta  có $\lim\limits_{x\to\pm\infty}\dfrac{2x^2-3x+7}{3-4x^2}=\lim\limits_{x\to\pm\infty}\dfrac{2-\dfrac{3}{x}+\dfrac{7}{x^2}}{\dfrac{3}{x^2}-4}=-\dfrac{1}{2}$.
	} 
\end{ex}

% nguồn: Ok_Toan11_TracNghiem_Lan2
\begin{ex}%[1D3H2-3]%[ma:049a5696]
	 $\lim\limits_{x\to-\infty}\dfrac{\sqrt{9x^2+1}}{x+1}$ bằng
	\choice 
		 {$-9$}
		 {$3$}
		{\True $-3$} 
		 {$9$}
	\loigiai{ 
	Ta có  $K=\lim\limits_{x\to-\infty}\dfrac{\sqrt{9x^2+1}}{x+1}=\lim\limits_{x\to-\infty}\dfrac{-x\sqrt{9+\dfrac{1}{x^2}}}{x+1}=\lim\limits_{x\to-\infty}\dfrac{-\sqrt{9+\dfrac{1}{x^2}}}{1+\dfrac{1}{x}}=-3$.
	} 
\end{ex}

% nguồn: Ok_Toan11_TracNghiem_Lan2
\begin{ex}%[1D3H2-3]%[ma:e55393c8]
	Kết quả  của $\lim\limits_{x \to 1} \dfrac{x^3-1}{x-1}$ bằng
	 \choice
		 {$-\infty$}
		{$0$} 
		{\True $3$} 
		{$+\infty$} 
	\loigiai{ 
		Ta có $A=\lim\limits_{x \to 1} \dfrac{x^3-1}{x-1}=\lim\limits_{x  \to 1} \dfrac{(x-1)\left(x^2+x+1\right)}{x-1}=\lim\limits_{x \to 1}\left(x^2+x+1\right)=3$.
	 }
\end{ex}

% nguồn: Ok_Toan11_TracNghiem_Lan2
\begin{ex}%[1D3H2-3]%[ma:f0e11a52]
	Tính  $\lim\limits_{x\to+\infty}\dfrac{3x^2-x}{x^2+1}$.
	\choice 
		 {$+\infty$}
		 {$-1$}
		 {$-\infty$}
		 {\True $3$}
	 \loigiai{
	Ta  có $\lim\limits_{x\to+\infty}\dfrac{3x^2-x}{x^2+1}=\lim\limits_{x\to+\infty}\dfrac{3-\dfrac{1}{x}}{1+\dfrac{1}{x^2}}=3$.
	} 
\end{ex}

% nguồn: Ok_Toan11_TracNghiem_Lan2
\begin{ex}%[1D3H2-3]%[ma:a12f3442]
	 Tìm $\lim\limits_{x \to -2}\dfrac{x^2+5x+6}{x+2}$.
	 \choice
		 {$+\infty$}
		{$-2$} 
		 {$-1$}
		 {\True $1$}
	 \loigiai{
		Ta có $\lim\limits_{x \to -2} \dfrac{x^2+5x+6}{x+2}=\lim\limits_{x \to -2} \dfrac{(x+2)(x+3)}{x+2}=\lim\limits_{x \to -2}(x+3)=1$. 
	 }
\end{ex}

% nguồn: Ok_Toan11_TracNghiem_Lan2
\begin{ex}%[1D3H2-3]%[ma:7b47b4ce]
	Giới hạn $\lim\limits_{x  \to 4} \dfrac{x^2+2x-15}{x-3}$ bằng
	 \choice
		{$\dfrac{1}{8}$} 
		 {\True $9$}
		{$+\infty$} 
		{$8$} 
	 \loigiai{
		Ta có $\lim\limits_{x  \to 4} \dfrac{x^2+2x-15}{x-3}=\lim\limits_{x \to 4} \dfrac{(x-3)(x+5)}{x-3}=9$.
	 }
\end{ex}

% nguồn: Ok_Toan11_TracNghiem_Lan2
\begin{ex}%[1D3H2-3]%[ma:03aea61e]
	Cho $a$, $b$ là các số thực dương thỏa  mãn $a+b=8$ và $\lim\limits_{x\to 0} \dfrac{\sqrt{x^2+2ax+1}-\sqrt{bx+1}}{x}=5$. Trong các mệnh đề dưới đây, mệnh đề nào đúng?
	\choice 
		{$a\in \left(2;4\right)$} 
		 {\True $a\in \left(3;8\right)$}
		{$b\in  \left(3;5\right)$}
		{$b\in  \left(4;9\right)$}
	 \loigiai{
		Ta  có
		\begin{align*}
			\lim\limits_{x\to 0} \dfrac{\sqrt{x^2+2ax+1}-\sqrt{bx+1}}{x}&=\lim\limits_{x\to 0} \dfrac{\left(x^2+2ax+1\right)-\left(bx+1\right)}{x\left(\sqrt{x^2+2ax+1}+\sqrt{bx+1}  \right)}\\
			 &=\lim\limits_{x\to 0} \dfrac{x^2+\left(2a-b\right)x}{x\left(\sqrt{x^2+2ax+1}+\sqrt{bx+1} \right)}\\
			&=\lim\limits_{x\to  0} \dfrac{x+\left(2a-b\right)}{\sqrt{x^2+2ax+1}+\sqrt{bx+1}}=\dfrac{2a-b}{2}.
		\end{align*}
		 Mà
		\begin{align*}
			&\lim\limits_{x\to 0}  \dfrac{\sqrt{x^2+2ax+1}-\sqrt{bx+1}}{x}=5\\
			&\Leftrightarrow  \dfrac{2a-b}{2}=5\\
			&\Leftrightarrow  2a-b=10.
		\end{align*}
		Từ đó ta có hệ phương  trình $\heva{&a+b=8\\&2a-b=10}\Leftrightarrow \heva{&a=6\\&b=2.}$\\
		Vậy $a\in  \left(3;8\right)$.
	 }
\end{ex}

% nguồn: Ok_Toan11_TracNghiem_Lan2
\begin{ex}%[1D3H2-3]%[ma:52661b75]
	Giới hạn $\displaystyle\lim\limits _{x  \to+\infty}\dfrac{x^{2}-3 x+2}{2 x^{2}+1}$ có kết quả là
	 \choice
		{$+\infty$} 
		 {$-\infty$}
		 {$2$}
		{\True $\dfrac{1}{2}$} 
	 \loigiai{
	Ta có $\displaystyle\lim\limits _{x \to+\infty}\dfrac{x^{2}-3 x+2}{2 x^{2}+1}=\displaystyle\lim\limits _{x  \to+\infty}\dfrac{1-\dfrac{3}{x}+\dfrac{2}{x^{2}}}{2+\dfrac{1}{x^{2}}}=\dfrac{1}{2}$.}
\end{ex}

% nguồn: Ok_Toan11_TracNghiem_Lan2
\begin{ex}%[1D3H2-3]%[ma:dcd28ab6]
	Kết quả của giới hạn $\lim\limits_{x \to 2}  \dfrac{x^2-5x+6}{x-2}$ là
	\choice 
		{$1$} 
		 {$-2$}
		{$0$} 
		{\True  $-1$}
	 \loigiai{
		Ta có $\lim\limits_{x \to 2}  \dfrac{(x-2)(x-3)}{x-2}=\lim\limits_{x \to 2}(x-3)=2-3=-1$.
	 }
\end{ex}

% nguồn: Ok_Toan11_TracNghiem_Lan2
\begin{ex}%[1D3H2-3]%[ma:631c9777]
	Giá trị $\displaystyle\lim\limits _{x \to+\infty}\dfrac{x-2}{x+3}$ bằng 
	 \choice
		{$-\dfrac{2}{3}$} 
		 {\True $1$}
		{$2$} 
		 {$-3$}
	 \loigiai{
	Chia cả tử và mẫu cho $x$, ta có $\displaystyle\lim\limits _{x  \to+\infty}\dfrac{x-2}{x+3}=\displaystyle\lim\limits _{x \to+\infty}\dfrac{1-\dfrac{2}{x}}{1+\dfrac{3}{x}}=\dfrac{1}{1}=1$.}
\end{ex}

% nguồn: Ok_Toan11_TracNghiem_Lan2
\begin{ex}%[1D3H2-3]%[ma:348d6bc2]
	$\lim\limits_{x\to  -\infty}\dfrac{(-x^2+1)^3(1-2x)^{94}}{2x^{100}+3}$ có kết quả là
	 \choice
		 {$-\infty$}
		 {$-2^{100}$}
		{$2^{93}$} 
		{\True $-2^{93}$} 
	\loigiai{ 
		Ta có 
		\begin{eqnarray*}
			\lim\limits_{x\to  -\infty}\dfrac{(-x^2+1)^3(1-2x)^{94}}{2x^{100}+3}
			&=&\lim\limits_{x\to -\infty}\dfrac{\left[x^2\left(-1+\dfrac{1}{x^2}\right)\right]^3\left[x\left(\dfrac{1}{x}-2\right)\right]^{94}}{x^{100}\left(2+\dfrac{3}{x^{100}}\right)}  \\
			&=&\lim\limits_{x\to  -\infty}\dfrac{x^6\left(-1+\dfrac{1}{x^2}\right)^3x^{94}\left(\dfrac{1}{x}-2\right)^{^{94}}}{x^{100}\left(2+\dfrac{3}{x^{100}}\right)} \\
			&=&\lim\limits_{x\to -\infty}\dfrac{\left(-1+\dfrac{1}{x^2}\right)^3\left(\dfrac{1}{x}-2\right)^{94}}{2+\dfrac{3}{x^{100}}}  \\
			 &=&\dfrac{(-1)^3\cdot{(-2)^{94}}}{2}=-2^{93}.
		\end{eqnarray*}
	 }
\end{ex}

% nguồn: Ok_Toan11_TracNghiem_Lan2
\begin{ex}%[1D3H2-3]%[ma:bd7255e4]
	Tính $\lim\limits_{x\to+\infty}\dfrac{x+3}{\sqrt{4x^2+1}-2}$. 
	 \choice
		 {$\dfrac{1}{4}$}
		 {$-\dfrac{3}{2}$}
		{\True  $\dfrac{1}{2}$}
		{$0$} 
	 \loigiai{
	Ta có $\lim\limits_{x\to+\infty}\dfrac{x+3}{\sqrt{4x^2+1}-2}=\lim\limits_{x\to+\infty}\dfrac{1+\dfrac{3}{x}}{\sqrt{4+\dfrac{1}{x^2}}-\dfrac{2}{x}}=\dfrac{1}{2}$. 
	 }
\end{ex}

% nguồn: Ok_Toan11_TracNghiem_Lan2
\begin{ex}%[1D3H2-3]%[ma:7e6e7c31]
	Kết quả $L=\lim\limits_{x \to -1} \dfrac{x^2-x-2}{3x^2+8x+5}$  là
	 \choice
		{\True  $L=-\dfrac{3}{2}$}
		 {$L=\dfrac{1}{2}$}
		{$L=-\infty$} 
		{$L=0$} 
	 \loigiai{
		Ta có $L=\lim\limits_{x \to  -1} \dfrac{x^2-x-2}{3x^2+8x+5}=\lim\limits_{x \to -1} \dfrac{(x+1)(x-2)}{(x+1)(3x+5)}=\lim\limits_{x \to -1} \dfrac{x-2}{3x+5}=-\dfrac{3}{2}$.
	 }
\end{ex}

% nguồn: Ok_Toan11_TracNghiem_Lan2
\begin{ex}%[1D3H2-3]%[ma:3c3a026b]
	$\lim\limits_{x\to+\infty}\dfrac{5}{3x+2}$ bằng  bao nhiêu?
	 \choice
		{\True $0$} 
		{$1$} 
		 {$+\infty$}
		 {$\dfrac{5}{3}$}
	 \loigiai{
	Ta có $\lim\limits_{x\to+\infty}\dfrac{5}{3x+2}=\lim\limits_{x\to+\infty}\dfrac{\dfrac{5}{x}}{3+\dfrac{2}{x}}=0$. 
	 }
\end{ex}

% nguồn: Ok_Toan11_TracNghiem_Lan2
\begin{ex}%[1D3H2-3]%[ma:b9527c5b]
	Cho hàm số $f(x)$ xác định trên  $\mathbb{R}$ thỏa mãn $\lim\limits_{x\to -3}\dfrac{f(x)-2}{x+2}=5$. Tính $\lim\limits_{x\to -3}\dfrac{-4+\sqrt[3]{5f(x)+7}}{x^2+x}$.
	\choice 
		{$4$} 
		 {$-2$}
		{\True  $-1$}
		 {$-\dfrac{1}{4}$}
	 \loigiai{
		Ta có  $\lim\limits_{x\to -3}\dfrac{f(x)-2}{x+2}=5
		\Leftrightarrow \dfrac{\lim\limits_{x\to  -3}\left[f(x)-2\right]}{\lim\limits_{x\to -3}(x+2)}=5
		\Leftrightarrow  \dfrac{\lim\limits_{x\to -3}f(x)-2}{-1}=5
		\Leftrightarrow \lim\limits_{x\to -3}f(x)=-3$.\\ 
		Vậy  $\lim\limits_{x\to -3}\dfrac{-4+\sqrt[3]{5f(x)+7}}{x^2+x}
		=\dfrac{-4+\sqrt[3]{5\cdot (-3)+7}}{(-3)^2-3}=-1$. 
	 }
\end{ex}

% nguồn: Ok_Toan11_TracNghiem_Lan2
\begin{ex}%[1D3H2-3]%[ma:8d856014]
	Cho $a\in \mathbb{R}$ thỏa mãn $\lim\limits_{x\to +\infty}\dfrac{ax+\sqrt{x^2+x+1}}{2x-1}=2$.  Khẳng định nào sau đây là khẳng định đúng?
	 \choice
		 {$a<-1$}
		{$-1\le a\le  1$}
		{$1\le  a<2$}
		 {\True $a\ge 2$}
	\loigiai{ 
		 Ta có $\lim\limits_{x\to +\infty}\dfrac{ax+\sqrt{x^2+x+1}}{2x-1}=2
		 \Leftrightarrow \lim\limits_{x\to +\infty} \dfrac{x\left(a+\sqrt{1+\dfrac{1}{x}+\dfrac{1}{x^2}}\right)}{x\left(2-\dfrac{1}{x}\right)}=2
		\Leftrightarrow \dfrac{a+1}{2}=2\Leftrightarrow  a=3$.\\
		Vậy $a=3$. 
	 }
\end{ex}

% nguồn: Ok_Toan11_TracNghiem_Lan2
\begin{ex}%[1D3H2-3]%[ma:418303c4]
	Giá trị của  giới hạn $\lim\limits_{x\to 3}\sqrt{\dfrac{9x^2-x}{(2x-1)(x^4-3)}}$ là
	\choice 
		{$\dfrac{1}{5}$} 
		 {$\sqrt{5}$}
		{\True  $\dfrac{1}{\sqrt{5}}$}
		 {$5$}
	\loigiai{ 
		Ta  có $\lim\limits_{x\to 3}\sqrt{\dfrac{9x^2-x}{(2x-1)(x^4-3)}}
		=\sqrt{\dfrac{9\cdot  3^2-3}{(2\cdot 3-1)(3^4-3)}}=\dfrac{1}{\sqrt{5}}$.
	 }
\end{ex}

% nguồn: Ok_Toan11_TracNghiem_Lan2
\begin{ex}%[1D3H2-3]%[ma:65f54b92]
	Tính giới hạn $\lim\limits_{x\to  -\infty}\dfrac{\sqrt{x^2-x+1}}{2x}$.
	 \choice
		{$\dfrac{1}{2}$} 
		 {$+\infty $}
		{$-\infty $} 
		 {\True $-\dfrac{1}{2}$}
	\loigiai{ 
	$\lim\limits_{x\to  -\infty}\dfrac{\sqrt{x^2-x+1}}{2x}=\lim\limits_{x\to -\infty}\dfrac{-x\sqrt{1-\dfrac{1}{x}+\dfrac{1}{x^2}}}{2x}=\lim\limits_{x\to -\infty}\dfrac{-\sqrt{1-\dfrac{1}{x}+\dfrac{1}{x^2}}}{2}=-\dfrac{1}{2}$.}
\end{ex}

% nguồn: Ok_Toan11_TracNghiem_Lan2
\begin{ex}%[1D3H2-3]%[ma:f0f9ee08]
	Giá trị $\displaystyle\lim\limits _{x  \to-\infty}\dfrac{x+1}{6 x-2}$ bằng
	\choice 
		{$\dfrac{1}{2}$} 
		{\True  $\dfrac{1}{6}$}
		{$\dfrac{1}{3}$} 
		{$1$} 
	\loigiai{ 
		Ta có $\displaystyle\lim\limits _{x \to-\infty}\dfrac{x+1}{6 x-2}=\displaystyle\lim\limits _{x \to-\infty}\dfrac{1+\dfrac{1}{x}}{6-\dfrac{2}{x}}=\dfrac{1}{6}$. 
	 }
\end{ex}

% nguồn: Ok_Toan11_TracNghiem_Lan2
\begin{ex}%[1D3H2-3]%[ma:05841255]
	Tìm giới hạn $\lim\limits_{x\to  1}\dfrac{x^2+\left(3a+2\right)x-3a-3}{x-1}$.
	\choice 
		 {$4-3a$}
		{\True $3a+4$} 
		{$3a-4$} 
		 {$3a$}
	\loigiai{ 
		Ta có $\lim\limits_{x\to 1} \dfrac{x^2+\left(3a+2\right)x-3a-3}{x-1}=\lim\limits_{x\to 1} \dfrac{\left(x-1\right)\left(x+3a+3\right)}{x-1} =\lim\limits_{x\to  1} \left(x+3a+3\right)=3a+4$.
	 }
\end{ex}

% nguồn: Ok_Toan11_TracNghiem_Lan2
\begin{ex}%[1D3H2-3]%[ma:7411538d]
	 Tìm $\lim\limits_{x\to +\infty}(x+1-\sqrt[3]{x^3+2})$.
	 \choice
		 {$-1$}
		{$-\infty$} 
		 {$+\infty$}
		{\True  $1$}
	\loigiai{ 
		Ta có $\lim\limits_{x\to  +\infty}(1+x-\sqrt[3]{x^3+2})=\lim\limits_{x\to +\infty}\left(1+\dfrac{-2}{x^2+x\sqrt[3]{x^3+2}+\left(\sqrt[3]{x^3+2}\right)^2}\right)$\\
		=  $\lim\limits_{x\to +\infty}\left[1+\dfrac{-2}{x^2\left(1+\sqrt[3]{1+\dfrac{2}{x^3}}+\left(\sqrt[3]{1+\dfrac{2}{x^3}}\right)^2\right)}\right]=\lim\limits_{x\to +\infty}\left[1+\dfrac{\dfrac{-2}{x^2}}{1+\sqrt[3]{1+\dfrac{2}{x^3}}+\left(\sqrt[3]{1+\dfrac{2}{x^3}}\right)^2}\right]=1$.\\
	Vậy $\lim\limits_{x\to +\infty}(x+1-\sqrt[3]{x^3+2})=1$} 
\end{ex}

% nguồn: Ok_Toan11_TracNghiem_Lan2
\begin{ex}%[1D3H2-3]%[ma:7becf991]
	Tính $\displaystyle\lim\limits  _{x \to+\infty}\dfrac{x+\sin x}{x}$.
	 \choice
		{$\dfrac{1}{2}$} 
		 {$+\infty$}
		{\True  $1$}
		 {$0$}
	\loigiai{ 
		Ta có $\left|\dfrac{\sin x}{x}\right|=\dfrac{\left|\sin x \right|}{\left|x\right|} \leq  \dfrac{1}{\left|x\right|}$. Vì $\lim\limits_{x\to +\infty} \dfrac{1}{\left|x\right|} = 0$ nên $\lim\limits_{x\to +\infty} \dfrac{\sin x}{x}=0$.\\
	Khi đó $\displaystyle\lim\limits _{x  \to+\infty}\dfrac{x+\sin x}{x}=\displaystyle\lim\limits _{x \to+\infty}\dfrac{x}{x}+\displaystyle\lim\limits _{x \to+\infty}\dfrac{\sin x}{x}=1+\displaystyle\lim\limits _{x \to+\infty}\dfrac{\sin x}{x}=1+0=1$.}
\end{ex}

% nguồn: Ok_Toan11_TracNghiem_Lan2
\begin{ex}%[1D3H2-3]%[ma:0d5e4c05]
	Kết quả  của $\lim\limits_{x \to -1} \dfrac{x^2+2x+1}{2x+2}$ là
	 \choice
		{$-\infty$} 
		 {\True $0$}
		{$\dfrac{1}{2}$} 
		 {$+\infty$}
	\loigiai{ 
		Ta có $\lim\limits_{x \to -1} \dfrac{x^2+2x+1}{2x+2}=\lim\limits_{x \to -1} \dfrac{(x+1)^2}{2(x+1)}=\lim\limits_{x  \to -1} \dfrac{x+1}{2}=0$.
	 }
\end{ex}

% nguồn: Ok_Toan11_TracNghiem_Lan2
\begin{ex}%[1D3H2-3]%[ma:3c939745]
	Biết $\lim\limits_{x\to -\infty}\dfrac{\sqrt{4x^2+x+1}+4}{mx-2}=\dfrac{1}{2}$. Khi đó $m$  thuộc tập hợp nào?
	\choice 
		{$\left[3;6\right]$} 
		 {$\left[-3;0\right]$}
		{\True  $\left[-6;-3\right]$}
		 {$\left[1;3\right]$}
	\loigiai{ 
		Ta có 
		\begin{eqnarray*}
			\lim\limits_{x\to  -\infty}\dfrac{\sqrt{4x^2+x+1}+4}{mx-2}=\dfrac{1}{2}
			 &\Leftrightarrow& \lim\limits_{x\to -\infty}\dfrac{x\left(-\sqrt{4+\dfrac{1}{x}+\dfrac{1}{x^2}}+\dfrac{4}{x}\right)}{x\left(m-\dfrac{2}{x}\right)}=\dfrac{1}{2} \\
			&\Leftrightarrow& \dfrac{-2}{m}=\dfrac{1}{2}\Leftrightarrow m=-4. 
		\end{eqnarray*}
		Vậy  $m=-4\in \left[-6;-3\right]$.
	} 
\end{ex}

% nguồn: Ok_Toan11_TracNghiem_Lan2
\begin{ex}%[1D3H2-3]%[ma:d8c6605a]
	Giá  trị $\displaystyle\lim\limits _{x \to-\infty}\dfrac{4 x+1}{-x+1}$ bằng
	 \choice
		{$2$} 
		{$4$} 
		{$-1$} 
		 {\True $-4$}
	 \loigiai{
	Ta có  $\displaystyle\lim\limits _{x \to-\infty}\dfrac{4 x+1}{-x+1}=\displaystyle\lim\limits _{x \to-\infty}\dfrac{4+\dfrac{1}{x}}{-1+\dfrac{1}{x}}=-4$.}
\end{ex}

% nguồn: Ok_Toan11_TracNghiem_Lan2
\begin{ex}%[1D3H2-3]%[ma:39c6f6fe]
	 Biết $\lim\limits_{x\to 1} \dfrac{x^3-ax+a-1}{x-1}=2$. Tính $M=a^2+2a$.
	\choice 
		 {\True $M=3$}
		{$M=1$} 
		{$M=-1$} 
		 {$M=8$}
	 \loigiai{
		Ta  có
		\begin{eqnarray*}
			&  &\lim\limits_{x\to 1} \dfrac{x^3-ax+a-1}{x-1}\\
			&=&\lim\limits_{x\to  1} \dfrac{\left(x-1\right)\left(x^2+x+1\right)-a\left(x-1\right)}{x-1}\\
			&=&\lim\limits_{x\to  1} \left(x^2+x+1-a\right)=3-a\\
			 &\Rightarrow& a=1.
		\end{eqnarray*}
		Vậy  $M=a^2+2a=3$.
	 }
\end{ex}

% nguồn: Ok_Toan11_TracNghiem_Lan2
\begin{ex}%[1D3H2-3]%[ma:b145c616]
	Tính  $\lim\limits_{x\to-\infty}\dfrac{\sqrt{x^2-x}+5x}{x+1}$.
	\choice 
		{$-\infty$} 
		 {$-1$}
		 {$6$}
		{\True $4$} 
	 \loigiai{
	Ta có  $\lim\limits_{x\to-\infty}\dfrac{\sqrt{x^2-x}+5x}{x+1}=\lim\limits_{x\to-\infty}\dfrac{-\sqrt{1-\dfrac{1}{x}}+5}{1+\dfrac{1}{x}}=4$.
	} 
\end{ex}

% nguồn: Ok_Toan11_TracNghiem_Lan2
\begin{ex}%[1D3H2-3]%[ma:7af74250]
	Tính giới hạn $I=\displaystyle\lim\limits _{x \to-\infty}\dfrac{3 x-2}{2 x+1}$. 
	 \choice
		 {$I=-2$}
		 {$I=-\dfrac{3}{2}$}
		{$I=2$} 
		 {\True $I=\dfrac{3}{2}$}
	\loigiai{ 
	Ta có $I=\displaystyle\lim\limits _{x \to-\infty}\dfrac{3  x-2}{2 x+1}=\displaystyle\lim\limits _{x \to-\infty}\dfrac{3-\dfrac{2}{x}}{2+\dfrac{1}{x}}=\dfrac{3}{2}$.}
\end{ex}

% nguồn: Ok_Toan11_TracNghiem_Lan2
\begin{ex}%[1D3H2-3]%[ma:522aeb63]
	Tính 
	 $\lim\limits_{x\to+\infty}\left(\dfrac{1}{2x^2+1}\right)\sqrt{\dfrac{4x^5-x^4+5}{2x-3}}$.
	 \choice
		{$+\infty$} 
		 {\True $\dfrac{\sqrt{2}}{2}$}
		{$1$} 
		 {$\dfrac{1}{2}$}
	\loigiai{ 
	$\begin{aligned}
	\lim\limits_{x\to+\infty}\left(\dfrac{1}{2 x^2+1}\right) \sqrt{\dfrac{4 x^5-x^4+5}{2  x-3}}&=\lim\limits_{x\to+\infty} \dfrac{x^2}{2 x^2+1} \sqrt{\dfrac{4 x-1+\dfrac{5}{x^4}}{2 x-3}}\\&=\lim\limits_{x\to+\infty} \dfrac{1}{2+\dfrac{1}{x^2}} \sqrt{\dfrac{4-\dfrac{1}{x}+\dfrac{5}{x^5}}{2-\dfrac{3}{x}}}=\dfrac{\sqrt{2}}{2}.
\end{aligned}$
} 
\end{ex}

% nguồn: Ok_Toan11_TracNghiem_Lan2
\begin{ex}%[1D3H2-3]%[ma:5791527a]
	 Giới hạn $\lim\limits_{x\to-\infty}\dfrac{1-x}{3x+2}$ bằng
	 \choice
		 {$\dfrac{1}{3}$}
		{$\dfrac{1}{2}$} 
		{\True  $-\dfrac{1}{3}$}
		 {$-\dfrac{1}{2}$}
	 \loigiai{
	Có  $\lim\limits_{x\to-\infty}\dfrac{1-x}{3x+2}=\lim\limits_{x\to-\infty}\dfrac{x\left(\dfrac{1}{x}-1\right)}{x\left(3+\dfrac{2}{x}\right)}=\lim\limits_{x\to-\infty}\dfrac{\dfrac{1}{x}-1}{3+\dfrac{2}{x}}=\dfrac{0-1}{3+0}=\dfrac{-1}{3}$.
	} 
\end{ex}

% nguồn: Ok_Toan11_TracNghiem_Lan2
\begin{ex}%[1D3H2-3]%[ma:cfe37216]
	 Tính giới hạn $L=\lim\limits_{x\to 2} \dfrac{\sqrt{2x^2-x+3}-3}{4-x^2}$.
	\choice 
		 {$L=-\dfrac{2}{7}$}
		{\True  $L=-\dfrac{7}{24}$}
		 {$L=\dfrac{9}{31}$}
		 {$L=1$}
	\loigiai{ 
		Ta có 
		\begin{align*}
			 L&=\lim\limits_{x\to 2} \dfrac{\left(\sqrt{2x^2-x+3}-3\right)\left(\sqrt{2x^2-x+3}+3\right)}{-\left(x^2-4\right)\left(\sqrt{2x^2-x+3}+3\right)}\\
			&=\lim\limits_{x\to  2} \dfrac{2x^2-x-6}{-\left(x^2-4\right)\left(\sqrt{2x^2-x+3}+3\right)}\\
			&=\lim\limits_{x\to  2} \dfrac{\left(2x+3\right)\left(x-2\right)}{-\left(x-2\right)\left(x+2\right)\left(\sqrt{2x^2-x+3}+3\right)}\\
			&=\lim\limits_{x\to 2}  \dfrac{-2x-3}{\left(x+2\right)\left(\sqrt{2x^2-x+3}+3\right)}\\&=-\dfrac{7}{24}.
		\end{align*}
	} 
\end{ex}

% nguồn: Ok_Toan11_TracNghiem_Lan2
\begin{ex}%[1D3H2-3]%[ma:372c9e73]
	 $\lim\limits_{x\to+\infty}\dfrac{2x-5}{-x+3}$ bằng
	\choice 
		{$\dfrac{-5}{3}$} 
		 {$-1$}
		{$3$} 
		 {\True $-2$}
	\loigiai{ 
	$\lim\limits_{x\to+\infty} \dfrac{2 x-5}{-x+3}=\lim\limits_{x\to+\infty} \dfrac{2-\dfrac{5}{x}}{-1+\dfrac{3}{x}}=\dfrac{2}{-1}=-2$. 
	 }
\end{ex}

% nguồn: Ok_Toan11_TracNghiem_Lan2
\begin{ex}%[1D3H2-3]%[ma:a6d52bc4]
	Cho hàm số $y=f(x)=\dfrac{2\sqrt{1+x}-\sqrt[3]{8-x}}{x}$. Tính $\lim\limits_{x\to 0}  f(x)$.
	\choice 
		 {$\dfrac{1}{12}$}
		{\True $\dfrac{13}{12}$} 
		{$+\infty$} 
		{$\dfrac{10}{11}$} 
	 \loigiai{
		 Ta có
		\begin{align*}
			 \dfrac{2\sqrt{1+x}-\sqrt[3]{8-x}}{x}=&\dfrac{\left(2\sqrt{1+x}-2\right)+\left(2-\sqrt[3]{8-x} \right)}{x}\\ =&\dfrac{2\left(\sqrt{1+x}-1\right)}{x}+\dfrac{2-\sqrt[3]{8-x}}{x}\\
			 =&\dfrac{2}{\sqrt{1+x}+1}+\dfrac{1}{4+2\sqrt[3]{8-x}+\sqrt[3]{\left(8-x\right)^2}}.
		\end{align*}
		 Do vậy
		\begin{eqnarray*}
			& &\lim\limits_{x\to 0} f(x)=\lim\limits_{x\to 0} \left[\dfrac{2}{\sqrt{1+x}+1}+\dfrac{1}{4+2\sqrt[3]{8-x}+\sqrt[3]{\left(8-x\right)^2}}  \right]\\
			&=&\lim\limits_{x\to 0}  \dfrac{2}{\sqrt{1+x}+1}+\lim\limits_{x\to 0} \dfrac{1}{4+2\sqrt[3]{8-x}+\sqrt[3]{\left(8-x\right)^2}}\\
			 &=&1+\dfrac{1}{12}=\dfrac{13}{12}.
		\end{eqnarray*}
	 }
\end{ex}

% nguồn: Ok_Toan11_TracNghiem_Lan2
\begin{ex}%[1D3H2-3]%[ma:395592fa]
	Tìm $\lim\limits_{x\to 1}\dfrac{x^4-3x^2+2}{x^3+2x-3}$. 
	 \choice
		 {$-\dfrac{5}{2}$}
		{\True $-\dfrac{2}{5}$} 
		 {$\dfrac{1}{5}$}
		{$+\infty$} 
	 \loigiai{Ta có
		$\lim\limits_{x\to  1}\dfrac{x^4-3x^2+2}{x^3+2x-3}=\lim\limits_{x\to 1}\dfrac{(x-1)(x+1)(x^2-2)}{(x-1)(x^2+x+3)}=\lim\limits_{x\to 1}\dfrac{(x+1)(x^2-2)}{x^2+x+3}=\dfrac{2\cdot(-1)}{5}=-\dfrac{2}{5}$.
	} 
\end{ex}

% nguồn: Ok_Toan11_TracNghiem_Lan2
\begin{ex}%[1D3H2-3]%[ma:895416ea]
	Giới hạn $\lim\limits_{x\to 3}  \dfrac{x+1-\sqrt{5x+1}}{x-\sqrt{4x-3}}= \dfrac{a}{b}$ trong đó $\dfrac{a}{b}$ là phân số tối giản, $a$, $b$ là các số nguyên dương. Giá trị của $a-b$ là
	\choice 
		 {\True $1$}
		 {$\dfrac{1}{9}$}
		{$-1$} 
		 {$\dfrac{9}{8}$}
	 \loigiai{
		Ta có 
		\begin{align*}
			 \lim\limits_{x\to 3} \dfrac{x+1-\sqrt{5x+1}}{x-\sqrt{4x-3}}&=\lim\limits_{x\to 3} \dfrac{\left[\left(x+1\right)^2-\left(5x+1\right)\right]\left[x+\sqrt{4x-3} \right]}{\left[x^2-4x+3\right]\left[x+1+\sqrt{5x+1} \right]}\\
			&= \lim\limits_{x\to 3} \dfrac{\left[x^2-3x\right]\left[x+\sqrt{4x-3}  \right]}{\left[x^2-4x+3\right]\left[x+1+\sqrt{5x+1} \right]}\\
			&=\lim\limits_{x\to 3} \dfrac{x\left(x+\sqrt{4x-3} \right)}{\left(x-1\right)\left(x+1+\sqrt{5x+1} \right)}=\dfrac{9}{8}. 
		\end{align*}
		 Do đó $a=9$; $b=8$ nên $a-b=1$.
	} 
\end{ex}

% nguồn: Ok_Toan11_TracNghiem_Lan2
\begin{ex}%[1D3H2-3]%[ma:b35ca3bb]
	Giới hạn $\lim\limits_{x\to-4} \dfrac{x^2+3x-4}{x^2+4x}$  bằng
	\choice 
		{\True  $\dfrac{5}{4}$}
		{$-\dfrac{5}{4}$} 
		{$1$} 
		 {$-1$}
	 \loigiai{
		 Ta có $\lim\limits_{x\to-4} \dfrac{x^2+3x-4}{x^2+4x}=\lim\limits_{x\to-4} \dfrac{\left(x+4\right)\left(x-1\right)}{x\left(x+4\right)}=\lim\limits_{x\to-4} \dfrac{x-1}{x}=\dfrac{-4-1}{-4}=\dfrac{5}{4}$.
	 }
\end{ex}

% nguồn: Ok_Toan11_TracNghiem_Lan2
\begin{ex}%[1D3H2-3]%[ma:aee3ee98]
	Giới hạn $\lim\limits_{x \to  -1} \dfrac{x^2+x}{x^2+3x+2}$ bằng
	 \choice
		 {$0$}
		{$\dfrac{2}{3}$} 
		{\True $-1$} 
		 {$2$}
	\loigiai{ 
		Ta có $\lim\limits_{x \to -1} \dfrac{x^2+x}{x^2+3x+2}=\lim\limits_{x \to -1} \dfrac{x(x+1)}{(x+1)(x+2)}=\lim\limits_{x \to  -1} \dfrac{x}{x+2}=-1$.
	} 
\end{ex}

% nguồn: Ok_Toan11_TracNghiem_Lan2
\begin{ex}%[1D3H2-3]%[ma:f8369487]
	Kết quả của $\lim\limits_{x\to  +\infty}\dfrac{1}{x}(\sqrt{x^2+16}-x)$ bằng
	\choice 
		 {$-\infty$}
		 {4}
		 {\True 0}
		 {$+\infty$}
	\loigiai{ 
	Ta có $\lim\limits_{x\to +\infty}\dfrac{1}{x}(\sqrt{x^2+16}-x)=\lim\limits_{x\to  +\infty}\dfrac{1}{x}\left(\dfrac{x^2+16-x^2}{\sqrt{x^2+16}+x}\right)=\lim\limits_{x\to +\infty}\dfrac{1}{x}\left(\dfrac{16}{\sqrt{x^2+16}+x}\right)=0$.}
\end{ex}

% nguồn: Ok_Toan11_TracNghiem_Lan2
\begin{ex}%[1D3H2-3]%[ma:3e2a344a]
	Giới hạn $\lim\limits_{x \to 3} \dfrac{x+1-\sqrt{5x+1}}{x-\sqrt{4x-3}}=\dfrac{a}{b}$, với $a$, $b \in Z$, $b > 0$ và  $\dfrac{a}{b}$ là phân số tối giản. Giá trị của $a-b$là
	\choice 
		{\True  $1$}
		{$-1$} 
		{$\dfrac{8}{9}$} 
		 {$\dfrac{1}{9}$}
	 \loigiai{
		\begin{align*}
			\lim\limits_{x \to  3} \dfrac{x+1-\sqrt{5x+1}}{x-\sqrt{4x-3}}
			 &=\lim\limits_{x \to 3} \left[\dfrac{x+\sqrt{4x-3}}{x+1+\sqrt{5x+1}}\cdot \dfrac{\left(x+1\right)^2-\left(5x+1\right)}{x^2-4x+3}\right]\\
			 &=\lim\limits_{x \to 3} \left[\dfrac{x+\sqrt{4x-3}}{x+1+\sqrt{5x+1}}\cdot \dfrac{x^2-3x}{x^2-4x+3}\right]\\
			&=\lim\limits_{x \to 3} \left[\dfrac{x+\sqrt{4x-3}}{x+1+\sqrt{5x+1}}\cdot  \dfrac{x}{x-1}\right]\\
			&=\dfrac{6}{8}\cdot  \dfrac{3}{2}=\dfrac{9}{8}.
		\end{align*}
		Do đó,  $a=9$, $b=8$.\\
		Suy ra  $a-b=1$.
	} 
\end{ex}

% nguồn: Ok_Toan11_TracNghiem_Lan2
\begin{ex}%[1D3H2-3]%[ma:e517e8da]
	$\lim\limits_{x\to +\infty}\dfrac{2x^2+x}{x^2-1}$  bằng
	\choice 
		{$-2$} 
		 {$1$}
		 {\True $2$}
		{$-1$} 
	\loigiai{ 
	$\lim\limits_{x\to +\infty}\dfrac{2x^2+x}{x^2-1}=\lim\limits_{x\to +\infty} \dfrac{2x(x+1)}{(x-1)(x+1)}=\lim\limits_{x\to +\infty} \dfrac{2x}{(x-1)}=\lim\limits_{x\to +\infty}  \dfrac{2}{1-\dfrac{1}{x}}=2$.}
\end{ex}

% nguồn: Ok_Toan11_TracNghiem_Lan2
\begin{ex}%[1D3H2-3]%[ma:f2f50719]
	Tính $\lim\limits_{x\to 3} \dfrac{x^2-4x+3}{x^2-9}$ 
	\choice 
		 {$\dfrac{1}{2}$}
		 {$\dfrac{2}{5}$}
		{\True $\dfrac{1}{3}$} 
		{$\dfrac{1}{5}$} 
	 \loigiai{
		Ta có $\lim\limits_{x\to 3}  \dfrac{x^2-4x+3}{x^2-9}=\lim\limits_{x\to 3} \dfrac{\left(x-3\right)\left(x-1\right)}{\left(x-3\right)\left(x+3\right)}=\lim\limits_{x\to 3} \dfrac{x-1}{x+3}=\dfrac{1}{3}$.
	} 
\end{ex}

% nguồn: Ok_Toan11_TracNghiem_Lan2
\begin{ex}%[1D3H2-3]%[ma:9cc88259]
	Giá trị  $\displaystyle\lim\limits _{x \to+\infty}\dfrac{x+1}{4 x+3}$ bằng
	\choice 
		{$\dfrac{1}{3}$} 
		{\True $\dfrac{1}{4}$} 
		 {$3$}
		{$1$} 
	\loigiai{ 
	Ta có $\displaystyle\lim\limits _{x  \to+\infty}\dfrac{x+1}{4 x+3}=\displaystyle\lim\limits _{x \to+\infty}\dfrac{1+\dfrac{1}{x}}{4+\dfrac{3}{x}}=\dfrac{1}{4}$.}
\end{ex}

% nguồn: Ok_Toan11_TracNghiem_Lan2
\begin{ex}%[1D3H2-3]%[ma:e4fb4cdb]
	Biết rằng $b<0$, $a+b=5$ và $\lim\limits_{x\to 0}\dfrac{\sqrt[3]{ax+1}-\sqrt{1-b+x}}{x+1}=-2$. Khẳng định nào dưới đây \textbf{sai}? 
	 \choice
		{\True $8<a<13$} 
		 {$b<-1$}
		 {$a^2+b^2>10$}
		 {$a-b>0$}
	\loigiai{ 
		Ta có  $\lim\limits_{x\to 0}\dfrac{\sqrt[3]{ax+1}-\sqrt{1-b+x}}{x+1}
		=\dfrac{1-\sqrt{1-b}}{1}=-2 
		\Leftrightarrow  1-\sqrt{1-b}=-2
		 \Leftrightarrow b=-8$.\\
		Mà $a+b=5\Leftrightarrow  a+(-8)=5\Leftrightarrow a=13$.
	 }
\end{ex}

% nguồn: Ok_Toan11_TracNghiem_Lan2
\begin{ex}%[1D3H2-3]%[ma:dd511143]
	 Tính $I=\lim\limits_{x\to 2} \dfrac{x^3-8}{x^2-3x+2}$.
	 \choice
		 {$I=-12$}
		{\True  $I=12$}
		 {$I=-8$}
		 {$I=8$}
	\loigiai{ 
		Ta  có
		\begin{align*}
			I=\lim\limits_{x\to 2} \dfrac{x^3-8}{x^2-3x+2}&=\lim\limits_{x\to 2}  \dfrac{\left(x-2\right)\left(x^2+2x+4\right)}{\left(x-1\right)\left(x-2\right)}\\
			&=\lim\limits_{x\to 2} \dfrac{x^2+2x+4}{x-1}\\ 
			 &=\dfrac{2^2+2\cdot 2+4}{2-1}=12.
		\end{align*}
	} 
\end{ex}

% nguồn: Ok_Toan11_TracNghiem_Lan2
\begin{ex}%[1D3H2-3]%[ma:87a2711e]
	 Kết quả của $\lim\limits_{x\to -\infty}\dfrac{1}{x}(\sqrt{x^2+16})$ bằng
	\choice 
		 {\True $-1$}
		{$-4$} 
		{$-\infty$} 
		 {$+\infty$}
	\loigiai{ 
	Ta có $\lim\limits_{x\to -\infty}\dfrac{1}{x}\sqrt{x^2+16}=\lim\limits_{x\to -\infty}\dfrac{1}{x}\cdot  |x|\sqrt{1+\dfrac{16}{x^2}}=\lim\limits_{x\to -\infty}\dfrac{1}{x}\cdot (-x)\sqrt{1+\dfrac{16}{x^2}}=\lim\limits_{x\to -\infty}-\sqrt{1+\dfrac{16}{x^2}}=-1$.}
\end{ex}

% nguồn: Ok_Toan11_TracNghiem_Lan2
\begin{ex}%[1D3H2-3]%[ma:578fc30a]
	Biết  $\lim\limits_{x\to 1}\dfrac{\sqrt[3]{x}}{\sqrt[3]{4x+4}+2}=\dfrac{a}{b}$. Khẳng định nào sau đây là đúng?
	 \choice
		 {$a+b=3$}
		 {\True $b=4a$}
		 {$4b=a$}
		 {$a\cdot b=1$}
	 \loigiai{
		Ta có  $\lim\limits_{x\to 1}\dfrac{\sqrt[3]{x}}{\sqrt[3]{4x+4}+2}
		=\dfrac{1}{4}\Rightarrow \heva{&a=1\\&b=4}$ nên $b=4a$. 
	} 
\end{ex}

% nguồn: Ok_Toan11_TracNghiem_Lan2
\begin{ex}%[1D3H2-3]%[ma:ff55605b]
	Biết rằng $\lim\limits_{x \to -\sqrt{3}} \dfrac{2x^3+6 \sqrt{3}}{3-x^2}=a\sqrt{a}$ với $a$ là số  nguyên dương. Tính $a^2$.
	 \choice
		{$49$} 
		{$25$} 
		 {$4$}
		 {\True $9$}
	 \loigiai{
		Ta có $\lim\limits_{x \to -\sqrt{3}} \dfrac{2x^3+6 \sqrt{3}}{3-x^2}=\lim\limits_{x \to -\sqrt{3}} \dfrac{2\left(x^3+3  \sqrt{3}\right)}{\left(\sqrt{3}-x\right) \left(\sqrt{3}+x\right)}=\lim\limits_{x \to -\sqrt{3}} \dfrac{2\left(x^2-\sqrt{3}x+3\right)}{\sqrt{3}-x}=3 \sqrt{3}$.\\
		Suy ra  $a=3$.\\
		Vậy  $a^2=9$.
	 }
\end{ex}

% nguồn: Ok_Toan11_TracNghiem_Lan2
\begin{ex}%[1D3H2-3]%[ma:86de6415]
	Tính giới hạn  sau: $\lim\limits_{x\to \tfrac{1}{2}}\dfrac{8x^3-1}{6x^2-5x+1}$.
	\choice 
		{\True  $6$}
		 {$8$}
		{$1$} 
		 {$10$}
	 \loigiai{Ta có
		\begin{align*}
			\lim\limits_{x\to \tfrac{1}{2}}\dfrac{8x^3-1}{6x^2-5x+1}=\lim\limits_{x\to \tfrac{1}{2}}\dfrac{(2x-1)(4x^2+2x+1)}{(2x-1)(3x-1)}=\lim\limits_{x\to  \tfrac{1}{2}}\dfrac{4x^2+2x+1}{3x-1}=\dfrac{4\cdot\left(\dfrac{1}{2}\right)^2+2\cdot\dfrac{1}{2}+1}{3\cdot\dfrac{1}{2}-1}=6.
		\end{align*}
	 }
\end{ex}

% nguồn: Ok_Toan11_TracNghiem_Lan2
\begin{ex}%[1D3H2-3]%[ma:c72a967d]
	Giá trị  $\displaystyle\lim\limits _{x \to+\infty}\dfrac{2 x+8}{x-2}$ bằng
	\choice 
		{$-2$} 
		 {$4$}
		 {$-4$}
		{\True  $2$}
	\loigiai{Ta có $\displaystyle\lim\limits _{x \to+\infty}\dfrac{2 x+8}{x-2}=\displaystyle\lim\limits _{x  \to+\infty}\dfrac{x\left(2+\dfrac{8}{x}\right)}{x\left(1-\dfrac{2}{x}\right)}=\displaystyle\lim\limits _{x \to+\infty}\dfrac{2+\dfrac{8}{x}}{1-\dfrac{2}{x}}=2$.}
\end{ex}

% nguồn: Ok_Toan11_TracNghiem_Lan2
\begin{ex}%[1D3H2-3]%[ma:6c734ef1]
	Tìm $\lim\limits_{x\to  +\infty}(\sqrt{x^2+3x}-\sqrt{x^2+4x})$.
	 \choice
		{$\dfrac{7}{2}$} 
		{\True  $-\dfrac{1}{2}$}
		 {$+\infty$}
		 {$-\infty$}
	\loigiai{ 
	$\lim\limits_{x\to +\infty}(\sqrt{x^2+3x}-\sqrt{x^2+4x})=$  $\lim\limits_{x\to +\infty}\dfrac{-x}{\sqrt{x^2+3x}+\sqrt{x^2+4x}}=\lim\limits_{x\to +\infty}\dfrac{-1}{\sqrt{1+\dfrac{3}{x}}+\sqrt{1+\dfrac{4}{x}}}=-\dfrac{1}{2}$.}
\end{ex}

% nguồn: Ok_Toan11_TracNghiem_Lan2
\begin{ex}%[1D3H2-3]%[ma:804731ca]
	 Giá trị $\displaystyle\lim\limits _{x \to-\infty}\dfrac{1-x}{3 x+2}$ bằng
	 \choice
		{$\dfrac{1}{3}$} 
		 {$\dfrac{1}{2}$}
		 {\True $-\dfrac{1}{3}$}
		{$-\dfrac{1}{2}$} 
	\loigiai{ 
	 Ta có $\displaystyle\lim\limits _{x \to-\infty}\dfrac{1-x}{3 x+2}=\displaystyle\lim\limits _{x \to-\infty}\dfrac{\dfrac{1}{x}-1}{3+\dfrac{2}{x}}=-\dfrac{1}{3}$.}
\end{ex}

% nguồn: Ok_Toan11_TracNghiem_Lan2
\begin{ex}%[1D3H2-3]%[ma:315c8885]
	Biết rằng  $\lim\limits_{x\to +\infty}\left(\dfrac{x^2+1}{x-2}+ax-b\right)=-5$. Tính tổng $a+b$.
	\choice 
		{\True 6} 
		{7} 
		 {8}
		{5} 
	\loigiai{ 
		$\lim\limits_{x\to  +\infty}\left(\dfrac{x^2+1}{x-2}+ax-b\right)=\lim\limits_{x\to +\infty}\left(\dfrac{\left(a+1\right)x^2-(2a+b)x+2b+1}{x-2}\right)=-5$.\\
		Khi $a+1>0\Leftrightarrow a>-1$, ta  được $\lim\limits_{x\to +\infty}\left(\dfrac{\left(a+1\right)x^2-(2a+b)x+2b+1}{x-2}\right)=+\infty$ (không thỏa mãn).\\
		Khi $a+1<0\Leftrightarrow a<-1$,  ta được $\lim\limits_{x\to +\infty}\left(\dfrac{\left(a+1\right)x^2-(2a+b)x+2b+1}{x-2}\right)=-\infty$ (không thỏa mãn).\\
		Khi $a+1=0\Leftrightarrow  a=-1$, ta được\\ \hspace*{3cm}
		$\lim\limits_{x\to +\infty}\left(\dfrac{\left(a+1\right)x^2-(2a+b)x+2b+1}{x-2}\right)=\lim\limits_{x\to +\infty}\left(\dfrac{-\left(-2+b\right)x+2b+1}{x-2}\right)=2-b$.\\ 
	Do đó, $2-b=-5\Leftrightarrow b=7$.\\ Vậy  $a+b=6$}
\end{ex}

% nguồn: Ok_Toan11_TracNghiem_Lan2
\begin{ex}%[1D3H2-3]%[ma:db35fd69]
	Giới hạn $\lim\limits_{x\to 2}\dfrac{x^2-5x+6}{x^3-x^2-x-2}$  bằng
	\choice 
		{$0$} 
		{\True  $-\dfrac{1}{7}$}
		 {$-7$}
		 {$+\infty$}
	 \loigiai{
		 Ta có $\lim\limits_{x\to 2}\dfrac{x^2-5x+6}{x^3-x^2-x-2}=\lim\limits_{x\to 2}\dfrac{(x-2)(x-3)}{(x-2)(x^2+x+1)}=\lim\limits_{x\to 2}\dfrac{x-3}{x^2+x+1}=\dfrac{2-3}{2^2+2+1}=-\dfrac{1}{7}$.
	} 
\end{ex}

% nguồn: Ok_Toan11_TracNghiem_Lan2
\begin{ex}%[1D3H2-3]%[ma:9fe53bca]
	$\lim\limits_{x\to -\infty}\dfrac{\sqrt{4x^2-2x+1}+2-x}{\sqrt{9x^2-3x}+2x}$ có kết  quả là
	 \choice
		 {$-\dfrac{1}{5}$}
		 {\True 3}
		 {$-3$}
		 {$\dfrac{1}{5}$}
	\loigiai{ 
		Ta có $\lim\limits_{x\to -\infty}\dfrac{\sqrt{4x^2-2x+1}+2-x}{\sqrt{9x^2-3x}+2x} 
		 =\lim\limits_{x\to -\infty}\dfrac{-\sqrt{4-\dfrac{2}{x}+\dfrac{1}{x^2}}+\dfrac{2}{x}-1}{-\sqrt{9-\dfrac{3}{x}}+2}=\dfrac{-3}{-1}=3$.
	 }
\end{ex}

% nguồn: Ok_Toan11_TracNghiem_Lan2
\begin{ex}%[1D3H2-3]%[ma:a1cde4b5]
	Tính  $\lim\limits_{x\to +\infty}\dfrac{\sqrt{x+1}}{x^{2018}-1}$.
	 \choice
		{$-1$} 
		{$1$} 
		 {$2$}
		{\True $0$} 
	\loigiai{ 
	 $\lim\limits_{x\to +\infty}\dfrac{\sqrt{x+1}}{x^{2018}-1}=\lim\limits_{x\to +\infty}\dfrac{1}{x^{2017}}\cdot \dfrac{\sqrt{\dfrac{1}{x}+\dfrac{1}{x^2}}}{1-\dfrac{1}{x^{2017}}}=0$.}
\end{ex}

% nguồn: Ok_Toan11_TracNghiem_Lan2
\begin{ex}%[1D3H2-3]%[ma:3042370e]
	 Kết quả của $\lim\limits_{x\to +\infty}(1-2\mathrm{x})\sqrt{\dfrac{3\mathrm{x}-11}{x^3-1}}$ bằng
	\choice 
		 {\True $-2\sqrt{3}$}
		 {$2\sqrt{3}$}
		{$-\infty$} 
		{0} 
	\loigiai{ 
		Khi $x\to +\infty$ thì $1-2x<0$  nên ta có $\lim\limits_{x\to \to +\infty}(1-2x)\sqrt{\dfrac{3x-11}{x^3-1}}$ $=\lim\limits_{x\to \to +\infty}-\sqrt{\dfrac{(1-2x)^2(3x-11)}{x^3-1}}$\\ \hspace*{6cm}
	$=\lim\limits_{x\to  \to +\infty}-\sqrt{\dfrac{\left(\dfrac{1}{x^2}-2\right)^2\left(3-\dfrac{11}{x}\right)}{1-\dfrac{1}{x^3}}}=-\sqrt{12}=-2\sqrt{3}$.}
\end{ex}

% nguồn: Ok_Toan11_TracNghiem_Lan2
\begin{ex}%[1D3H2-3]%[ma:a3e6bffb]
	 Tìm tất cả các giá trị của $a$ để $\lim\limits_{x\to 2}(\sqrt{2x^2+1}+ax)$ là $0$.
	 \choice
		 {$-\dfrac{2}{3}$}
		 {\True $-\dfrac{9}{6}$}
		 {$0$}
		{Không có giá trị  $a$}
	 \loigiai{
		Ta có $\lim\limits_{x\to  2}(\sqrt{2x^2+1}+ax)=3+2a=0
		 \Leftrightarrow a=-\dfrac{3}{2}$.
	 }
\end{ex}

% nguồn: Ok_Toan11_TracNghiem_Lan2
\begin{ex}%[1D3H2-3]%[ma:8681aa06]
	Tính $\lim\limits_{x \to 1} \dfrac{(x-1)(x-3)}{1-x^2}$. 
	 \choice
		 {$0$}
		{$-1$} 
		 {\True $1$}
		{$-2$} 
	\loigiai{ 
		Ta  có $\lim\limits_{x \to 1} \dfrac{(x-1)(x-3)}{1-x^2}=\lim\limits_{x \to 1} \dfrac{(x-1)(x-3)}{(1-x)(1+x)}=\lim\limits_{x \to 1} \dfrac{-(x-3)}{(1+x)}=1$.
	 }
\end{ex}

% nguồn: Ok_Toan11_TracNghiem_Lan2
\begin{ex}%[1D3H2-3]%[ma:8608c4d5]
	Biết $\lim\limits_{x\to -\tfrac{1}{2} } \dfrac{2x^2-5x-3}{2x^2+11x+5}=-\dfrac{a}{b}$;\, $a,\,b\in \mathbb{Z}; \,b\ne 0$. Giá trị nhỏ nhất  của $a\cdot b$ bằng
	\choice 
		{\True $63$} 
		 {$16$}
		{$2$} 
		 {$-2$}
	\loigiai{ 
		Ta có $\lim\limits_{x\to -\tfrac{1}{2} } \dfrac{2x^2-5x-3}{2x^2+11x+5}=\lim\limits_{x\to  -\tfrac{1}{2} } \dfrac{\left(2x+1\right)\left(x-3\right)}{\left(2x+1\right)\left(x+5\right)}=\lim\limits_{x\to -\tfrac{1}{2} } \dfrac{x-3}{x+5}=-\dfrac{7}{9}$.\\
		Do đó, giá trị nhỏ nhất của $a\cdot b$ bằng $63$. 
	 }
\end{ex}

% nguồn: Ok_Toan11_TracNghiem_Lan2
\begin{ex}%[1D3H2-3]%[ma:7ec0bc88]
	Kết quả của $\lim\limits_{x\to{2^{+}}}(x-2)\sqrt{\dfrac{x}{x^2-4}}$  bằng
	\choice 
		{$+\infty$} 
		{\True $0$} 
		 {$-\infty$}
		 {1}
	 \loigiai{
	 Ta có $\lim\limits_{x\to{2^{+}}}(x-2)\sqrt{\dfrac{x}{x^2-4}}=\lim\limits_{x\to{2^{+}}}(x-2)\sqrt{\dfrac{x}{(x-2)(x+2)}}=\lim\limits_{x\to{2^{+}}}\sqrt{x-2}\sqrt{\dfrac{x}{x+2}}=0$.}
\end{ex}

% nguồn: Ok_Toan11_TracNghiem_Lan2
\begin{ex}%[1D3H2-3]%[ma:3a3b10ba]
	Chọn kết quả đúng của $\displaystyle\lim\limits _{x \to+\infty}\dfrac{1+3 x}{\sqrt{2 x^{2}+3}}$. 
	 \choice
		 {$-\dfrac{3 \sqrt{2}}{2}$}
		 {$-\dfrac{\sqrt{2}}{2}$}
		{\True  $\dfrac{3 \sqrt{2}}{2}$}
		{$\dfrac{\sqrt{2}}{2}$} 
	\loigiai{ 
	Ta có $\displaystyle\lim\limits _{x \to+\infty}\dfrac{1+3 x}{\sqrt{2  x^{2}+3}}=\displaystyle\lim\limits _{x \to+\infty}\dfrac{x\left(\dfrac{1}{x}+3\right)}{|x| \sqrt{2+\dfrac{3}{x^{2}}}}=\displaystyle\lim\limits _{x \to+\infty}\dfrac{\dfrac{1}{x}+3}{\sqrt{2+\dfrac{3}{x^{2}}}}=\dfrac{3}{\sqrt{2}}=\dfrac{3 \sqrt{2}}{2}$.}
\end{ex}

% nguồn: Ok_Toan11_TracNghiem_Lan2
\begin{ex}%[1D3H2-3]%[ma:b83bae86]
	Giới hạn $T=\lim\limits_{x\to  1}\dfrac{x^2-3x+2}{x^3+2x-3}$ bằng
	\choice 
		{$\dfrac{2}{9}$} 
		 {$\dfrac{2}{5}$}
		 {\True $-\dfrac{1}{5}$}
		 {$+\infty$}
	\loigiai{Ta có 
		$T=\lim\limits_{x\to  1}\dfrac{x^2-3x+2}{x^3+2x-3}=\lim\limits_{x\to 1}\dfrac{(x-1)(x-2)}{(x-1)(x^2+x+3)}=\lim\limits_{x\to 1}\dfrac{x-2}{x^2+x+3}=\dfrac{1-2}{1+1+3}=-\dfrac{1}{5}$.
	 }
\end{ex}

% nguồn: Ok_Toan11_TracNghiem_Lan2
\begin{ex}%[1D3H2-3]%[ma:600c10e1]
	Giới hạn $\displaystyle\lim\limits _{x \to+\infty}\dfrac{2 x^{5}-3 x^{3}+1}{4 x^{3}-2 x^{4}-x^{5}-3}$ bằng 
	\choice 
		{\True  $-2$}
		 {$\dfrac{1}{2}$}
		 {$-3$}
		 {$\dfrac{3}{2}$}
	\loigiai{ 
	Ta có $\displaystyle\lim\limits _{x \to+\infty}\dfrac{2  x^{5}-3 x^{3}+1}{4 x^{3}-2 x^{4}-x^{5}-3}=\displaystyle\lim\limits _{x \to+\infty}\dfrac{2-\dfrac{3}{x^{2}}+\dfrac{1}{x^{5}}}{\dfrac{4}{x^{2}}-\dfrac{2}{x}-1-\dfrac{3}{x^{5}}}=-2$.}
\end{ex}

% nguồn: Ok_Toan11_TracNghiem_Lan2
\begin{ex}%[1D3H2-3]%[ma:860e9e0f]
	Biết $\lim\limits_{x\to +\infty}\left(\sqrt{4x^2-3x+1}-\left(ax+b\right)\right)=0$.  Tính $a-4b$.
	\choice 
		{3} 
		 {\True 5}
		 {$-1$}
		 {2}
	 \loigiai{
		 Ta có $L=\lim\limits_{x\to +\infty}\left(\sqrt{4x^2-3x+1}-\left(ax+b\right)\right)=\lim\limits_{x\to +\infty}x\left(\sqrt{4-\dfrac{3}{x}+\dfrac{1}{x^2}}-a-\dfrac{b}{x}\right)$.\\
		Nếu $a>2$ thì $L=-\infty$ vì $\lim\limits_{x\to +\infty}x=+\infty$  và $\lim\limits_{x\to +\infty}\left(\sqrt{4-\dfrac{3}{x}+\dfrac{1}{x^2}}-a-\dfrac{b}{x}\right)=2-a<0$.\\
		 Tương tự nếu $a<2$ thì $L=+\infty$.\\
		Từ đó suy ra $a=2$.\\ Khi đó $L=\lim\limits_{x\to +\infty}\left(\sqrt{4x^2-3x+1}-\left(2x+b\right)\right)=\lim\limits_{x\to  +\infty}\dfrac{(-3-4b)x+1-b^2}{\sqrt{4x^2-3x+1}+2x+b}$\\
		 \hspace*{1cm}
		 $=\lim\limits_{x\to +\infty}\dfrac{-3-4b+\dfrac{1-b^2}{x}}{\sqrt{4-\dfrac{3}{x}+\dfrac{1}{x^2}}+2+\dfrac{b}{x}}=\dfrac{-3-4b}{4}$.\\ Suy ra $-3-4b=0\Leftrightarrow b=-\dfrac{3}{4}$.\\
	 Vậy $a-4b=5$.}
\end{ex}

% nguồn: Ok_Toan11_TracNghiem_Lan2
\begin{ex}%[1D3H2-3]%[ma:662041ac]
	Tính  $\lim\limits_{x\to +\infty}\dfrac{\sin x+1}{x}$ bằng
	\choice 
		{$+\infty $} 
		{$1$} 
		 {$-\infty $}
		 {\True $0$}
	 \loigiai{
	$\lim\limits_{x\to +\infty}\dfrac{\sin x+1}{x}=\lim\limits_{x\to +\infty}\dfrac{\sin x}{x}+\lim\limits_{x\to +\infty}\dfrac{1}{x}=0+0=0$.} 
\end{ex}

% nguồn: Ok_Toan11_TracNghiem_Lan2
\begin{ex}%[1D3H2-3]%[ma:71c0c88e]
	Cho hàm số $f(x)=\dfrac{(4x+1)^3(2x+1)^4}{(3+2x)^7}$.  Tính $\lim\limits_{x\to-\infty} f(x)$.
	\choice 
		 {$2$}
		{\True $8$} 
		{$4$} 
		{$0$} 
	 \loigiai{
	\[\lim\limits_{x\to-\infty}  f(x)=\lim\limits_{x\to-\infty} \dfrac{(4 x+1)^3(2 x+1)^4}{(3+2 x)^7}=\lim\limits_{x\to-\infty} \dfrac{\left(4+\dfrac{1}{x}\right)^3\left(2+\dfrac{1}{x}\right)^4}{\left(\dfrac{3}{x}+2\right)^7}=2^3=8.\]
	} 
\end{ex}

% nguồn: Ok_Toan11_TracNghiem_Lan2
\begin{ex}%[1D3H2-3]%[ma:c0fc6295]
	Tính  $\lim\limits_{x\to{1^{+}}}(x^3-1)\sqrt{\dfrac{x}{x^2-1}}$.
	 \choice
		 {$+\infty$}
		{\True $0$} 
		{$-\infty$} 
		{$\dfrac{3}{2}$} 
	 \loigiai{
		Ta có  $\lim\limits_{x\to{1^{+}}}(x^3-1)\cdot \sqrt{\dfrac{x}{x^2-1}}=\lim\limits_{x\to{1^{+}}}\sqrt{\dfrac{(x^3-1)^2x}{x^2-1}}$\\ \hspace*{1cm}
	$=\lim\limits_{x\to{1^{+}}}\sqrt{\dfrac{(x-1)^2(x^2+x+1)^2x}{(x-1)(x+1)}}=\lim\limits_{x\to{1^{+}}}\sqrt{\dfrac{(x-1)(x^2+x+1)^2x}{(x+1)}}=0$.} 
\end{ex}

% nguồn: Ok_Toan11_TracNghiem_Lan2
\begin{ex}%[1D3H2-3]%[ma:68c5378e]
	Tính ${\lim\limits_{x\to \sqrt{3}}} \dfrac{2x^2-6}{x-\sqrt{3}}=a\sqrt{b}$ ($a$, $b$ nguyên, $a>3$). Khi đó giá trị của $P=a+b$  bằng
	 \choice
		 {\True $7$}
		{$10$} 
		{$5$} 
		 {$6$}
	\loigiai{ 
		 Ta có ${\lim\limits_{x\to \sqrt{3}}} \dfrac{2x^2-6}{x-\sqrt{3}}={\lim\limits_{x\to \sqrt{3}}} \dfrac{2\left(x^2-3\right)}{x-\sqrt{3}}={\lim\limits_{x\to \sqrt{3}}} 2\left(x+\sqrt{3} \right)=4\sqrt{3}$.\\
		 Suy ra $a=4$, $b=3$. Vậy $P=a+b=7$.
	} 
\end{ex}

% nguồn: Ok_Toan11_TracNghiem_Lan2
\begin{ex}%[1D3H2-3]%[ma:84ffea2a]
	 Giá trị ${\lim\limits_{x\to-1}} \dfrac{x^2-1}{x+1}$ bằng
	 \choice
		{$2$} 
		 {$1$}
		{$0$} 
		{\True $-2$} 
	\loigiai{ 
		Ta có $\lim\limits_{x\to-1} \dfrac{x^2-1}{x+1}=\lim\limits_{x\to-1} \dfrac{\left(x+1\right)\left(x-1\right)}{x+1}=\lim\limits_{x\to-1} \left(x-1\right)=-2$. 
	} 
\end{ex}

% nguồn: Ok_Toan11_TracNghiem_Lan2
\begin{ex}%[1D3H2-3]%[ma:01bd37b1]
	$\lim\limits_{x\to -\infty}\dfrac{2x-3}{\sqrt{x^2+1}-x}$ có  kết quả là
	\choice 
		 {$-2$}
		 {$+\infty$}
		 {$3$}
		 {\True $-1$}
	\loigiai{ 
		Ta có $\lim\limits_{x\to -\infty}\dfrac{2x-3}{\sqrt{x^2+1}-x} 
		 =\lim\limits_{x\to -\infty}\dfrac{2-\dfrac{3}{x}}{-\sqrt{1+\dfrac{1}{x^2}}-1} =-1$.
	} 
\end{ex}

% nguồn: Ok_Toan11_TracNghiem_Lan2
\begin{ex}%[1D3H2-3]%[ma:8a076530]
	 Tìm $\lim\limits_{x\to a} \dfrac{x^3-\left(1+a^2\right)x+a}{x^3-a^3}$.
	\choice 
		{$\dfrac{2a^2}{a^2+3}$} 
		 {\True $\dfrac{2a^2-1}{3a^2}$}
		{$\dfrac{2}{3}$} 
		{$\dfrac{2a^2-1}{3}$} 
	 \loigiai{
		Ta có $\lim\limits_{x\to a} \dfrac{x^3-\left(1+a^2\right)x+a}{x^3-a^3}=\lim\limits_{x\to  a}\dfrac{x^3-a^2x-x+a}{\left(x-a\right)\left(x^2+ax+a^2\right)}=\lim\limits_{x\to a} \dfrac{x\left(x+a\right)-1}{x^2+ax+a^2}=\dfrac{2a^2-1}{3a^2}$.
	} 
\end{ex}

% nguồn: Ok_Toan11_TracNghiem_Lan2
\begin{ex}%[1D3H2-3]%[ma:533e1d04]
	Giá trị $\displaystyle\lim\limits  _{x \to-\infty}\dfrac{x}{x^{2}+1}$ bằng.
	\choice 
		{$-\infty$} 
		{$1$} 
		{$+\infty$} 
		{\True $0$} 
	 \loigiai{
	Ta có $\displaystyle\lim\limits _{x \to-\infty}\dfrac{x}{x^{2}+1}=\displaystyle\lim\limits _{x  \to-\infty}\dfrac{\dfrac{1}{x}}{1+\dfrac{1}{x^{2}}}=0$.}
\end{ex}

% nguồn: Ok_Toan11_TracNghiem_Lan2
\begin{ex}%[1D3H2-3]%[ma:bc1eefe4]
	Cho biết $\lim\limits_{x\to -\infty}(\sqrt{5x^2+2x}+x\sqrt{5})=a\sqrt{5}+b$ với $a$, $b\in \mathbb{Q}$. Tính $S=5a+b$. 
	 \choice
		{$-5$} 
		{\True $-1$} 
		 {$1$}
		 {$5$}
	\loigiai{ 
		$\lim\limits_{x\to -\infty}(\sqrt{5x^2+2x}+x\sqrt{5})=\lim\limits_{x\to -\infty}\dfrac{2x}{\sqrt{5x^2+2x}-x\sqrt{5}}=\lim\limits_{x\to  -\infty}\dfrac{2}{-\sqrt{5+\dfrac{2}{x}}-\sqrt{5}}$ $=-\dfrac{1}{5}\sqrt{5}$.\\
		Suy ra: $a=-\dfrac{1}{5}$, $b=0$. \\Vậy $S=-1$. 
	} 
\end{ex}

% nguồn: Ok_Toan11_TracNghiem_Lan2
\begin{ex}%[1D3H2-3]%[ma:d9b84cd4]
	Kết quả của $\lim\limits_{x \to  -1} \dfrac{3x^2-2x-5}{x^2-1}$ là
	 \choice
		 {$3$}
		{$+\infty$} 
		 {$0$}
		{\True  $4$}
	 \loigiai{
		Ta có $\lim\limits_{x \to  -1} \dfrac{3x^2-2x-5}{x^2-1}=\lim\limits_{x \to -1} \dfrac{(3x-5)(x+1)}{(x-1)(x+1)}=\lim\limits_{x \to -1} \dfrac{3x-5}{x-1}=4$.
	} 
\end{ex}

% nguồn: Ok_Toan11_TracNghiem_Lan2
\begin{ex}%[1D3H2-3]%[ma:9b36a29a]
	$\lim\limits_{x\to  1} \dfrac{x^2-3x+2}{x^3-1}$ bằng
	 \choice
		 {$\dfrac{1}{3}$}
		{\True $-\dfrac{1}{3}$} 
		{$0$} 
		 {$-\dfrac{2}{3}$}
	 \loigiai{
		Ta có $\lim\limits_{x\to 1} \dfrac{x^2-3x+2}{x^3-1}=\lim\limits_{x\to 1} \dfrac{\left(x-1\right)\left(x-2\right)}{\left(x-1\right)\left(x^2+x+1\right)}$ $=\lim\limits_{x\to 1} \dfrac{x-2}{x^2+x+1}=-\dfrac{1}{3}$. 
	 }
\end{ex}

% nguồn: Ok_Toan11_TracNghiem_Lan2
\begin{ex}%[1D3H2-3]%[ma:06279b38]
	Kết quả của  $\lim\limits_{x\to +\infty}x\sqrt{\dfrac{2\mathrm{x}+1}{3\mathrm{x^3}+x^2+2}}$ bằng
	\choice 
		 {$\dfrac{2}{3}$}
		{\True $\dfrac{\sqrt{6}}{3}$} 
		 {$-\infty$}
		 {$+\infty$}
	 \loigiai{
	Ta có $\lim\limits_{x\to  +\infty}x\sqrt{\dfrac{2\mathrm{x}+1}{3\mathrm{x^3}+x^2+2}}=\lim\limits_{x\to +\infty}\sqrt{\dfrac{2\mathrm{x^3}+x^2}{3\mathrm{x^3}+x^2+2}}=\lim\limits_{x\to +\infty}\sqrt{\dfrac{2+\dfrac{1}{x}}{3+\dfrac{1}{x}+\dfrac{2}{x^2}}}=\dfrac{\sqrt{6}}{3}$.}
\end{ex}

% nguồn: Ok_Toan11_TracNghiem_Lan2
\begin{ex}%[1D3H2-3]%[ma:658906f4]
	Biết $\lim\limits_{x\to 2} \dfrac{x^2+3x-10}{x^2+2x-3}=\dfrac{a}{b}$, $a,b\in \mathbb{Z};\, b\ne  0$. Giá trị nhỏ nhất của $a\cdot b$ bằng
	 \choice
		{$-10$} 
		{$10$} 
		{$-15$} 
		{\True  $7$}
	\loigiai{ 
		 Ta có $\lim\limits_{x\to 2} \dfrac{x^2+3x-10}{x^2-3x+2}=\lim\limits_{x\to 2} \dfrac{\left(x-2\right)\left(x+5\right)}{\left(x-2\right)\left(x-1\right)}=\lim\limits_{x\to 2} \dfrac{x+5}{x-1}=7$.\\
		Theo  giả thiết thì $\dfrac{a}{b}=7\Rightarrow a=7b$.\\
		Lúc  đó $a\cdot b=7b^2$ đạt giá trị nhỏ nhất bằng $7$ khi $b=\pm 1$.
	} 
\end{ex}

% nguồn: Ok_Toan11_TracNghiem_Lan2
\begin{ex}%[1D3H2-3]%[ma:adec1cf8]
	Có bao nhiêu giá trị nguyên của tham  số thực $m$ để $I<12$ biết $I=\lim\limits_{x\to -1}(x^4-2mx+m^2+3)$
	 \choice
		 {6}
		{\True 5} 
		{8} 
		 {7}
	 \loigiai{
		 Ta có $I=\lim\limits_{x\to -1}(x^4-2mx+m^2+3)=(-1)^4-2m(-1)+m^2+3=m^2+2m+4$.\\
		Khi đó  $I<12\Leftrightarrow m^2+2m+4<12\Leftrightarrow m^2+2m-8<0\Leftrightarrow -4<m<2$.\\
	Do $m$  nguyên nên $m\in \left\{-3,-2,-1,0,1\right\}$.}
\end{ex}

% nguồn: Ok_Toan11_TracNghiem_Lan2
\begin{ex}%[1D3H2-3]%[ma:d00a68e5]
	Tính giới hạn  $K=\lim\limits_{x\to -\infty}\dfrac{\sqrt{4x^2+1}}{x+1}$.
	 \choice
		{$K=0$} 
		{$K=1$} 
		{\True  $K=-2$}
		 {$K=4$}
	\loigiai{ 
	Ta  có $K=\lim\limits_{x\to -\infty}\dfrac{\sqrt{4x^2+1}}{x+1}=\lim\limits_{x\to -\infty}\dfrac{-x\sqrt{4+\dfrac{1}{x^2}}}{x+1}=\lim\limits_{x\to -\infty}\dfrac{-\sqrt{4+\dfrac{1}{x^2}}}{1+\dfrac{1}{x}}=-2$.}
\end{ex}

% nguồn: Ok_Toan11_TracNghiem_Lan2
\begin{ex}%[1D3H2-3]%[ma:688eb916]
	Giới hạn $\displaystyle\lim\limits _{x  \to+\infty}\dfrac{\sqrt{x^{2}+2}-2}{x-2}$ bằng
	\choice 
		 {$-\infty$}
		 {\True $1$}
		{$+\infty$} 
		{$-1$} 
	\loigiai{ 
	 Ta có $\displaystyle\lim\limits _{x \to+\infty}\dfrac{\sqrt{x^{2}+2}-2}{x-2}=\displaystyle\lim\limits _{x \to+\infty}\dfrac{x \sqrt{1+\dfrac{2}{x^{2}}}-2}{x-2}=\displaystyle\lim\limits _{x \to+\infty}\dfrac{\sqrt{1+\dfrac{2}{x^{2}}}-\dfrac{2}{x}}{1-\dfrac{2}{x}}=1$.}
\end{ex}

% nguồn: Ok_Toan11_TracNghiem_Lan2
\begin{ex}%[1D3H2-3]%[ma:2e0d0212]
	Biết $\lim\limits_{x \to 1} \dfrac{x^2-2x+1}{x^2-1}=a$. Hỏi $a$ không là  nghiệm của bất phương trình nào sau đây?
	 \choice
		 {$x^2+x+1>0$}
		 {$-2x+1>0$}
		 {\True $x^2-5x+6<0$}
		 {$x^2-3x \geq 0$}
	 \loigiai{
		Ta có  $\lim\limits_{x \to 1} \dfrac{x^2-2x+1}{x^2-1}=\lim\limits_{x \to 1} \dfrac{x-1}{x+1}=0$ suy ra $a=0$.\\
		Thay vào các bất  phương trình đã cho ta thấy $a=0$ không phải là nghiệm của bất phương trình $x^2-5x+6<0$.
	} 
\end{ex}

% nguồn: Ok_Toan11_TracNghiem_Lan2
\begin{ex}%[1D3H2-3]%[ma:8fcd4a81]
	Giới hạn $\lim\limits_{x\to 1}\dfrac{x^2-3x+2}{x^3-x^2+x-1}$  bằng
	 \choice
		{$-2$} 
		 {$-1$}
		{\True  $-\dfrac{1}{2}$}
		{$\dfrac{1}{2}$} 
	 \loigiai{
		Ta có $\lim\limits_{x\to 1}\dfrac{x^2-3x+2}{x^3-x^2+x-1}=\lim\limits_{x\to 1}\dfrac{(x-1)(x-2)}{(x-1)(x^2+1)}=\lim\limits_{x\to  1}\dfrac{x-2}{x^2+1}=\dfrac{1-2}{1+1}=-\dfrac{1}{2}$.
	 }
\end{ex}

% nguồn: Ok_Toan11_TracNghiem_Lan2
\begin{ex}%[1D3H2-3]%[ma:4a7f8266]
	Biết $\lim\limits_{x\to -\infty}\dfrac{ax-\sqrt{4x^2-x+1}}{a^2x-1}=\dfrac{2}{3}$ với $a=\dfrac{m+\sqrt{57}}{n}>0$ và $m; n$ là số nguyên dương. Khi đó  giá trị của biểu thức $S=m+n$ là
	\choice 
		 {$5$}
		 {$6$}
		{\True $7$} 
		 {$8$}
	\loigiai{ 
		 Ta có
		\begin{eqnarray*}
			\lim\limits_{x\to -\infty}\dfrac{ax-\sqrt{4x^2-x+1}}{a^2x-1}=\dfrac{2}{3} 
			&\Leftrightarrow& \lim\limits_{x\to  -\infty}\dfrac{x\left(a+\sqrt{4-\dfrac{1}{x}+\dfrac{1}{x^2}}\right)}{x\left(a^2-\dfrac{1}{x}\right)}=\dfrac{2}{3} \\
			&\Leftrightarrow& \dfrac{a+2}{a^2}=\dfrac{2}{3} 
			\Leftrightarrow 2a^2-3a-6=0 
			\Leftrightarrow \hoac{&a=\dfrac{3+\sqrt{57}}{4} \text{ (nhận)}  \\&a=\dfrac{3-\sqrt{57}}{4}\text{ (loại)}.}
		\end{eqnarray*}
		Suy ra $\heva{&m=3 \\&n=4}\Leftrightarrow S=7$. Vậy $S=7$. 
	 }
\end{ex}

% nguồn: Ok_Toan11_TracNghiem_Lan2
\begin{ex}%[1D3H2-3]%[ma:4d5c7136]
	 Kết quả của $\lim\limits_{x \to -1} \dfrac{x^3+1}{x^2+x}$ bằng
	 \choice
		 {\True $-3$}
		{$-1$} 
		 {$0$}
		 {$1$}
	 \loigiai{
		Ta có $\lim\limits_{x \to -1} \dfrac{x^3+1}{x^2+x}=\lim\limits_{x \to -1} \dfrac{(x+1)\left(x^2-x+1\right)}{x(x+1)}=\lim\limits_{x \to  -1} \dfrac{x^2-x+1}{x}=\dfrac{3}{-1}=-3$.
	 }
\end{ex}

% nguồn: Ok_Toan11_TracNghiem_Lan2
\begin{ex}%[1D3H2-3]%[ma:62d41739]
	 Tính
	 $L=\lim\limits_{x\to-\infty} \dfrac{(2x-1)(3x+2)(4 x-5)}{8 x^3+2 x+7}$.
	 \choice
		{\True  $L=3$}
		 {$L=\dfrac{3}{4}$}
		 {$L=\dfrac{4}{3}$}
		{$L=-3$} 
	 \loigiai{
	\[L=\lim\limits_{x\to-\infty} \dfrac{x^3\left(2-\dfrac{1}{x}\right)\left(3+\dfrac{2}{x}\right)\left(4-\dfrac{5}{x}\right)}{x^3\left(8+\dfrac{2}{x^2}+\dfrac{7}{x^3}\right)}=\lim\limits_{x\to-\infty} \dfrac{\left(2-\dfrac{1}{x}\right)\left(3+\dfrac{2}{x}\right)\left(4-\dfrac{5}{x}\right)}{8+\dfrac{2}{x^2}+\dfrac{7}{x^3}}=3.\] 
	} 
\end{ex}

% nguồn: Ok_Toan11_TracNghiem_Lan2
\begin{ex}%[1D3H2-3]%[ma:dc6ac5a2]
	$\lim\limits_{x\to  -\infty}\dfrac{2x^2+5x-3}{x^2+6x+3}$ có kết quả là
	 \choice
		 {$-2$}
		{$+\infty$} 
		{$3$} 
		 {\True $2$}
	\loigiai{ 
		 Ta có $\lim\limits_{x\to -\infty}\dfrac{2x^2+5x-3}{x^2+6x+3}
		=\lim\limits_{x\to  +\infty}\dfrac{2+\dfrac{5}{x}-\dfrac{3}{x^2}}{1+\dfrac{6}{x}+\dfrac{3}{x^2}}
		 =2$.
	} 
\end{ex}

% nguồn: Ok_Toan11_TracNghiem_Lan2
\begin{ex}%[1D3H2-3]%[ma:d3ae4db2]
	 Tính $\lim\limits_{x\to -\infty}\dfrac{1}{x}(\sqrt{4x^2+1}-x)$
	\choice 
		 {$+\infty$}
		 {0}
		{$\dfrac{1}{2}$} 
		{\True  $-3$}
	\loigiai{ 
		Ta có $\lim\limits_{x\to -\infty}\dfrac{1}{x}(\sqrt{4x^2+1}-x)$ $=\lim\limits_{x\to  -\infty}\dfrac{1}{x}\left(\dfrac{4x^2+1-x^2}{\sqrt{4x^2+1}+x}\right)$\\\hspace*{1cm}
	$=\lim\limits_{x\to -\infty}\dfrac{3x^2+1}{x\cdot (\sqrt{4x^2+1}+x)}$  $=\lim\limits_{x\to -\infty}\dfrac{3+\dfrac{1}{x^2}}{1\cdot \left(-\sqrt{4+\dfrac{1}{x^2}}+1\right)}=-3$.}
\end{ex}

% nguồn: Ok_Toan11_TracNghiem_Lan2
\begin{ex}%[1D3H2-3]%[ma:83796b36]
	Kết quả  $\lim\limits_{x \to 1} \dfrac{x^2-3x+2}{x-1}$ là
	\choice 
		{\True $-1$} 
		{$3$} 
		 {$0$}
		 {$+\infty$}
	 \loigiai{
		Ta có $\lim\limits_{x \to 1} \dfrac{x^2-3x+2}{x-1}=\lim\limits_{x \to  1} \dfrac{(x-1)(x-2)}{x-1}=\lim\limits_{x \to 1}(x-2)=-1$.
	 }
\end{ex}

% nguồn: Ok_Toan11_TracNghiem_Lan2
\begin{ex}%[1D3H2-3]%[ma:2cee0242]
	 Với $a$ là số thực khác $0$, $\lim\limits_{x\to a} \dfrac{x^2-\left(a+1\right)x+a}{x^2-a^2}$ bằng
	 \choice
		 {$a-1$}
		{$a+1$} 
		 {\True $\dfrac{a-1}{2a}$}
		{$\dfrac{a+1}{2a}$} 
	 \loigiai{
		Ta có $\lim\limits_{x\to a} \dfrac{x^2-\left(a+1\right)x+a}{x^2-a^2}=\lim\limits_{x\to  a} \dfrac{\left(x-1\right)\left(x-a\right)}{\left(x+a\right)\left(x-a\right)}=\lim\limits_{x\to a} \dfrac{x-1}{x+a}=\dfrac{a-1}{2a}$.
	 }
\end{ex}

% nguồn: Ok_TracNghiemLop11_Lan1
\begin{ex}%[1D3H2-3]%[ma:75f10c6c]
	Tìm  số thực $a$ khác $0$ sao cho $\lim\limits_{n \to +\infty}{\dfrac{n^2-2}{an^2-1}}=2$.
	 \choice
		 {$a=-\dfrac{1}{2}$}
		 {$a=-2$}
		 {$a=2$}
		{\True  $a=\dfrac{1}{2}$}
	\loigiai{ 
		Ta có: $\lim\limits_{n \to +\infty}{\dfrac{n^2-2}{a n^2-1}}=\lim\limits_{n \to +\infty}{\dfrac{n^2\left(\dfrac{n^2}{n^2}-\dfrac{2}{n^2}\right)}{n^2\left(\dfrac{a^2}{n^2}-\dfrac{1}{n^2}\right)} }=\lim\limits_{n \to +\infty}{\dfrac{1-\dfrac{2}{n^2}}{a-\dfrac{1}{n^2}} }=\dfrac{1}{a}$.\\ 
		Do đó,  $\dfrac{1}{a}=2 \Rightarrow a=\dfrac{1}{2}(\mathrm{tm})$.
	} 
\end{ex}

% nguồn: Ok_TracNghiemLop11_Lan1
\begin{ex}%[1D3H2-3]%[ma:3da20c73]
	Tính $\lim\limits_{x\to -\infty} \dfrac{-3x+1}{x-2}$. 
	 \choice
		{\True $-3$} 
		{$+\infty $} 
		 {$-\infty $}
		 {$-\dfrac{1}{2}$}
	 \loigiai{
		Ta có  $\lim\limits_{x\to -\infty} \dfrac{-3x+1}{x-2}=\lim\limits_{x\to -\infty} \dfrac{-3+\dfrac{1}{x}}{1-\dfrac{2}{x}}=\dfrac{-3}{1}=-3$.
	} 
\end{ex}

% nguồn: Ok_TracNghiemLop11_Lan1
\begin{ex}%[1D3H2-3]%[ma:285482f6]
	Tính $\lim\limits_{x\to 5}  \dfrac{x^2-12x+35}{25-5x}$.
	\choice 
		{  $-\dfrac{2}{5}$}
		{  $+\infty $}
		{\True  $\dfrac{2}{5}$}
		 { $-\infty $}
	\loigiai{ 
		Ta có $\lim\limits_{x\to 5} \dfrac{x^2-12x+35}{25-5x}=\lim\limits_{x\to  5} \dfrac{\left( x-7 \right)\left( x-5 \right)}{-5\left( x-5 \right)}=\lim\limits_{x\to 5} \dfrac{x-7}{-5}=\dfrac{2}{5}$.
	} 
\end{ex}

% nguồn: Ok_TracNghiemLop11_Lan1
\begin{ex}%[1D3H2-3]%[ma:338dd159]
	Với giá trị nào của tham số $m$ thì $\lim\limits_{x \to 2}\dfrac{x^2-3x+m}{x-2}  = 1$?
	 \choice
		 {$m=-1$}
		{\True $m=2$} 
		 {$m=-4$}
		{Không  có $m$ thỏa mãn}
	\loigiai{ 
	Vì giới hạn đã cho là hữu hạn nên $x^2 - 3x + m$ phải có nhân tử $x-2$  hay phương trình $x^2 - 3x + m = 0$ có một nghiệm $x = 2$. Suy ra $m = 2$.\\
	Thử lại ta nhận $m =  2$.
	} 
\end{ex}

% nguồn: Ok_Toan11_TracNghiem_Lan2
\begin{ex}%[1D3V2-3]%[ma:c38063e2]
	Biết $\lim\limits_{x\to 3} \dfrac{x+1-\sqrt{5x+1}}{x-\sqrt{4x-3}}=\dfrac{a}{b}$ trong đó $\dfrac{a}{b}$ là phân số tối giản, $a$, $b$ là các số nguyên dương.  Giá trị của $a-b$ bằng
	 \choice
		{$\dfrac{1}{9}$} 
		{$\dfrac{9}{8}$} 
		 {\True $1$}
		{$-1$} 
	\loigiai{ 
		Ta  có
		\begin{align*}
			\lim\limits_{x\to 3} \dfrac{x+1-\sqrt{5x+1}}{x-\sqrt{4x-3}}&=\lim\limits_{x\to 3}  \dfrac{\left[\left(x+1\right)^2-\left(5x+1\right)\right]\left[x+\sqrt{4x-3} \right]}{\left[x^2-\left(4x-3\right)\right]\left[x+1+\sqrt{5x+1} \right]}\\
			&=\lim\limits_{x\to  3} \dfrac{\left(x^2-3x\right)\left[x+\sqrt{4x-3} \right]}{\left(x^2-4x+3\right)\left[x+1+\sqrt{5x+1} \right]}\\
			&=\lim\limits_{x\to  3} \dfrac{x\left(x-3\right)\left[x+\sqrt{4x-3} \right]}{\left(x-1\right)\left(x-3\right)\left[x+1+\sqrt{5x+1} \right]}\\
			 &=\lim\limits_{x\to 3} \dfrac{x\left[x+\sqrt{4x-3} \right]}{\left(x-1\right)\left[x+1+\sqrt{5x+1} \right]}=\dfrac{3\cdot 6}{2\cdot 8}=\dfrac{9}{8}.
		\end{align*}
		 Do đó $a=9$, $b=8$. Vậy $a-b=1$.
	 }
\end{ex}

% nguồn: Ok_Toan11_TracNghiem_Lan2
\begin{ex}%[1D3V2-3]%[ma:72485f46]
	Tìm $\lim\limits_{x\to 1} \dfrac{\sqrt[3]{x+7}-\sqrt{x^2+x+2}}{x-1}$. 
	 \choice
		 {$\dfrac{1}{12}$}
		{$+\infty$} 
		{$-\dfrac{3}{2}$} 
		{\True  $-\dfrac{2}{3}$}
	\loigiai{ 
		 Ta có
		\begin{align*}
			\lim\limits_{x\to 1} \dfrac{\sqrt[3]{x+7}-\sqrt{x^2+x+2}}{x-1}&=\lim\limits_{x\to 1}  \left(\dfrac{\sqrt[3]{x+7}-2}{x-1}+\dfrac{2-\sqrt{x^2+x+2}}{x-1} \right)\\
			&=\lim\limits_{x\to 1} \left[\dfrac{x-1}{x-1}\cdot \dfrac{1}{\sqrt[3]{\left(x+7\right)^2}+2\sqrt[3]{x+7}+4}+\dfrac{2-x^2-x}{\left(x-1\right)\left(2+\sqrt{x^2+x+2} \right)}  \right]\\
			&=\lim\limits_{x\to 1}  \left[\dfrac{1}{\sqrt[3]{\left(x+7\right)^2}+2\sqrt[3]{x+7}+4}-\dfrac{x+2}{2+\sqrt{x^2+x+2}} \right]\\
			&=-\dfrac{2}{3}. 
		\end{align*}
	 }
\end{ex}

% nguồn: Ok_Toan11_TracNghiem_Lan2
\begin{ex}%[1D3V2-3]%[ma:e0ea0ab5]
	Cho $a\in\mathbb{R}$, $a\ne 0$. Khi đó $\lim\limits_{x\to+\infty}\left(\dfrac{x^2-2}{a x^2-1}\right)=3$ thì giá trị của $a$ bằng 
	\choice 
		 {$-1$}
		{$1$} 
		 {$2$}
		 {\True $\dfrac{1}{3}$}
	\loigiai{ 
	Ta có  $
	\lim\limits_{x\to+\infty}\left(\dfrac{x^2-2}{a  x^2-1}\right)=\lim\limits_{x\to+\infty}\left(\dfrac{1-\dfrac{2}{x^2}}{a-\dfrac{1}{x^2}}\right)=\dfrac{1}{a}=3\Leftrightarrow a=\dfrac{1}{3}$.
	} 
\end{ex}

% nguồn: Ok_Toan11_TracNghiem_Lan2
\begin{ex}%[1D3V2-3]%[ma:db38f484]
	Giá trị của $\displaystyle\lim\limits _{x \to-\infty}\dfrac{2  x-1}{\sqrt{x^{2}+1}-1}$ bằng
	\choice 
		 {$0$}
		{\True  $-2$}
		 {$-\infty$}
		 {$2$}
	\loigiai{ 
	Ta có $\displaystyle\lim\limits _{x \to-\infty}\dfrac{2 x-1}{\sqrt{x^{2}+1}-1}=\displaystyle\lim\limits _{x \to-\infty}\dfrac{2 x-1}{-x \sqrt{1+\dfrac{1}{x^{2}}}-1}=\displaystyle\lim\limits _{x \to-\infty}\dfrac{2-\dfrac{1}{x}}{-\sqrt{1+\dfrac{1}{x^{2}}}-\dfrac{1}{x}}=-2$.} 
\end{ex}

% nguồn: Ok_Toan11_TracNghiem_Lan2
\begin{ex}%[1D3V2-3]%[ma:8624e12a]
	 Cho $a$ và $b$ là các số thực thỏa mãn $\lim\limits_{x\to-1}\dfrac{x^2+ax+b}{x+1}=3$. Tính $a+b$.
	 \choice
		 {\True $9$}
		 {$6$}
		{$8$} 
		{$7$} 
	 \loigiai{
		Ta có $\lim\limits_{x\to-1}(x+1)=0$ và $\lim\limits_{x\to-1}\dfrac{x^2+ax+b}{x+1}=3$ nên $\lim\limits_{x\to-1}(x^2+ax+b)=0$.\\ 
		Từ đó ta có $1-a+b=0\Rightarrow b=a-1$ và  $x^2+ax+b=(x+1)(x+a-1)$.\\
		 Khi đó $\lim\limits_{x\to-1}\dfrac{x^2+ax+b}{x+1}=\lim\limits_{x\to-1}\dfrac{(x+1)(x+a-1)}{x+1}=\lim\limits_{x\to-1}(x+a-1)=a-2$.\\
		Suy ra $a-2=3\Leftrightarrow a=5, b=4$.\\ 
		Vậy  $a+b=9$.
	} 
\end{ex}

% nguồn: Ok_Toan11_TracNghiem_Lan2
\begin{ex}%[1D3V2-3]%[ma:6310f94d]
	Cho $f(x)$ là đa thức thỏa mãn $\lim\limits_{x\to 2} \dfrac{f(x)-20}{x-2}=10$.  Tìm $\lim\limits_{x\to 2} \dfrac{\sqrt[3]{6f(x)+5}-5}{x^2+x-6}$.
	\choice 
		 {$T=\dfrac{4}{15}$}
		 {$T=\dfrac{12}{25}$}
		{$T=\dfrac{6}{25}$} 
		 {\True $T=\dfrac{4}{25}$}
	 \loigiai{
		Vì $\lim\limits_{x\to 2} \dfrac{f(x)-20}{x-2}=10$ nên $f(x)\to 20$  khi $x\to 2$.\\
		Ta  có
		\begin{align*}
			\lim\limits_{x\to 2} \dfrac{\sqrt[3]{6f(x)+5}-5}{x^2+x-6}&=\lim\limits_{x\to 2} \dfrac{6f(x)+5-125}{\left(x-2\right)\left(x+3\right)\left[\sqrt[3]{\left(6f(x)+5\right)^2}+5\cdot  \sqrt[3]{6f(x)+5}+25\right]}\\
			&=\lim\limits_{x\to 2}  \left[\dfrac{f(x)-20}{x-2}\cdot\dfrac{6}{\left(x+3\right)\left[\sqrt[3]{\left(6f(x)+5\right)^2}+5\cdot \sqrt[3]{6f(x)+5}+25\right]} \right]\\
			&=10\cdot \dfrac{6}{5\cdot \left[\sqrt[3]{\left(6\cdot 20+5\right)^2}+5\cdot\sqrt[3]{6\cdot 20+5}+25\right]}=\dfrac{4}{25}. 
		\end{align*}
	} 
\end{ex}

% nguồn: Ok_Toan11_TracNghiem_Lan2
\begin{ex}%[1D3V2-3]%[ma:322d194d]
	 Biết rằng $b > 0$, $a+b=5$ và $\lim\limits_{x\to 0} \dfrac{\sqrt[3]{ax+1}-\sqrt{1-bx}}{x}=2$. Khẳng định nào dưới đây \textbf{sai}?
	\choice 
		{$a^2+b^2> 10$} 
		{\True $a^2-b^2> 6$} 
		{$a-b\ge 0$} 
		{$1\le  a\le 3$}
	\loigiai{ 
		\begin{align*}
			\lim\limits_{x\to 0} \dfrac{\sqrt[3]{ax+1}-\sqrt{1-bx}}{x}&=\lim\limits_{x\to 0} (\dfrac{\sqrt[3]{ax+1}-1}{x}+\dfrac{1-\sqrt{1-bx}}{x})\\ 
			 &=\lim\limits_{x\to 0} (\dfrac{a}{\sqrt[3]{(ax+1)^2}+\sqrt[3]{ax+1}+1}+\dfrac{b}{1+\sqrt{bx+1}})\\
			 &=\dfrac{a}{3}+\dfrac{b}{2}.
		\end{align*}
		 Từ giả thiết ta có
		 $\heva{&\dfrac{a}{3}+\dfrac{b}{2}=2\\&a+b=5}\Leftrightarrow \heva{&a=3\\&b=2.}$\\
		Đối chiếu với  các phương án thì $a^2-b^2=3^2-2^2=5> 6$ là khẳng định sai.
	} 
\end{ex}

% nguồn: Ok_Toan11_TracNghiem_Lan2
\begin{ex}%[1D3V2-3]%[ma:dd70ffd1]
	Giới  hạn $\lim\limits_{x\to-\infty}\dfrac{\sqrt{x^2+3}-x}{3x-2}$ bằng
	\choice 
		{$-\dfrac{1}{3}$} 
		{\True  $-\dfrac{2}{3}$}
		 {$+\infty$}
		{$0$} 
	\loigiai{ 
	Khi $x\to-\infty$  thì $x<0\Rightarrow|x|=-x$.
	\[\lim\limits_{x\to-\infty} \dfrac{\sqrt{x^2+3}-x}{3  x-2}=\lim\limits_{x\to-\infty} \dfrac{|x| \sqrt{1+\dfrac{3}{x^2}}-x}{x\left(3-\dfrac{2}{x}\right)}=\lim\limits_{x\to-\infty} \dfrac{-x \sqrt{1+\dfrac{3}{x^2}}-x}{x\left(3-\dfrac{2}{x}\right)}=\lim\limits_{x\to-\infty} \dfrac{-\sqrt{1+\dfrac{3}{x^2}}-1}{\left(3-\dfrac{2}{x}\right)}=-\dfrac{2}{3}.\]
	 }
\end{ex}

% nguồn: Ok_Toan11_TracNghiem_Lan2
\begin{ex}%[1D3V2-3]%[ma:835d0f95]
	Cho $\lim\limits_{x\to 0} \left(\dfrac{x}{\sqrt[7]{x+1}\cdot \sqrt{x+4}-2} \right)=\dfrac{a}{b}$ trong đó $\dfrac{a}{b}$ là phân số tối giản, $a$, $b$ là  các số nguyên dương. Tính tổng $T=a+b$.
	 \choice
		{\True $T=43$} 
		{$T=23$} 
		{$T=13$} 
		 {$T=53$}
	 \loigiai{
		Đặt $L=\lim\limits_{x\to 0} \left(\dfrac{x}{\sqrt[7]{x+1}\cdot \sqrt{x+4}-2} \right)=\dfrac{a}{b}$ thì  $\dfrac{1}{L}=\lim\left(\dfrac{\sqrt[7]{x+1}\cdot \sqrt{x+4}-2}{x} \right)=\dfrac{b}{a}$.\\
		Ta có 
		\begin{align*}
			\dfrac{b}{a}&=\lim\limits_{x\to 0} \left(\dfrac{\sqrt[7]{x+1}\cdot \sqrt{x+4}-\sqrt{x+4}+\sqrt{x+4}-2}{x} \right)\\ 
			&=\lim\limits_{x\to 0} \left(\dfrac{\sqrt[7]{x+1}\cdot \sqrt{x+4}-\sqrt{x+4}}{x} \right)+\lim\limits_{x\to 0}  \left(\dfrac{\sqrt{x+4}-2}{x} \right).
		\end{align*}
		\begin{itemize}
			\item Xét $L_1=\lim\limits_{x\to  0} \left(\dfrac{\sqrt{x+4} \left(\sqrt[7]{x+1}-1\right)}{x} \right)$.\\
			Đặt $t=\sqrt[7]{x+1}$. Khi đó  $\heva{&x=t^7-1\\&x\to 0\Rightarrow t\to 1.}$\\
			Do đó 
			 \[L_1=\lim\limits_{t\to 1} \dfrac{\sqrt{t^7+3} \left(t-1\right)}{t^7-1}=\lim\limits_{t\to 1} \dfrac{\sqrt{t^7+3}}{\left(t^6+t^5+t^4+t^3+t^2+t+1\right)}=\dfrac{2}{7}.\]
			\item Xét $L_2=\lim\limits_{x\to 0}  \left(\dfrac{\sqrt{x+4}-2}{x} \right)=\lim\limits_{x\to 0} \dfrac{\left(\sqrt{x+4}-2\right)\left(\sqrt{x+4}+2\right)}{x\left(\sqrt{x+4}+2\right)}=\lim\limits_{x\to 0} \dfrac{1}{\sqrt{x+4}+2}=\dfrac{1}{4}$.
		\end{itemize}
		Từ đó  $\dfrac{b}{a}=\dfrac{2}{7}+\dfrac{1}{4}=\dfrac{15}{28}$ suy ra $a=28$, $b=15$. Vậy $T=a+b=43$.
	 }
\end{ex}

% nguồn: Ok_Toan11_TracNghiem_Lan2
\begin{ex}%[1D3V2-3]%[ma:3b00316a]
	$\lim\limits_{x\to 0} \dfrac{\sqrt{x+1}-\sqrt{x^2+x+1}}{x}$ bằng 
	\choice 
		{\True  $0$}
		 {$-1$}
		{$5$} 
		{$1$} 
	\loigiai{ 
		Ta  có
		\begin{align*}
			\lim\limits_{x\to-1} \dfrac{\sqrt{x+1}-\sqrt{x^2+x+1}}{x}&=\lim\limits_{x\to  0} \dfrac{x+1-x^2-x-1}{x\left(\sqrt{x+1}+\sqrt{x^2+x+1} \right)}\\
			&=\lim\limits_{x\to  0} \dfrac{-x}{\sqrt{x+1}+\sqrt{x^2+x+1}}=\dfrac{0}{2}=0.
		\end{align*}
		Vậy  $\lim\limits_{x\to 0} \dfrac{\sqrt{x+1}-\sqrt{x^2+x+1}}{x}=0$.
	} 
\end{ex}

% nguồn: Ok_Toan11_TracNghiem_Lan2
\begin{ex}%[1D3V2-3]%[ma:6415d42d]
	Biết rằng $\lim\limits_{x\to 0} \dfrac{2\sqrt{1+x}-\sqrt[3]{8-x}}{x}=\dfrac{a}{b}$ trong đó $\dfrac{a}{b}$ là phân số tối  giản, $a$, $b$ là các số nguyên dương. Tính $a-b$.
	 \choice
		{$25$} 
		{\True $1$} 
		 {$-1$}
		{$-\dfrac{13}{12}$} 
	 \loigiai{
		Ta  có
		\begin{align*}
			\lim\limits_{x\to 0} \dfrac{2\sqrt{1+x}-\sqrt[3]{8-x}}{x}&=\lim\limits_{x\to 0}  \left(\dfrac{2\sqrt{1+x}-2}{x}+\dfrac{2-\sqrt[3]{8-x}}{x} \right)\\
			 &=\lim\limits_{x\to 0} \left(\dfrac{2}{\sqrt{x+1}+1}+\dfrac{1}{4+2\sqrt[3]{8-x}+\sqrt[3]{\left(8-x\right)^2}} \right)\\
			&=1+\dfrac{1}{12}=\dfrac{13}{12}. 
		\end{align*}
		Suy  ra $\dfrac{a}{b}=\dfrac{13}{12} \Rightarrow a-b=1$.
	} 
\end{ex}

% nguồn: Ok_Toan11_TracNghiem_Lan2
\begin{ex}%[1D3V2-3]%[ma:f649be2c]
	Tính  $\lim\limits_{x\to 1}\dfrac{x^2-(a+2)x+a+1}{x^3-1}$.
	\choice 
		 {$\dfrac{2-a}{3}$}
		{$\dfrac{-2-a}{3}$} 
		{\True  $-\dfrac{a}{3}$}
		{$\dfrac{a}{3}$} 
	 \loigiai{
		Ta có $\lim\limits_{x\to 1}\dfrac{x^2-(a+2)x+a+1}{x^3-1}=\lim\limits_{x\to 1}\dfrac{(x-1)^2-a(x-1)}{(x-1)(x^2+x+1)}=\lim\limits_{x\to  1}\dfrac{x-1-a}{x^2+x+1}=-\dfrac{a}{3}$.
	} 
\end{ex}

% nguồn: Ok_Toan11_TracNghiem_Lan2
\begin{ex}%[1D3V2-3]%[ma:cfcd84e3]
	Tìm  $\displaystyle\lim\limits _{x \to-\infty}\dfrac{\sqrt{x^{2}+3 x+5}}{4 x-1}$.
	\choice 
		{\True  $-\dfrac{1}{4}$}
		{$1$} 
		{$0$} 
		{$\dfrac{1}{4}$} 
	\loigiai{ 
	Ta có $\displaystyle\lim\limits _{x \to-\infty}\dfrac{\sqrt{x^{2}+3 x+5}}{4 x-1}=\displaystyle\lim\limits _{x \to-\infty}\dfrac{-\sqrt{1+\dfrac{3}{x}+\dfrac{5}{x^{2}}}}{4-\dfrac{1}{x}}=-\dfrac{1}{4}$.} 
\end{ex}

% nguồn: Ok_Toan11_TracNghiem_Lan2
\begin{ex}%[1D3V2-3]%[ma:760c6e3b]
	$\lim\limits_{x\to+\infty}\dfrac{2x^4-5x^3}{3x^4+3x+2}=\dfrac{a}{b}$ ($a$, $b$ là các số nguyên dương, $a$ và $b$ là các số nguyên tố cùng nhau). Khẳng  định nào \textbf{đúng}?
	\choice 
		{$a=-5$, $b=3$} 
		{$a=5$, $b=3$} 
		{$a=-2$,  $b=3$}
		 {\True $a=2$, $b=3$}
	\loigiai{ 
	\textbf{\textit{Cách 1.}}\\ 
	Ta có  $\lim\limits_{x\to+\infty}\dfrac{2x^4-5x^3}{3x^4+3x+2}=\lim\limits_{x\to+\infty}\dfrac{x^4\left(2-\dfrac{5}{x}\right)}{x^4\left(3+\dfrac{3}{x^3}+\dfrac{2}{x^4}\right)}=\lim\limits_{x\to+\infty}\dfrac{\left(2-\dfrac{5}{x}\right)}{\left(3+\dfrac{3}{x^3}+\dfrac{2}{x^4}\right)}=\dfrac{2-0}{3+0+0}=\dfrac{2}{3}$.\\
	Suy  ra $a=2,b=3$.\\
	 \textbf{\textit{Cách 2.}}\\
	Nhập biểu  thức $\dfrac{2x^4-5x^3}{3x^4+3x+2}$ vào máy tính bỏ túi.\\
	Vì $x\to+\infty$ nên bấm CALC với $x=10^{20}$, ta thu được  kết quả $\dfrac{2}{3}$, suy ra $a=2$ và $b=3$.\\
	} 
\end{ex}

% nguồn: Ok_Toan11_TracNghiem_Lan2
\begin{ex}%[1D3V2-3]%[ma:48f7ac5f]
	Cho $\lim\limits_{x\to+\infty}\left(\dfrac{x^2+3x+1}{x+1}+a  x+b\right)=1$. Khi đó giá trị của biểu thức $T=a+b$ bằng
	\choice 
		 {\True $-2$}
		{$0$} 
		{$1$} 
		 {$2$}
	 \loigiai{
	\[\begin{aligned}
	\lim\limits_{x\to+\infty}\left(\dfrac{x^2+3 x+1}{x+1}+a  x+b\right)=1 &\Leftrightarrow\lim\limits_{x\to+\infty}\left(\dfrac{(a+1) x^2+(a+b+3) x+b+1}{x+1}\right)=1\\
	 &\Leftrightarrow\lim\limits_{x\to+\infty}\left(\dfrac{(a+1) x+(a+b+3)+\dfrac{b+1}{x}}{1+\dfrac{1}{x}}\right)=1\\
	 &\Leftrightarrow\heva{&{a+1=0}\\&{a+b+3=1}}\\&\Leftrightarrow\heva{&a=-1\\&b=-1.}
\end{aligned}\]
 Vậy $T=a+b=-2$.
} 
\end{ex}

% nguồn: Ok_Toan11_TracNghiem_Lan2
\begin{ex}%[1D3V2-3]%[ma:0dfaca1c]
	Để $\lim\limits_{x\to -\infty}\dfrac{\sqrt{4x^2+x+1}+4}{mx-2}=\dfrac{1}{2}$. Giá trị  của $m$ thuộc tập hợp nào sau đây?
	 \choice
		{$[3;6]$} 
		{$[-3;0]$} 
		{\True  $[-6;-3]$}
		{$[1;3]$} 
	 \loigiai{
		Ta có 
		\begin{eqnarray*}
			 &&\lim\limits_{x\to -\infty}\dfrac{\sqrt{4x^2+x+1}+4}{mx-2}\\
			 &=&\lim\limits_{x\to -\infty}\dfrac{-\sqrt{4+\dfrac{1}{x}+\dfrac{1}{x^2}}+\dfrac{4}{x}}{m-\dfrac{2}{x}}\\
			 &=&-\dfrac{2}{m}.
		\end{eqnarray*}
	Theo bài  ra ta có $-\dfrac{2}{m}=\dfrac{1}{2}\Leftrightarrow m=-4\in [-6;-3]$.}
\end{ex}

% nguồn: Ok_Toan11_TracNghiem_Lan2
\begin{ex}%[1D3V2-3]%[ma:ee5459c9]
	Cho  số thực $a$ thỏa mãn $\lim\limits_{x\to +\infty}\dfrac{a\sqrt{2x^2+3}+2\,017}{2x+2\,018}=\dfrac{1}{2}$. Khi đó giá trị của $a$ là
	 \choice
		{\True $a=\dfrac{\sqrt{2}}{2}$} 
		 {$a=\dfrac{-\sqrt{2}}{2}$}
		 {$a=\dfrac{1}{2}$}
		{$a=-\dfrac{1}{2}$} 
	\loigiai{ 
	Ta có $\lim\limits_{x\to +\infty}\dfrac{a\sqrt{2x^2+3}+2\,017}{2x+2\,018}=\dfrac{1}{2}$$\Leftrightarrow \lim\limits_{x\to  +\infty}\dfrac{a\sqrt{2+\dfrac{3}{x^2}}+\dfrac{2\,017}{x}}{2+\dfrac{2\,018}{x}}=\dfrac{1}{2}$$\Leftrightarrow \dfrac{a\sqrt{2}}{2}=\dfrac{1}{2}$$\Leftrightarrow a=\dfrac{\sqrt{2}}{2}$.}
\end{ex}

% nguồn: Ok_Toan11_TracNghiem_Lan2
\begin{ex}%[1D3V2-3]%[ma:fb6096da]
	Biết $\lim\limits_{x\to 0} \dfrac{\sqrt{5}-\sqrt{5-x^2}}{\sqrt{x^2+16}-4}=\dfrac{a}{\sqrt{b}}$, trong đó $a$ là số nguyên,  $b$ là số nguyên tố. Ta có tổng $a+2b$ bằng
	 \choice
		 {$13$}
		{$3$} 
		{\True $14$} 
		{$8$} 
	 \loigiai{
		Ta  có
		\begin{eqnarray*}
			& &\dfrac{\sqrt{5}-\sqrt{5-x^2}}{\sqrt{x^2+16}-4}=\dfrac{\left(\sqrt{5}-\sqrt{5-x^2} \right)\left(\sqrt{x^2+16}+4\right)}{x^2}\\ &=&\dfrac{\left(\sqrt{5}-\sqrt{5-x^2} \right)\left(\sqrt{5}+\sqrt{5-x^2} \right)}{x^2\left(\sqrt{5}+\sqrt{5-x^2} \right)}=\dfrac{x^2\left(\sqrt{x^2+16}+4\right)}{x^2\left(\sqrt{5}+\sqrt{5-x^2} \right)}=\dfrac{\left(\sqrt{x^2+16}+4\right)}{\left(\sqrt{5}+\sqrt{5-x^2}  \right)}.
		\end{eqnarray*}
		Khi đó ta có 
		 $\lim\limits_{x\to 0} \dfrac{\sqrt{5}-\sqrt{5-x^2}}{\sqrt{x^2+16}-4}=\lim\limits_{x\to 0} \dfrac{\left(\sqrt{x^2+16}+4\right)}{\left(\sqrt{5}+\sqrt{5-x^2} \right)}=\dfrac{4}{\sqrt{5}} \Rightarrow a+2b=14$.
	 }
\end{ex}

% nguồn: Ok_Toan11_TracNghiem_Lan2
\begin{ex}%[1D3V2-3]%[ma:32c56c5e]
	Biết $\lim\limits_{x\to 1} \dfrac{x+2-\sqrt{x+8}}{\sqrt{x+3}+x-3}=\dfrac{a}{b}$ trong đó $\dfrac{a}{b}$ là phân số tối giản, $a$, $b$ là các số  nguyên dương. Giá trị của $3a+2b$ bằng
	\choice 
		{\True $12$} 
		{$13$} 
		{$10$} 
		{$5$} 
	\loigiai{ 
		Ta  có
		\begin{align*}
			 \lim\limits_{x\to 1} \dfrac{x+2-\sqrt{x+8}}{\sqrt{x+3}+x-3}&=\lim\limits_{x\to 1} \dfrac{\left[\left(x+2\right)^2-\left(x+8\right)\right]\left[\sqrt{x+3}-\left(x-3\right)\right]}{\left[\left(x+3\right)-\left(x-3\right)^2\right]\left[x+2+\sqrt{x+8} \right]}\\
			&=\lim\limits_{x\to 1}  \dfrac{\left(x^2+3x-4\right)\left[\sqrt{x+3}-\left(x-3\right)\right]}{\left(-x^2+7x-6\right)\left[x+2+\sqrt{x+8} \right]}\\
			&=\lim\limits_{x\to 1}  \dfrac{\left(x-1\right)\left(x+4\right)\left[\sqrt{x+3}-\left(x-3\right)\right]}{\left(x-1\right)\left(6-x\right)\left[x+2+\sqrt{x+8} \right]}\\
			 &=\lim\limits_{x\to 1} \dfrac{\left(x+4\right)\left[\sqrt{x+3}-\left(x-3\right)\right]}{\left(6-x\right)\left[x+2+\sqrt{x+8} \right]}\\
			 &=\dfrac{5\cdot 4}{5\cdot 6}=\dfrac{2}{3}.
		\end{align*}
		Do đó $a=2$, $b=3$. Vậy $3a+2b=12$. 
	} 
\end{ex}

% nguồn: Ok_Toan11_TracNghiem_Lan2
\begin{ex}%[1D3V2-3]%[ma:e5227706]
	Cho $a$, $b$, $c$ là các số thực khác  $0$. Để $\lim\limits_{x\to -\infty}\dfrac{\sqrt{x^2-3x}+ax}{bx-1}=3$ thì
	 \choice
		{\True  $\dfrac{a-1}{b}=3$}
		 {$\dfrac{a+1}{b}=3$}
		{$\dfrac{-a-1}{b}=3$} 
		 {$M$}
	 \loigiai{
		Ta có 
		\begin{eqnarray*}
			&&\lim\limits_{x\to -\infty}\dfrac{\sqrt{x^2-3x}+ax}{bx-1}\\ 
			&=&\lim\limits_{x\to -\infty}\dfrac{x^2-3x-{{\left( ax \right)}^{2}}}{\left(  bx-1 \right)\left( \sqrt{x^2-3x}-ax \right)}\\
			&=&\lim\limits_{x\to -\infty}\dfrac{x\left[\left( 1-a^2 \right)x-3 \right]}{\left( bx-1 \right)\left( \sqrt{x^2-3x}-ax \right)}\\ 
			 &=&\lim\limits_{x\to -\infty}\dfrac{\left( 1-a^2 \right)-\dfrac{3}{x}}{\left( b-\dfrac{1}{x} \right)\left( -\sqrt{1-\dfrac{3}{x}}-a \right)}\\
			&=&\dfrac{\left(  1-a^2 \right)}{b\left( -1-a \right)}=\dfrac{a-1}{b}=3.
	\end{eqnarray*}}
\end{ex}

% nguồn: Ok_Toan11_TracNghiem_Lan2
\begin{ex}%[1D3V2-3]%[ma:cf5c6588]
	Biết $\lim\limits_{x\to 1} \dfrac{\sqrt{x^2+x+2}-\sqrt[3]{7x+1}}{\sqrt{2} \left(x-1\right)}=\dfrac{a\sqrt{2}}{b}+c$ với $a$, $b$ nguyên dương, $c$ $\in  \mathbb{Z}$ và $\dfrac{a}{b}$ là phân số tối giản. Giá trị của $a+b+c$ bằng
	 \choice
		 {$5$}
		 {$37$}
		 {\True $13$}
		 {$51$}
	 \loigiai{
		Ta  có
		\begin{align*}
			\lim\limits_{x\to 1} \dfrac{\sqrt{x^2+x+2}-\sqrt[3]{7x+1}}{\sqrt{2} \left(x-1\right)}&=\lim\limits_{x\to 1} \dfrac{\sqrt{x^2+x+2}-2+2-\sqrt[3]{7x+1}}{\sqrt{2}  \left(x-1\right)}\\
			&=\lim\limits_{x\to  1} \dfrac{\sqrt{x^2+x+2}-2}{\sqrt{2} \left(x-1\right)}+\lim\limits_{x\to 1} \dfrac{2-\sqrt[3]{7x+1}}{\sqrt{2} \left(x-1\right)}\\
			 &=I+J.
		\end{align*}
		\begin{itemize}
			\item Tính $I$ 
			\begin{align*}
				I&=\lim\limits_{x\to 1} \dfrac{\sqrt{x^2+x+2}-2}{\sqrt{2} \left(x-1\right)}\\ 
				 &=\lim\limits_{x\to 1} \dfrac{x^2+x+2-4}{\sqrt{2} \left(x-1\right)\left(\sqrt{x^2+x+2}+2\right)}\\
				&=\lim\limits_{x\to 1}  \dfrac{\left(x-1\right)\left(x+2\right)}{\sqrt{2} \left(x-1\right)\left(\sqrt{x^2+x+2}+2\right)}\\
				&=\lim\limits_{x\to  1} \dfrac{x+2}{\sqrt{2} \left(\sqrt{x^2+x+2}+2\right)}=\dfrac{3}{4\sqrt{2}}.
			\end{align*}
			\item Tính  $J$
			\begin{align*}
				J&=\lim\limits_{x\to 1} \dfrac{2-\sqrt[3]{7x+1}}{\sqrt{2}  \left(x-1\right)}\\
				&=\lim\limits_{x\to 1} \dfrac{8-7x-1}{\sqrt{2} \left(x-1\right)\left[4+2\sqrt[3]{7x+1}+\left(\sqrt[3]{7x+1}  \right)^2\right]}\\
				 &=\lim\limits_{x\to 1} \dfrac{-7}{\sqrt{2} \left[4+2\sqrt[3]{7x+1}+\left(\sqrt[3]{7x+1} \right)^2\right]}=-\dfrac{7}{12\sqrt{2}}.
			\end{align*}
		\end{itemize}
		Do  đó $\lim\limits_{x\to 1} \dfrac{\sqrt{x^2+x+2}-\sqrt[3]{7x+1}}{\sqrt{2} \left(x-1\right)}=I+J=\dfrac{\sqrt{2}}{12}$.\\
		Suy ra $a=1$,  $b=12$, $c=0$. Vậy $a+b+c=13$.
	} 
\end{ex}

% nguồn: Ok_Toan11_TracNghiem_Lan2
\begin{ex}%[1D3V2-3]%[ma:cdf26151]
	Tính giới hạn $\lim\limits_{x\to-\infty}\dfrac{\sqrt{x^2-2x}-\sqrt{9x^2+1}}{2x+3}$. 
	 \choice
		 {$-\dfrac{3}{2}$}
		{$2$} 
		 {$-1$}
		{\True  $1$}
	\loigiai{ 
	Ta có $\begin{aligned}
	 \lim\limits_{x\to-\infty}\dfrac{\sqrt{x^2-2x}-\sqrt{9x^2+1}}{2x+3}&=\lim\limits_{x\to-\infty}\dfrac{\sqrt{x^2\left(1-\dfrac{2}{x}\right)}-\sqrt{x^2\left(9+\dfrac{1}{x^2}\right)}}{2x+3}\\&=\lim\limits_{x\to-\infty}\dfrac{-x\sqrt{1-\dfrac{2}{x}}+x\sqrt{9+\dfrac{1}{x^2}}}{2x+3}
	 \\&=\lim\limits_{x\to-\infty}\dfrac{x\left(-\sqrt{1-\dfrac{2}{x}}+\sqrt{9+\dfrac{1}{x^2}}\right)}{x\left(2+\dfrac{3}{x}\right)}\\&=\lim\limits_{x\to-\infty}\dfrac{-\sqrt{1-\dfrac{2}{x}}+\sqrt{9+\dfrac{1}{x^2}}}{2+\dfrac{3}{x}}=1.
\end{aligned}$
 }
\end{ex}

% nguồn: Ok_Toan11_TracNghiem_Lan2
\begin{ex}%[1D3V2-3]%[ma:64ac70fc]
	Có bao nhiêu giá trị $a > 0$ sao cho  $\lim\limits_{x\to a}\dfrac{x^3-(1+a^2)x+a}{x^3-a^3}=\dfrac{1}{3}$?
	 \choice
		 {$0$}
		 {\True $2$}
		{$4$} 
		 {$1$}
	\loigiai{ 
		 Ta có
		\begin{align*}
			\lim\limits_{x\to  a}\dfrac{x^3-(1+a^2)x+a}{x^3-a^3}&=\lim\limits_{x\to a}\dfrac{x(x^2-a^2)-(x-a)}{x^3-a^3}\\
			&=\lim\limits_{x\to a}\dfrac{x(x+a)(x-a)-(x-a)}{(x-a)(x^2+ax+a^2)}\\ 
			&=\dfrac{2a^2-1}{3a^2} 
		\end{align*}
		\begin{align*}
			\lim\limits_{x\to a}\dfrac{x^3-(1+a^2)x+a}{x^3-a^3}=\dfrac{1}{3}\Leftrightarrow \dfrac{2a^2-1}{3a^2}=\dfrac{1}{3}\Rightarrow a=\pm  1.
		\end{align*}
		Thử lại thấy thỏa mãn. Vậy  có $2$ giá trị của $a$.
	 }
\end{ex}

% nguồn: Ok_Toan11_TracNghiem_Lan2
\begin{ex}%[1D3V2-3]%[ma:56850a85]
	Cho $a$, $b$ là số nguyên và  $\lim\limits_{x\to 1} \dfrac{ax^2+bx-5}{x-1}=7$. Tính $a^2+b^2+a+b$.
	 \choice
		{\True$18$} 
		 {$1$}
		 {$15$}
		{$5$} 
	 \loigiai{
		Vì  $\lim\limits_{x\to 1} \dfrac{ax^2+bx-5}{x-1}=7$ hữu hạn nên $x=1$ phải là nghiệm của phương trình $ax^2+bx-5=0$.\\
		 Suy ra $a+b-5=0\Rightarrow b=5-a$. \hfill (1)\\
		Khi đó $\lim\limits_{x\to 1}  \dfrac{ax^2+\left(5-a\right)x-5}{x-1}=\lim\limits_{x\to 1} \dfrac{\left(x-1\right)\left(ax+5\right)}{x-1}=a+5=7$. \hfill (2)\\
		Từ (1) và (2) $\Rightarrow a=2$  nên $b=3$.\\
		Suy  ra $a^2+b^2+a+b=18$.
	 }
\end{ex}

% nguồn: Ok_Toan11_TracNghiem_Lan2
\begin{ex}%[1D3V2-3]%[ma:7ea196f5]
	Cho $\lim\limits_{x\to-\infty}\dfrac{2x+\sqrt{3x^2-2}}{\sqrt{4x^2+1}-|x|}=\dfrac{\sqrt{3}+b}{c}$ ($a$, $b\in\mathbb{Z}$). Giá trị của $A=b  c$?
	 \choice
		{$A=6$} 
		 {$A=-6$}
		 {$A=2$}
		{\True $A=-2$} 
	\loigiai{ 
	Ta có  $\lim\limits_{x\to-\infty}\dfrac{2x+\sqrt{3x^2-2}}{\sqrt{4x^2+1}-|x|}=\lim\limits_{x\to-\infty}\dfrac{2x-x\sqrt{3-\dfrac{2}{x^2}}}{-x\sqrt{4+\dfrac{1}{x^2}}+x}=\lim\limits_{x\to-\infty}\dfrac{2-\sqrt{3-\dfrac{2}{x^2}}}{-\sqrt{4+\dfrac{1}{x^2}}+1}=\dfrac{2-\sqrt{3}}{-1}=\dfrac{\sqrt{3}+(-2)}{1}$.\\
	Suy ra  $b=-2$, $c=1$, $A=bc=-2$.
	 }
\end{ex}

% nguồn: Ok_Toan11_TracNghiem_Lan2
\begin{ex}%[1D3V2-3]%[ma:c79be608]
	 Giá trị của giới hạn $\lim\limits_{x\to 0} \dfrac{\sqrt{4+x}-\sqrt[3]{8-x}}{x}$ là
	 \choice
		{\True  $\dfrac{1}{3}$}
		{$-\dfrac{13}{12}$} 
		 {$\dfrac{11}{12}$}
		 {$\dfrac{5}{6}$}
	 \loigiai{
		\begin{align*}
			\lim\limits_{x\to 0} \dfrac{\sqrt{4+x}-\sqrt[3]{8-x}}{x}&=\lim\limits_{x\to 0} \left(\dfrac{\sqrt{4+x}-2}{x}+\dfrac{2-\sqrt[3]{8-x}}{x} \right)\\ 
			&=\lim\limits_{x\to  0} \left(\dfrac{x}{x\left(\sqrt{x+4}+2\right)}+\dfrac{x}{x\left(4+2\sqrt[3]{8-x}+\left(\sqrt[3]{8-x} \right)^2\right)} \right)\\
			&=\lim\limits_{x\to 0} \left(\dfrac{1}{\sqrt{x+4}+2}+\dfrac{1}{4+2\sqrt[3]{8-x}+\left(\sqrt[3]{8-x}  \right)^2} \right)\\
			&=\dfrac{1}{4}+\dfrac{1}{12}=\dfrac{1}{3}. 
		\end{align*}
	 }
\end{ex}

% nguồn: Ok_Toan11_TracNghiem_Lan2
\begin{ex}%[1D3V2-3]%[ma:650e4343]
	Cho $\lim\limits_{x\to 1} \dfrac{x^2+ax+b}{x^2-1}=\dfrac{-1}{2}$, ($a$, $b\in \mathbb{R}$).  Tổng $S=a^2+b^2$ bằng
	\choice 
		{\True  $S=13$}
		{$S=9$} 
		{$S=4$} 
		 {$S=1$}
	\loigiai{ 
		Vì hàm số có giới hạn hữu hạn tại $x=1$ nên biểu  thức tử nhận $x=1$ làm nghiệm hay $1+a+b=0$.\\
		Áp dụng vào giả thiết, ta  được
		\begin{eqnarray*}
			&&\lim\limits_{x\to  1} \dfrac{x^2+ax-1-a}{x^2-1}=\dfrac{-1}{2}\\
			 &\Leftrightarrow&\lim\limits_{x\to 1} \dfrac{\left(x-1\right)\left(x+1+a\right)}{\left(x-1\right)\left(x+1\right)}=-\dfrac{1}{2}\\
			 &\Leftrightarrow&\lim\limits_{x\to 1} \dfrac{x+1+a}{x+1}=-\dfrac{1}{2}\\
			&\Leftrightarrow&\dfrac{2+a}{2}=-\dfrac{1}{2}\\ 
			 &\Leftrightarrow& a=-3.
		\end{eqnarray*}
		Suy  ra $b=2$.
		Vậy  $a^2+b^2=13$.
	} 
\end{ex}

% nguồn: Ok_Toan11_TracNghiem_Lan2
\begin{ex}%[1D3V2-3]%[ma:ce820cb4]
	Biết $\lim\limits_{x\to  1} \dfrac{\sqrt{3x^2+2}-\sqrt{4+x}}{x^2-1}=\dfrac{\sqrt{a}}{b}$ với $a,b$ là các số nguyên dương và nguyên tố cùng nhau. Tính $P=a-b$.
	 \choice
		 {$P=5$}
		{\True $P=1$} 
		 {$P=2$}
		 {$P=3$}
	 \loigiai{
		\begin{align*}
			\lim\limits_{x\to 1} \dfrac{\sqrt{3x^2+2}-\sqrt{4+x}}{x^2-1}&=\lim\limits_{x\to  1} \dfrac{\sqrt{3x^2+2}-\sqrt{4+x}}{x^2-1}\\
			 &=\lim\limits_{x\to 1} \dfrac{\sqrt{3x^2+2}-\sqrt{5}+\sqrt{5}-\sqrt{4+x}}{x^2-1}\\
			&=\lim\limits_{x\to 1} \left(\dfrac{\sqrt{3x^2+2}-\sqrt{5}}{x^2-1}+\dfrac{\sqrt{5}-\sqrt{4+x}}{x^2-1} \right)\\ 
			 &=\lim\limits_{x\to 1} \left(\dfrac{3\left(x^2-1\right)}{\left(x^2-1\right)\left(\sqrt{3x^2+2}+\sqrt{5} \right)}+\dfrac{1-x}{\left(x^2-1\right)\left(\sqrt{5}+\sqrt{4+x} \right)} \right)\\
			&=\lim\limits_{x\to 1} \left(\dfrac{3}{\sqrt{3x^2+2}+\sqrt{5}}-\dfrac{1}{\left(x+1\right)\left(\sqrt{5}+\sqrt{4+x} \right)}  \right)\\
			 &=\dfrac{\sqrt{5}}{4}.
		\end{align*}
		Vậy $a=5$,  $b=4$ nên $P=1$.
	 }
\end{ex}

% nguồn: Ok_Toan11_TracNghiem_Lan2
\begin{ex}%[1D3V2-3]%[ma:0a4a2415]
	Biết $\lim\limits_{x\to  1} \dfrac{\sqrt{x+3}-\sqrt[3]{9-x}}{x-1}=\dfrac{a}{b}$ trong đó $\dfrac{a}{b}$ là phân số tối giản, $a$, $b$ là các số nguyên dương. Giá trị của $b-a$ bằng
	\choice 
		{\True $2$} 
		 {$\dfrac{1}{3}$}
		 {$1$}
		{$\dfrac{1}{12}$} 
	 \loigiai{
		 Ta có
		\begin{align*}
			\lim\limits_{x\to  1} \dfrac{\sqrt{x+3}-\sqrt[3]{9-x}}{x-1}&=\lim\limits_{x\to 1} \left(\dfrac{\sqrt{x+3}-2}{x-1}+\dfrac{2-\sqrt[3]{9-x}}{x-1} \right)\\
			 &=\lim\limits_{x\to 1} \left(\dfrac{x-1}{(x-1)\left(\sqrt{x+3}+2\right)}+\dfrac{x-1}{(x-1)\left(4+2\sqrt[3]{9-x}+\left(\sqrt[3]{9-x} \right)^2\right)} \right)\\
			&=\lim\limits_{x\to 1} \left(\dfrac{1}{\sqrt{x+3}+2}+\dfrac{1}{4+2\sqrt[3]{9-x}+\left(\sqrt[3]{9-x} \right)^2}  \right)\\
			 &=\dfrac{1}{4}+\dfrac{1}{12}=\dfrac{1}{3}.
		\end{align*}
		Do đó $ a=1$, $b=3$.  Vậy $b-a=2$.
	} 
\end{ex}

% nguồn: Ok_Toan11_TracNghiem_Lan2
\begin{ex}%[1D3V2-3]%[ma:488a4e21]
	Tính $\displaystyle\lim\limits _{x  \to-\infty}\left(\sqrt{2 x^{2}+x}+x\right)$.
	 \choice
		{\True $+\infty$} 
		{$-1$} 
		{$-\infty$} 
		 {$0$}
	 \loigiai{
		Ta  có
		\begin{eqnarray*}
			&\displaystyle\lim\limits _{x \to-\infty}\left(\sqrt{2x^{2}+x}+x\right)&=\displaystyle\lim\limits _{x \to-\infty}\left(\sqrt{x^{2}\left(2+\dfrac{1}{x}\right)}+x\right)\\ 
			&&=\displaystyle\lim\limits _{x \to-\infty}\left(-x  \sqrt{2+\dfrac{1}{x}}+x\right)=\displaystyle\lim\limits _{x \to-\infty}\left[x\left(-\sqrt{2+\dfrac{1}{x}}+1\right)\right].
		\end{eqnarray*}
	Vì $\displaystyle\lim\limits  _{x \to-\infty}x=-\infty$ và $\displaystyle\lim\limits _{x \to-\infty}\left(-\sqrt{2+\dfrac{1}{x}}+1\right)=1-\sqrt{2}<0$ nên $\displaystyle\lim\limits _{x \to-\infty}\left(\sqrt{2 x^{2}+x}+x\right)=+\infty$.}
\end{ex}

% nguồn: Ok_Toan11_TracNghiem_Lan2
\begin{ex}%[1D3V2-3]%[ma:ed7b5b80]
	Cho $\lim\limits_{x\to 1} \left(\dfrac{\sqrt{x^2+x+2}-\sqrt[3]{3x+5}}{x^2-3x+2} \right)=\dfrac{a}{b}$ trong đó $\dfrac{a}{b}$ là phân số tối giản, $a$, $b$ là  các số nguyên. Tính tổng $P=a^2+b^2$.
	 \choice
		{\True  $P=5$}
		 {$P=3$}
		 {$P=2$}
		 {$P=-2$}
	 \loigiai{
		 Ta có
		\begin{align*}
			&\lim\limits_{x\to 1} \left(\dfrac{\sqrt{x^2+x+2}-\sqrt[3]{3x+5}}{x^2-3x+2} \right)\\ 
			 &= \lim\limits_{x\to 1} \left(\dfrac{\sqrt{x^2+x+2}-2}{x^2-3x+2}+\dfrac{2-\sqrt[3]{3x+5}}{x^2-3x+2} \right)\\
			&=\lim\limits_{x\to 1} \left(\dfrac{x^2+x-2}{\left(x^2-3x+2\right)\left(\sqrt{x^2+x+2}+2\right)}+\dfrac{3-3x}{\left(x^2-3x+2\right)\left[4+2\sqrt[3]{3x+5}+\left(\sqrt[3]{3x+5}  \right)^2\right]} \right)\\
			&=\lim\limits_{x\to 1}  \left(\dfrac{\left(x-1\right)\left(x+2\right)}{\left(x-1\right)\left(x-2\right)\left(\sqrt{x^2+x+2}+2\right)}+\dfrac{-3\left(x-1\right)}{\left(x-1\right)\left(x-2\right)\left[4+2\sqrt[3]{3x+5}+\left(\sqrt[3]{3x+5} \right)^2\right]} \right)\\
			&=\lim\limits_{x\to  1} \left(\dfrac{x+2}{\left(x-2\right)\left(\sqrt{x^2+x+2}+2\right)}+\dfrac{-3}{\left(x-2\right)\left[4+2\sqrt[3]{3x+5}+\left(\sqrt[3]{3x+5} \right)^2\right]} \right)\\
			&=\dfrac{-3}{4}+\dfrac{3}{12}=\dfrac{-1}{2}. 
		\end{align*}
		Theo giả thiết  ta có $-\dfrac{1}{2}=\dfrac{a}{b}$.\\
		Vì $\dfrac{a}{b}$ là phân  số tối giản, $a$, $b$ là số nguyên nên
		 $\hoac{&\heva{&a=-1\\&b=2}\\&\heva{&a=1\\&b=-2.}}$\\
		Vậy $P=a^2+b^2=5$. 
	} 
\end{ex}

% nguồn: Ok_Toan11_TracNghiem_Lan2
\begin{ex}%[1D3C2-3]%[ma:3b836d29]
	Tính $\lim\limits_{x\to 4} \dfrac{2^{2018}-x^{1009}}{4-x}$. 
	 \choice
		{$+\infty$} 
		{\True  $1009\cdot 2^{2016}$}
		{$1009\cdot 2^{2018}$} 
		 {$1009\cdot 4^{2018}$}
	 \loigiai{
		Ta  có
		\begin{align*}
			\lim\limits_{x\to  4} \dfrac{2^{2018}-x^{1009}}{4-x}&=\lim\limits_{x\to 4} \dfrac{4^{1009}-x^{1009}}{4-x}\\
			&=\lim\limits_{x\to 4} \dfrac{\left(4-x\right)\left(4^{1008}+4^{1007}\cdot x+4^{1006}\cdot  x^2+\cdots+4\cdot x^{1007}+x^{1008} \right)}{4-x}\\
			&=\lim\limits_{x\to 4} \left(4^{1008}+4^{1007}\cdot  x+4^{1006}\cdot x^2+\cdots+4\cdot x^{1007}+x^{1008} \right)\\
			&=1009\cdot 4^{1008}=1009\cdot 2^{2016}. 
		\end{align*}
	 }
\end{ex}

% nguồn: Ok_Toan11_TracNghiem_Lan2
\begin{ex}%[1D3C2-3]%[ma:cef7f4ba]
	Tính giới hạn $\lim\limits_{x\to 1} \dfrac{x^m-x^n}{x-1}$, $\left(m;n\in  \mathbb{N}^{*} \right)$.
	\choice 
		 {$1$}
		{$+\infty$} 
		 {$m$}
		{\True $m-n$} 
	 \loigiai{
		\begin{align*}
			\lim\limits_{x\to 1} \dfrac{x^m-x^n}{x-1}&=\lim\limits_{x\to  1} \left(\dfrac{x^m-1}{x-1}-\dfrac{x^n-1}{x-1} \right)\\
			&=\lim\limits_{x\to  1} \left[\left(x^{m-1}+x^{m-2}+\cdots+1\right)-\left(x^{n-1}+x^{n-2}+\cdots+1\right)\right]\\
			 &=m-n.
		\end{align*}
	} 
\end{ex}

% nguồn: Ok_Toan11_TracNghiem_Lan2
\begin{ex}%[1D3C2-3]%[ma:2c3bdf52]
	Tính giới hạn của hàm số  $\lim\limits_{x\to 1} \dfrac{x^n-nx+n-1}{\left(x-1\right)^2}$.
	 \choice
		{\True  $\dfrac{n^2-n}{2}$}
		 {$\dfrac{n^2+n}{2}$}
		 {$\dfrac{n}{2}$}
		{$\dfrac{n^2}{2}$} 
	 \loigiai{
		Ta  có
		\begin{align*}
			\lim\limits_{x\to 2} \dfrac{x^n+nx+n-1}{\left(x-1\right)^2}&=\lim\limits_{x\to 2}  \dfrac{\left(x^n-1\right)-n\left(x-1\right)}{\left(x-1\right)^2}\\
			&=\lim\limits_{x\to 1} \dfrac{\left(x^{n-1}+x^{n-2}+\cdots+x+1\right)-n}{x-1}\\ 
			&=\lim\limits_{x\to  1} \dfrac{\left(x^{n-1}-1\right)+\left(x^{n-2}-1\right)+\cdots+\left(x-1\right)}{x-1}\\
			&=\lim\limits_{x\to 1}  \left(\left(x^{n-1}+x^{n-2}+\cdots+1\right)+\left(x^{n-3}+x^{n-4}+\cdots+1\right)+\cdots+1\right)\\
			 &=\left(n-1\right)+\left(n-2\right)+\cdots+1\\
			&=\dfrac{n\left(n-1\right)}{2}=\dfrac{n^2-n}{2}. 
		\end{align*}
	 }
\end{ex}

% nguồn: Ok_TracNghiemLop11_Lan1
\begin{ex}%[1D3B2-3]%[ma:e616416d]
	 Cho dãy số $( u_n)$ với $u_n=2n-1$. Dãy số $( u_n)$ là dãy số
	 \choice
		{ Bị chặn  trên bởi 1}
		{  Giảm}
		{ Bị chặn dưới  bởi 2}
		{ \True Tăng} 
	 \loigiai{
	 $\forall n\in {{\mathbb{N}}^{*}}$ ta có ${{u}_{n+1}}-u_n=2( n+1)-1-( 2n-1)=2>0$ nên ${{u}_{n+1}}>u_n$.\\ Vậy dãy số $( u_n)$ tăng.}
\end{ex}

% nguồn: Ok_TracNghiemLop11_Lan1
\begin{ex}%[1D3B2-3]%[ma:fdb8ca11]
	Cho dãy số $( u_n)$ có $u_n=-{{n}^2}+n+1$. Khẳng định nào  sau đây là đúng?
	\choice 
		{ 5  số hạng đầu của dãy là: $-1;1;5;-5;-11;-19$}
		 { ${{u}_{n+1}}=-{{n}^2}+n+2$}
		 { ${{u}_{n-1}}-u_n=1$}
		{ \True Là một  dãy số giảm}
	\loigiai{ 
		Ta có\\ 
	${{u}_{n+1}}-u_n=[-{{(  n+1)}^2}+n+1+1]-[-{{n}^2}+n+1]\\=-{{n}^2}-2n-1+n+2+{{n}^2}-n-1=-2n<0$ $\forall n\ge 1$.\\ Do đó $( u_n)$ là một dãy giảm.}
\end{ex}

% nguồn: Ok_TracNghiemLop11_Lan1
\begin{ex}%[1D3B2-3]%[ma:27e13b8f]
	Trong các dãy số $( u_n)$ cho bởi số  hạng tổng quát $u_n$ sau, dãy số nào là dãy số giảm?
	 \choice
		{  \True $u_n=\dfrac{1}{{{2}^{n}}}$}
		{ $u_n=\dfrac{3n-1}{n+1}$} 
		{ $u_n={{n}^2}$} 
		 { $u_n=\sqrt{n+2}$}
	 \loigiai{
	Ta có $u_n=\dfrac{1}{{{2}^{n}}}$$<\dfrac{1}{{{2}^{n+1}}}={{u}_{n+1}}$  $\forall n\in {{\mathbb{N}}^{*}}$}
\end{ex}

% nguồn: Ok_TracNghiemLop11_Lan1
\begin{ex}%[1D3B2-3]%[ma:a75ef01b]
	Dãy số  nào sau đây là dãy số giảm?
	\choice 
		{ \True $u_n=\dfrac{5-3n}{2n+3}$ $( n\in \mathbb{N}^*)$} 
		{ $u_n=\dfrac{n-5}{4n+1}$ $(  n\in \mathbb{N}^*)$}
		{ $u_n=2{{n}^3}+3$  $( n\in \mathbb{N}^*)$}
		{ $u_n=\cos (  2n+1)$ $( n\in \mathbb{N}^*)$}
	\loigiai{ 
		Xét $u_n=\dfrac{5-3n}{2n+3}$ $( n\in \mathbb{N}^*)$.\\ 
		Ta có ${{u}_{n+1}}-u_n=\dfrac{5-3( n+1)}{2( n+1)+3}-\dfrac{5-3n}{2n+3}$\\ $=\dfrac{2-3n}{2n+5}-\dfrac{5-3n}{2n+3}$ $=\dfrac{( 2-3n)( 2n+3)-( 2n+5)( 5-3n)}{( 2n+5)(  2n+3)}$ \\$=\dfrac{4n-6{n^2}+6-9n-10n+6{n^2}-25+15n}{( 2n+5)( 2n+3)}$$=\dfrac{-19}{( 2n+5)( 2n+3)}<0,\,\forall n\in \mathbb{N}^*$.\\
	Vậy $u_n=\dfrac{5-3n}{2n+3}$ $(n\in \mathbb{N}^*)$ là dãy  giảm.}
\end{ex}

% nguồn: Ok_TracNghiemLop11_Lan1
\begin{ex}%[1D3B2-3]%[ma:8a5e7403]
	Trong các dãy số sau, dãy số nào là dãy số giảm 
	\choice 
		{ $u_n=\dfrac{n-3}{n+1}$} 
		{ $u_n=\dfrac{n}{2}$} 
		{ \True  $u_n=\dfrac{2}{{{n}^2}}$}
		 { $u_n=\dfrac{{{( -1)}^{n}}}{{{3}^{n}}}$}
	 \loigiai{
		Xét $u_n=\dfrac{n-3}{n+1};$ ${{u}_{n+1}}=\dfrac{n-2}{n+2}$. Khi đó:  ${{u}_{n+1}}-u_n=\dfrac{n-2}{n+2}-\dfrac{n-3}{n+1}=\dfrac{4}{( n+1)( n+2)}>0$ $\forall n\in \mathbb{N}^*$
		Vậy $( u_n)$ là dãy  số tăng.\\
		Xét $u_n=\dfrac{n}{2};$ ${{u}_{n+1}}=\dfrac{n+1}{2}$. Khi đó: ${{u}_{n+1}}-u_n=\dfrac{n+1}{2}-\dfrac{n}{2}=\dfrac{1}{2}>0$ $\forall n\in \mathbb{N}^*$. 
		Vậy $( u_n)$ là dãy số tăng.\\ 
		Xét $u_n=\dfrac{2}{{{n}^2}}$, ${{u}_{n+1}}=\dfrac{2}{{{( n+1)}^2}}$. $\dfrac{{{u}_{n+1}}}{u_n}=\dfrac{{{n}^2}}{{{( n+1)}^2}}<\dfrac{{{n}^2}}{{{n}^2}}=1,\forall n\in {{\mathbb{N}}^{*}}$. Vậy $( u_n)$  là dãy giảm.\\
	Xét ${{u}_{1}}=\dfrac{-1}{3};$ ${{u}_{2}}=\dfrac{1}{9};$ ${{u}_{3}}=\dfrac{-1}{27}$. Vậy $( u_n)$ là dãy  số không tăng không giảm}
\end{ex}

% nguồn: Ok_TracNghiemLop11_Lan1
\begin{ex}%[1D3B2-3]%[ma:b4469b4b]
	Trong các dãy số sau, dãy số  nào là dãy số giảm?
	 \choice
		{  \True $u_n=\dfrac{2n+1}{n-1}$}
		 { $u_n={{n}^3}-1$}
		{  $u_n={{n}^2}$}
		 { $u_n=2n$}
	 \loigiai{
		Với  mọi $n\in \mathbb{N}$, $n>1$. Ta có\\
		${{u}_{n+1}}-u_n=\dfrac{2( n+1)+1}{( n+1)-1}-\dfrac{2n+1}{n-1}=\dfrac{2n+3}{n}-\dfrac{2n+1}{n-1}$ 
		$=\dfrac{( 2n+3)( n-1)-n( 2n+1)}{n( n-1)}\\=\dfrac{( 2n+3)( n-1)-n( 2n+1)}{n( n-1)}=\dfrac{-3}{n(  n-1)}<0$, với mọi $n\in \mathbb{N}$, $n>1$.\\
	Suy  ra dãy số giảm.}
\end{ex}

% nguồn: Ok_TracNghiemLop11_Lan1
\begin{ex}%[1D3B2-3]%[ma:f6f7bfc2]
	Cho dãy  số $( u_n)$ với $u_n=\dfrac{{{(-1)}^{n-1}}}{n+1}$. Khẳng định nào sau đây là sai?
	 \choice
		{Số hạng thứ $9$ của dãy số là  $\dfrac{1}{10}$}
		{Số hạng thứ 10 của dãy số là $\dfrac{-1}{11}$} 
		{\True Đây là một dãy số giảm} 
		{Bị chặn trên bởi số $M=1$} 
	\loigiai{ 
	Dãy $u_n$ là một dãy đan dấu  nên khẳng định dãy giảm là sai.}
\end{ex}

% nguồn: Ok_TracNghiemLop11_Lan1
\begin{ex}%[1D3B2-3]%[ma:44ce4150]
	Cho dãy số $( u_n)$ với $u_n=\dfrac{1}{{{n}^2}+n}$. Khẳng định nào sau đây là sai? 
	 \choice
		 { $5$ số hạng đầu của dãy là:$\dfrac{1}{2};\dfrac{1}{6};\dfrac{1}{12};\dfrac{1}{20};\dfrac{1}{30}$}
		{ \True Là  dãy số tăng}
		{ Bị chặn trên bởi số $M=\dfrac{1}{2}$} 
		{ Không bị  chặn}
	\loigiai{ 
		Ta có\\ ${{u}_{n+1}}-u_n=\dfrac{1}{{{( n+1)}^2}+( n+1)}-\dfrac{1}{{{n}^2}+n}\\=\dfrac{1}{( n+1)( n+2)}-\dfrac{1}{n(  n+1)}$
		$=\dfrac{-2}{n( n+1)( n+2)}<0$ với  $n\ge 1$.
	Do đó $(  u_n)$ là dãy giảm.}
\end{ex}


%%%%%%%%%%%%%%%% ĐÚNG/SAI %%%%%%%%%%%%%%%%

% nguồn: Ok_dungsai_lop11_Lan1
\begin{ex}%[1D3N2-3]%[ma:454d0b63]
 Mỗi kết  quả giới hạn sau đúng hay sai?
 \choiceTF 
  {\True $\lim\limits_{x \to 3}  \dfrac{x^2+2x-15}{x-3}=8$}
  {$\lim\limits_{x \to  0} \dfrac{x}{x^2-x}=2$}
  {\True $\lim\limits_{x \to 1} \dfrac{x^2-3x+2}{x^2-5x+4}=\dfrac{1}{3}$} 
  {$\lim\limits_{x \to  -1} \dfrac{2x^2+3x+1}{x^2+x}=\dfrac{1}{2}$}
  \loigiai{
  \begin{itemchoice}
   \itemch  $\lim\limits_{x \to 3} \dfrac{x^2+2x-15}{x-3}=\lim\limits_{x \to 3} \dfrac{(x-3)(x+5)}{x-3}=\lim\limits_{x \to 3} (x+5)=8$.
   \itemch $\lim\limits_{x \to  0} \dfrac{x}{x^2-x}=\lim\limits_{x \to 0} \dfrac{x}{x(x-1)}=\lim\limits_{x \to 0} \dfrac{1}{x-1}=-1$.
   \itemch  $\lim\limits_{x \to 1} \dfrac{x^2-3x+2}{x^2-5x+4}
   =\lim\limits_{x \to 1} \dfrac{(x-1)(x-2)}{(x-1)(x-4)}=\lim\limits_{x \to 1}  \dfrac{x-2}{x-4}=\dfrac{1}{3}$.
   \itemch $\lim\limits_{x \to  -1} \dfrac{2x^2+3x+1}{x^2+x}=\lim\limits_{x \to -1} \dfrac{(x+1)(2x+1)}{(x+1)x}=\lim\limits_{x \to -1} \dfrac{2x+1}{x}=1$.
  \end{itemchoice}
  }
\end{ex}

% nguồn: Ok_dungsai_lop11_Lan1
\begin{ex}%[1D3H2-3]%[ma:ed5a18d1]
 Tìm được các giới hạn sau. 
 \choiceTF 
  {\True  $\lim\limits_{x\to+\infty}\left(x^2+3\right)=+\infty$}
   {$\lim\limits_{x\to-\infty}\left(\sqrt{x^2+x}-x\right)=-\infty $}
  {\True  $\lim\limits_{x\to-\infty}\dfrac{1}{x+2}=0$}
  {$\lim\limits_{x\to+\infty}\sqrt{\dfrac{2x}{x+3}}=2$} 
 \loigiai{ 
  \begin{itemchoice}
    \itemch $\lim\limits_{x\to+\infty}\left(x^2+3\right)= \lim\limits_{x\to+\infty}{x^2}\left(1+\dfrac{3}{x^2}\right)=+\infty$, \\
    do $\lim\limits_{x\to+\infty}{x^2}=+\infty$ và $\lim\limits_{x\to+\infty}\left(1+\dfrac{3}{x^2}\right)=1$.
   \itemch $\lim\limits_{x\to-\infty}\left(\sqrt{x^2+x}-x\right)= \lim\limits_{x\to-\infty}\left(-x\sqrt{1+\dfrac{1}{x}}-x\right)\\ 
   =\lim\limits_{x\to-\infty}x\left(-\sqrt{1+\dfrac{1}{x}}-1\right)=+\infty$, \\ 
   do $\lim\limits_{x\to-\infty}x=-\infty$  và $\lim\limits_{x\to-\infty}\left(-\sqrt{1+\dfrac{1}{x}}-1\right)=-2$.
   \itemch $\lim\limits_{x\to-\infty}\dfrac{1}{x+2}= \lim\limits_{x\to-\infty}\dfrac{x\cdot\dfrac{1}{x}}{x\left(1+\dfrac{2}{x}\right)}=  \lim\limits_{x\to-\infty}\dfrac{\dfrac{1}{x}}{1+\dfrac{2}{x}}=0$.
   \itemch $\lim\limits_{x\to+\infty}\sqrt{\dfrac{2x}{x+3}}= \lim\limits_{x\to+\infty}\sqrt{\dfrac{2x}{x\left(1+\dfrac{3}{x}\right)}}= \lim\limits_{x\to+\infty}\sqrt{\dfrac{2}{1+\dfrac{3}{x}}}=\sqrt 2  $.
  \end{itemchoice}
  }
\end{ex}

% nguồn: Ok_dungsai_lop11_Lan1
\begin{ex}%[1D3H2-3]%[ma:d52eb21a]
 Xét  tính đúng sai trong các mệnh đề sau
  \choiceTF
  {\True $\lim\limits_{x\to  0}(-5x^3-4x+2)=2$}
   {$\lim\limits_{x\to-1}\dfrac{2x-3x^2}{4x+1}=-\dfrac{3}{4}$}
  {$\lim\limits_{x\to-5}\dfrac{x^2+2x-15}{x+5}=+\infty  $}
  {\True $\lim\limits_{x\to-4}\dfrac{x^2+3x-4}{x^2+4x}=\dfrac{5}{4}$} 
 \loigiai{ 
  \begin{itemchoice}
   \itemch $\lim\limits_{x\to 0}(-5x^3-4x+2)=-5\cdot{0^3}-4\cdot 0+2=2$. 
   \itemch $\lim\limits_{x\to-1}\dfrac{2x-3x^2}{4x+1}=\dfrac{2\cdot (-1)-3\cdot{(-1)^2}}{4\cdot (-1)+1}=\dfrac{5}{3}$. 
   \itemch  $\lim\limits_{x\to-5}\dfrac{x^2+2x-15}{x+5}=\lim\limits_{x\to-5}\dfrac{(x+5)(x-3)}{x+5}=\lim\limits_{x\to-5}(x-3)=-5-3=-8$.
    \itemch $\lim\limits_{x\to-4}\dfrac{x^2+3x-4}{x^2+4x}= \lim\limits_{x\to-4}\dfrac{(x-1)(x+4)}{x(x+4)}= \lim\limits_{x\to-4}\dfrac{x-1}{x}=\dfrac{-4-1}{-4}=\dfrac{5}{4}$.
  \end{itemchoice}
 } 
\end{ex}

% nguồn: Ok_dungsai_lop11_Lan1
\begin{ex}%[1D3H2-3]%[ma:dc7ddb5b]
 Cho hàm  số $ f(x)=\dfrac{x^2-1}{x^2-3x+2}$.
 \choiceTF 
  {\True $\lim\limits_{x\to  +\infty}f(x)=1$}
  {$\lim\limits_{x\to  2^-}f(x)=+\infty$}
  {\True Giới hạn $\lim\limits_{x  \to 1}f(x)$ bằng $-2$}
  {$\lim\limits_{x \to 5} \left[f(x)+\dfrac{x^2-12x+35}{25-5x}\right]=\dfrac{m}{n}$ (với $\dfrac{m}{n}$ là phân số tối giản và $m,\, n\in\mathbb{N}$, $n\ne 0$). Khi đó $m^2+n^3=296$} 
  \loigiai{
  \begin{itemchoice}
   \itemch \textbf{Đúng.}\\ Ta có $\lim\limits_{x \to +\infty}\dfrac{x^2-1}{x^2-3x+2}=\lim\limits_{x \to +\infty} \dfrac{x^2\left(1-\dfrac{1}{x^2}\right)}{x^2\left(1-\dfrac{3}{x}+\dfrac{2}{x^2}\right)}=\lim\limits_{x \to +\infty} \dfrac{1-\dfrac{1}{x^2}}{1-\dfrac{3}{x}+\dfrac{2}{x^2}}=1$. 
   \itemch \textbf{Sai.}\\ Ta có \\$\lim\limits_{x \to 2^-}\dfrac{x^2-1}{x^2-3x+2}=\lim\limits_{x \to 2^-} \dfrac{(x-1)(x+1)}{(x-1)(x-2)}=\lim\limits_{x \to 2^-} \dfrac{x+1}{x-2}=-\infty$ vì $\heva{&\lim\limits_{x \to 2^-} (x+1)=3>0\\&\lim\limits_{x \to 2^-}(x-2)=0\\&x-2<0\text{ khi }x\to 2^-.}$ 
   \itemch \textbf{Đúng.}\\ Ta có $\lim\limits_{x \to 1}\dfrac{x^2-1}{x^2-3x+2}=\lim\limits_{x \to 1}  \dfrac{(x-1)(x+1)}{(x-1)(x-2)}=\lim\limits_{x \to 1} \dfrac{x+1}{x-2}=\dfrac{1+1}{1-2}=-2 $.
   \itemch \textbf{Sai.}\\  Ta có
   \[\begin{aligned}
   \lim\limits_{x \to 5} \left[f(x)+\dfrac{x^2-12x+35}{25-5x}\right]&=\lim\limits_{x\to  5}\dfrac{x^2-1}{x^2-3x+2}+\lim\limits_{x \to 5} \dfrac{x^2-12x+35}{25-5x}\\
   &=\dfrac{5^2-1}{5^2-3\cdot5+2}+\lim\limits_{x \to  5} \dfrac{(x-5)(x-7)}{-5(x-5)}\\
   &=2+\lim\limits_{x  \to 5} \dfrac{x-7}{-5}\\
   &=2+\dfrac{5-7}{-5}\\ 
   &=\dfrac{12}{5}. 
  \end{aligned}\]
  Suy ra  $m=12$, $n=5$. Do đó $m^2+n^3=12^2+5^3=269$.
 \end{itemchoice}
 }
\end{ex}

% nguồn: Ok_dungsai_lop11_Lan1
\begin{ex}%[1D3H2-3]%[ma:d37b915c]
 Mỗi kết quả giới hạn  sau đúng hay sai?
  \choiceTF
   {$\lim\limits_{x \to -2} \dfrac{x^3+8}{x+2}=4$}
  {\True  $\lim\limits_{x \to 0} \dfrac{(x+3)^3-27}{x}=27$}
   {$\lim\limits_{x \to 1} \dfrac{x^3-1}{2x^2-2x}=3$}
  {\True $\lim\limits_{x  \to 3} \dfrac{x^2-9}{2x^2+x-21}=\dfrac{6}{13}$}
 \loigiai{ 
  \begin{itemchoice}
   \itemch $\lim\limits_{x \to -2} \dfrac{x^3+8}{x+2}=\lim\limits_{x \to -2} \dfrac{(x+2)(x^2-2x+4)}{x+2}=\lim\limits_{x \to -2}  (x^2-2x+4)=12$.
    \itemch $\lim\limits_{x \to 0} \dfrac{(x+3)^3-27}{x}
    =\lim\limits_{x \to 0} \dfrac{x^3+9x^2+27x}{x}
   =\lim\limits_{x \to  0} \left(x^2+9x+27\right)=27$.
   \itemch $\lim\limits_{x \to 1} \dfrac{x^3-1}{2x^2-2x} 
   =\lim\limits_{x \to  1} \dfrac{(x-1)(x^2+x+1)}{2x(x-1)}
   =\lim\limits_{x \to 1}  \dfrac{x^2+x+1}{2x}=\dfrac{3}{2}$.
   \itemch $\lim\limits_{x \to 3} \dfrac{x^2-9}{2x^2+x-21} 
   =\lim\limits_{x  \to 3} \dfrac{(x-3)(x+3)}{(x-3)(2x+7)}
   =\lim\limits_{x  \to 3} \dfrac{x+3}{2x+7}=\dfrac{6}{13}$.
  \end{itemchoice}
 } 
\end{ex}

% nguồn: Ok_dungsai_lop11_Lan1
\begin{ex}%[1D3H2-3]%[ma:24eef90f]
 Xét tính đúng  sai của các mệnh đề sau
  \choiceTF
  {\True $\lim\limits_{x\to-2}(x^2-x+3)=9$} 
  {$\lim\limits_{x\to 6}\sqrt{\dfrac{1}{x+3}}=3$} 
  {\True $\lim\limits_{x\to  2}\dfrac{x^2-3x+2}{x-2}=1$}
   {$\lim\limits_{x\to-1}\dfrac{2x^2+3x+1}{x^2-1}=\dfrac{1}{3}$}
 \loigiai{ 
  \begin{itemchoice}
   \itemch Ta có  $\lim\limits_{x\to-2}(x^2-x+3)=(-2)^2-(-2)+3=9$.
   \itemch Ta có $\lim\limits_{x\to 6}\sqrt{\dfrac{1}{x+3}}=\sqrt{\dfrac{1}{6+3}}=\dfrac{1}{3}$. 
   \itemch Ta có $\lim\limits_{x\to  2}\dfrac{x^2-3x+2}{x-2}= \lim\limits_{x\to 2}\dfrac{(x-2)(x-1)}{x-2}= \lim\limits_{x\to 2}(x-1)=2-1=1$.
   \itemch Ta  có $\lim\limits_{x\to-1}\dfrac{2x^2+3x+1}{x^2-1}= \lim\limits_{x\to-1}\dfrac{(2x+1)(x+1)}{(x-1)(x+1)}= \lim\limits_{x\to-1}\dfrac{2x+1}{x-1}=\dfrac{-2+1}{-1-1}=\dfrac{1}{2}$.
  \end{itemchoice}
  }
\end{ex}

% nguồn: Ok_dungsai_lop11_Lan2
\begin{ex}%[1D3H2-3]%[ma:9031abc6]
	 Cho giới hạn
	$\lim\limits_{x \rightarrow-\sqrt{3}} f(x)=\lim\limits_{x \rightarrow-\sqrt{3}} \dfrac{2 x^3+6 \sqrt{3}}{3-x^2}$ 
	và 
	$\lim\limits_{x \rightarrow x_0}  g(x)=\lim\limits_{x \rightarrow x_0} \dfrac{x^2+a x+b}{x-1} $.
	\choiceTF 
		 {$\lim\limits_{x \to -\sqrt{3}}f(x)=2$}
		{\True Nếu $a=1$  và $b=0$ thì $\lim\limits_{x \to 2}g(x)=5$}
		{\True Nếu $\lim\limits_{x \to -\sqrt{3}}f(x)=a \sqrt{3}+b$ thì $a+b=3$} 
		{\True Nếu $\lim\limits_{x \to 1}g(x)=\lim\limits_{x \to 1}\dfrac{x^2+a  x+b}{x-1}=1$ thì $a+b=-1$}
	 \loigiai{
		\begin{itemchoice}
			\itemch  \textbf{Sai.} Ta có
			\[ 
			\lim\limits_{x \to -\sqrt{3}}f(x) 
			=\lim\limits_{x \to -\sqrt{3}}\dfrac{2 x^3+6 \sqrt{3}}{3-x^2} 
			=\lim\limits_{x \to -\sqrt{3}}\dfrac{2 x^2-2 \sqrt{3}x+6}{\sqrt{3}-x} 
			=3 \sqrt{3}. 
			\] 
			\itemch \textbf{Đúng.} Nếu $a=1$  và $b=0$ thì $\lim\limits_{x \to 2}g(x)=\lim\limits_{x \to 2}\dfrac{x^2+x}{x-1}=5$.
			\itemch \textbf{Đúng.} Ta có $\lim\limits_{x \to  -\sqrt{3}}f(x)=a \sqrt{3}+b=3 \sqrt{3}$.
			\\ 
			Suy ra $a=3$ và  $b=0$.
			 \\
			 Vậy $a+b=3$.
			\itemch \textbf{Đúng.} Nếu $a+b+1 \neq  0$ thì $\lim\limits_{x \to 1^{-}}\dfrac{x^2+a x+b}{x-1}\neq \lim\limits_{x \to 1^{+}}\dfrac{x^2+a x+b}{x-1}$.
			 \\
			Suy ra không tồn tại $\lim\limits_{x \to  1}\dfrac{x^2+a x+b}{x-1}$ (trái giả thiết).
			 \\
			 Do đó $a+b+1=0 \Rightarrow b=-a-1$.
			\\ 
			Thay vào giới  hạn, ta được:
			\[ 
			 \lim\limits_{x \rightarrow 1} \dfrac{x^2+a x-a-1}{x-1}
			=\lim\limits_{x \rightarrow  1} \dfrac{(x-1)(x+1+a)}{x-1}
			=\lim\limits_{x \rightarrow 1}(x+1+a) 
			 =a+2.
			\] 
			Khi đó,  để $\lim\limits_{x \to 1}\dfrac{x^2+a x+b}{x-1}=1$ thì
			 \[
			\heva{&a+2=1 \\ &b=-a-1} 
			 \Leftrightarrow
			 \heva{&a=-1 \\ &b=0.}
			\] 
		\end{itemchoice}
	 }
\end{ex}

% nguồn: Ok_dungsai_lop11_Lan2
\begin{ex}%[1D3H2-3]%[ma:b3921dbb]
	Cho giới hạn $I=\lim\limits_{x \to 3}\dfrac{a  x^2+b x+c}{x-3}$.
	 \choiceTF
		{\True Khi $a=0$, $b=1$, $c=-3$ thì  $I=1$}
		{Khi $a=1$, $b=-4$,  $c=3$ thì $I=0$}
		{\True  Khi $a=1$, $b=0$, $c=-9$ thì $I=6$}
		{\True Khi $a=1$, biết rằng có duy nhất một cặp số thực $(b;c)$ thỏa mãn $I=3$. Khi đó $ b+c=  -3$}
	 \loigiai{
		\begin{itemchoice}
			\itemch \textbf{Đúng.} Khi $a=0$, $b=1$, $c=-3$ thì $I=\lim\limits_{x \to  3}\dfrac{x-3}{x-3}=1$.
			\itemch \textbf{Sai.} Khi $a=1$, $b=-4$, $c=3$ thì  $I=\lim\limits_{x \to 3}\dfrac{(x-1)(x-3)}{x-3}=\lim\limits_{x \to 3}(x-1)=2$.
			 \itemch \textbf{Đúng.} Khi $a=1$, $b=0$, $c=-9$ thì
			 \[
			I=\lim\limits_{x \to 3}\dfrac{x^2-9}{x-3}=\lim\limits_{x \to 3}\dfrac{(x-3)(x+3)}{x-3}=\lim\limits_{x \to  3}(x+3)=6.
			\] 
			\itemch \textbf{Đúng.} Để  $\lim\limits_{x \to 3}\dfrac{x^2+b x+c}{x-3}=3$ thì ta phải có $x^2+b x+c=(x-3)(x-m)$.
			 \\
			Khi đó $3-m=3  \Leftrightarrow m=0$.
			\\ 
			Vậy $x^2+b x+c=(x-3)  x=x^2-3 x$. \\
			Suy  ra $b=-3$ và $c=0$.
		\end{itemchoice}
	} 
\end{ex}

% nguồn: Ok_dungsai_lop11_Lan2
\begin{ex}%[1D3H2-3]%[ma:6a29fc3c]
	 Cho hàm số $f(x)=x^2-3 x+2$.
	\choiceTF 
		{$\lim\limits_{x \to 2}\dfrac{f(x)}{x-2}=-1$} 
		{\True  $\lim\limits_{x \to 2}\dfrac{f(x)}{x^2-4}=\dfrac{1}{4}$}
		{\True $\lim\limits_{x \to 2}\dfrac{f(x)}{x^3-x^2-x-2}>0$} 
		{Để $\lim\limits_{x  \to 2}\dfrac{f(x)}{a x+b}=2$ thì $a+b=\dfrac{3}{2}$}
	 \loigiai{
		\begin{itemchoice}
			\itemch  \textbf{Sai.} Ta có
			 \[
			\lim\limits_{x \to  2}\dfrac{f(x)}{x-2}
			=\lim\limits_{x \to 2}\dfrac{x^2-3  x+2}{x-2}
			=\lim\limits_{x \to 2}\dfrac{(x-2)(x-1)}{x-2} 
			=\lim\limits_{x \to 2}(x-1) 
			 =1.
			 \]
			\itemch  \textbf{Đúng.} Ta có
			 \[
			\lim\limits_{x  \to 2}\dfrac{f(x)}{x^2-4}
			=\lim\limits_{x  \to 2}\dfrac{(x-2)(x-1)}{(x-2)(x+2)}
			=\lim\limits_{x \to 2}\dfrac{x-1}{x+2} 
			=\dfrac{1}{4}. 
			 \]
			\itemch \textbf{Đúng.} Ta có 
			\[ 
			\lim\limits_{x  \to 2}\dfrac{f(x)}{x^3-x^2-x-2}
			=\lim\limits_{x  \to 2}\dfrac{(x-2)(x-1)}{(x-2)\left(x^2+x+1\right)}
			=\lim\limits_{x  \to 2}\dfrac{x-1}{x^2+x+1}
			=\dfrac{1}{7}>0. 
			 \]
			\itemch \textbf{Sai.} Để $\lim\limits_{x \to 2}\dfrac{f(x)}{a x+b}=2$ thì 
			\begin{itemize}
				 \item $a x+b$ có nghiệm bằng $2 \Leftrightarrow 2 a+b=0 \Leftrightarrow b=-2 a$.
				\item  Khi đó
				\[ 
				\lim\limits_{x \to 2}\dfrac{f(x)}{a  x+b}
				=\lim\limits_{x \to 2}\dfrac{(x-2)(x-1)}{a(x-2)} 
				=\lim\limits_{x \to 2}\dfrac{x-1}{a} 
				 =\dfrac{1}{a}=2.
				\] 
				 Suy ra
				\[ 
				a=\dfrac{1}{2} \Rightarrow b=-1. 
				 \]
			\end{itemize}
			 Vậy $a+b=-\dfrac{1}{2}$.
		\end{itemchoice}
	} 
\end{ex}

% nguồn: Ok_dungsai_lop11_Lan2
\begin{ex}%[1D3H2-3]%[ma:85b9c960]
	Cho hàm số $f(x)=\dfrac{\sqrt{4 x^2+1}}{x+1}$ và $g(x)=\dfrac{2 x}{x+1}$. 
	 \choiceTF
		{$\lim\limits_{x \to -\infty}f(x)=2$ và $\lim\limits_{x \to -\infty}g(x)=2$} 
		{$\lim\limits_{x  \to -\infty}[f(x)-g(x)]=0$}
		{\True $\lim\limits_{x \to  +\infty}g(x)=\lim\limits_{x \to +\infty}f(x)=2$}
		 {$\lim\limits_{x \to +\infty}\dfrac{f(x)}{g(x)}=0$}
	\loigiai{ 
		\begin{itemchoice}
			\itemch  \textbf{Sai.} Ta có
			\[ 
			\lim\limits_{x \to  -\infty}f(x)
			=\lim\limits_{x \to  -\infty}\dfrac{\sqrt{4 x^2+1}}{x+1}
			=\lim\limits_{x \to  -\infty}\dfrac{-x \sqrt{4+\dfrac{1}{x^2}}}{x+1}
			=\lim\limits_{x \to -\infty}\dfrac{-\sqrt{4+\dfrac{1}{x^2}}}{1+\dfrac{1}{x}} 
			 =-2.
			 \]
			 Và
			 \[
			\lim\limits_{x \to  -\infty}g(x)
			=\lim\limits_{x \to  -\infty}\dfrac{2 x}{x+1}
			=\lim\limits_{x  \to -\infty}\dfrac{2}{1+\dfrac{1}{x}}
			 =2.
			 \]
			\itemch \textbf{Sai.} Ta  có
			 \[
			\lim\limits_{x \to  -\infty}[f(x)-g(x)]
			=\lim\limits_{x \to  -\infty}f(x)-\lim\limits_{x \to -\infty}g(x)
			=-2-2 
			 =-4.
			\] 
			\itemch \textbf{Đúng.} Ta  có
			\[ 
			 \lim\limits_{x \to +\infty}g(x)
			=\lim\limits_{x  \to +\infty}\dfrac{2 x}{x+1}
			=\lim\limits_{x  \to +\infty}\dfrac{2}{1+\dfrac{1}{x}}
			 =2.
			\] 
			 Và
			\[ 
			\lim\limits_{x  \rightarrow+\infty} f(x)
			=\lim\limits_{x \rightarrow+\infty}  \dfrac{\sqrt{4 x^2+1}}{x+1}
			=\lim\limits_{x \rightarrow+\infty} \dfrac{x  \sqrt{4+\dfrac{1}{x^2}}}{x+1}
			=\lim\limits_{x \rightarrow-\infty} \dfrac{\sqrt{4+\dfrac{1}{x^2}}}{1+\dfrac{1}{x}}=2. 
			\] 
			\itemch \textbf{Sai.} Ta  có
			\[ 
			\lim\limits_{x \to +\infty}\dfrac{f(x)}{g(x)} 
			=\dfrac{\lim\limits_{x  \rightarrow+\infty}f(x)}{\lim\limits_{x \to +\infty}g(x)}
			=\dfrac{2}{2} 
			 =1.
			 \]
		\end{itemchoice}
	} 
\end{ex}

% nguồn: Ok_dungsai_lop11_Lan2
\begin{ex}%[1D3H2-3]%[ma:d85360e6]
	Cho hàm số $f(x)=x^2-5 x+4$. 
	\choiceTF 
		{\True $\lim\limits_{x \rightarrow 4} \dfrac{f(x)}{\sqrt{x}-2}=12$} 
		{$\lim\limits_{x \rightarrow 4} \dfrac{f(x)}{\sqrt{x^2-6  x+9}-1}<0$}
		{$\lim\limits_{x \rightarrow 4}  \dfrac{f(x)-\sqrt{x}+2}{\sqrt{x+5}-3}>18$}
		{\True Để  $\lim\limits_{x \rightarrow 4}\dfrac{\sqrt{f(x)+9}-3}{2}=2$ thì $4 a+2 b=-1$}
	\loigiai{ 
		\begin{itemchoice}
			\itemch \textbf{Đúng.} Ta có 
			\[ 
			 \lim\limits_{x \to 4}\dfrac{f(x)}{\sqrt{x}-2}
			=\lim\limits_{x \to  4}\dfrac{(x-1)(x-4)(\sqrt{x}+2)}{x-4}
			=\lim\limits_{x \to 4}(x-1)(\sqrt{x}+2) 
			=12. 
			\] 
			\itemch \textbf{Sai.} Ta có 
			\begin{align*}
				\lim\limits_{x  \to 4}\dfrac{f(x)}{\sqrt{x^2-6 x+9}-1}
				&=\lim\limits_{x \to 4}\dfrac{(x-1)(x-4)\left(\sqrt{x^2-6 x+9}+1\right)}{x^2-6 x+8} \\ 
				&=\lim\limits_{x \to 4}\dfrac{(x-1)(x-4)\left(\sqrt{x^2-6 x+9}+1\right)}{(x-4)(x-2)}  \\
				&=\lim\limits_{x \to 4}\dfrac{(x-1)\left(\sqrt{x^2-6  x+9}+1\right)}{x-2} \\
				&=3. 
			\end{align*}
			\itemch \textbf{Sai.}  Ta có
			\begin{align*}
				 \lim\limits_{x \rightarrow 4} \dfrac{f(x)-\sqrt{x}+2}{\sqrt{x+5}-3}
				 &=\lim\limits_{x \rightarrow 4}\left[\dfrac{f(x)}{\sqrt{x+5}-3}-\dfrac{\sqrt{x}-2}{\sqrt{x+5}-3}\right] \\
				&=\lim\limits_{x \rightarrow 4}\left[\dfrac{(x-1)(x-4)(\sqrt{x+5}+3)}{x-4}-\dfrac{(x-4)(\sqrt{x+5}+3)}{(x-4)(\sqrt{x}+2)}\right] \\ 
				&=\lim\limits_{x \rightarrow 4}\left[(x-1)(\sqrt{x+5}+3)-\dfrac{(\sqrt{x+5}+3)}{(\sqrt{x}+2)}\right]  \\
				&=18-\dfrac{3}{2} \\ 
				&=\dfrac{33}{2}. 
			\end{align*}
			\itemch  \textbf{Đúng.} Để $\lim\limits_{x \rightarrow 4} \dfrac{\sqrt{f(x)+9}-3}{a x+b}=2$ thì
			\begin{itemize}
				\item $a x+b$ có  nghiệm bằng $4 \Leftrightarrow 4 a+b=0 \Leftrightarrow b=-4 a$.
				\item  Ta thấy
				\begin{align*}
					\lim\limits_{x \rightarrow 4} \dfrac{\sqrt{f(x)+9}-3}{a  x+b}
					&=\lim\limits_{x \rightarrow 4} \dfrac{f(x)}{a(x-4)(\sqrt{f(x)+9}+3)} \\ 
					&=\lim\limits_{x \rightarrow 4}  \dfrac{(x-1)(x-4)}{a(x-4)(\sqrt{f(x)+9}+3)} \\
					&=\lim\limits_{x \rightarrow 4} \dfrac{x-1}{a(\sqrt{f(x)+9}+3)} \\ 
					&=\dfrac{3}{6  a} \\
					 &=2.
				\end{align*}
				Khi đó  $a=\dfrac{1}{4}\Rightarrow b=-1 \Rightarrow 4 a+2 b=-1$.
			\end{itemize}
		\end{itemchoice}
	} 
\end{ex}

% nguồn: Ok_dungsai_lop11_Lan2
\begin{ex}%[1D3H2-3]%[ma:33804cf7]
	Cho hàm số $f(x)=\dfrac{2 x^m-3 x+2}{3 x^n+5}$, với $m,  n \in \mathbb{N}$.
	\choiceTF 
		{Khi $m=n=2$ thì $\lim\limits_{x  \to +\infty}f(x)=+\infty$}
		{\True Khi $m=2,  n=3$ thì $\lim\limits_{x \to +\infty}f(x)=0$}
		{\True Khi  $m=3, n=2$ thì $\lim\limits_{x \to +\infty}f(x)=+\infty$}
		 {Khi $1<m<n$ thì $\lim\limits_{x \to +\infty}f(x)=+\infty$}
	 \loigiai{
		\begin{itemchoice}
			\itemch \textbf{Sai.}  Khi $m=n=2$ thì
			 \[
			\lim\limits_{x \to +\infty}f(x)=\lim\limits_{x \to +\infty}\dfrac{2 x^2-3 x+2}{3 x^2+5}=\lim\limits_{x \to  +\infty}\dfrac{2-\dfrac{3}{x}+\dfrac{2}{x^2}}{3+\dfrac{5}{x^2}}=\dfrac{2}{3}.
			\] 
			\itemch \textbf{Đúng.}  Khi $m=2, n=3$ thì
			\[ 
			\lim\limits_{x \to +\infty}f(x)=\lim\limits_{x  \to +\infty}\dfrac{2 x^2-3 x+2}{3 x^3+5}=\lim\limits_{x \to +\infty}\dfrac{\dfrac{2}{x}-\dfrac{3}{x^2}+\dfrac{2}{x^3}}{3+\dfrac{5}{x^3}}=0.
			 \]
			\itemch  \textbf{Đúng.} Khi $m=3, n=2$ thì
			\[ 
			\lim\limits_{x \to  +\infty}f(x)=\lim\limits_{x \to +\infty}\dfrac{2-\dfrac{3}{x^2}+\dfrac{2}{x^3}}{\dfrac{3}{x}+\dfrac{5}{x^3}}=+\infty.
			\] 
			 \itemch \textbf{Sai.} Khi $1<m<n$ thì
			\begin{align*}
				\lim\limits_{x  \to +\infty}f(x)
				&=\lim\limits_{x \to +\infty}\dfrac{2 x^m-3 x+2}{3 x^n+5}  \\
				&=\lim\limits_{x  \to +\infty}\dfrac{\dfrac{2 x^m}{x^n}-\dfrac{3 x}{x^n}+\dfrac{2}{x^n}}{\dfrac{3 x^n}{x^n}+\dfrac{5}{x^n}} \\
				 &=\lim\limits_{x \to +\infty}\dfrac{\dfrac{2}{x^{n-m}}-\dfrac{3}{x^{n-1}}+\dfrac{2}{x^n}}{3+\dfrac{5}{x^n}}\\
				 &=0.
			\end{align*}
		\end{itemchoice}
	} 
\end{ex}

% nguồn: Ok_dungsai_lop11_Lan1
\begin{ex}%[1D3V2-3]%[ma:52ab900f]
 Cho hàm số $f(x)=\dfrac{2 x^m-3 x+2}{3 x^n+5}$, với $m, n \in \mathbb{R}$. 
  \choiceTF
  {Khi $m=n=2$ thì $\lim \limits_{x  \rightarrow+\infty} f(x)=+\infty$}
  {\True Khi $m=2, n=3$ thì $\lim  \limits_{x \rightarrow+\infty} f(x)=0$}
  {\True Khi $m=3, n=2$ thì $\lim \limits_{x \rightarrow+\infty}  f(x)=+\infty$}
  {Khi $1<m<n$  thì $\lim \limits_{x \rightarrow+\infty} f(x)=+\infty$}
 \loigiai{ 
  \begin{itemchoice}
   \itemch Sai. Khi $m=n=2$ thì $\lim \limits_{x \rightarrow+\infty} f(x)=\lim \limits_{x \rightarrow+\infty} \dfrac{2 x^2-3 x+2}{3 x^2+5}=\lim \limits_{x \rightarrow+\infty}  \dfrac{2-\dfrac{3}{x}+\dfrac{2}{x^2}}{3+\dfrac{5}{x^2}}=\dfrac{2}{3}$.
   \itemch  Đúng. Khi $m=2, n=3$ thì
   \begin{align*}
    \lim \limits_{x \rightarrow+\infty} f(x)=\lim \limits_{x \rightarrow+\infty} \dfrac{2 x^2-3 x+2}{3 x^3+5}=\lim \limits_{x \rightarrow+\infty}  \dfrac{\dfrac{2}{x}-\dfrac{3}{x^2}+\dfrac{2}{x^3}}{3+\dfrac{5}{x^3}}=0.
   \end{align*}
   \itemch Đúng. Khi $m=3, n=2$ thì $\lim \limits_{x \rightarrow+\infty}  f(x)=\lim \limits_{x \rightarrow+\infty} \dfrac{2-\dfrac{3}{x^2}+\dfrac{2}{x^3}}{\dfrac{3}{x}+\dfrac{5}{x^3}}=+\infty$.
   \itemch Sai. Khi  $1<m<n$ thì
   \begin{align*}
    \lim \limits_{x \rightarrow+\infty} f(x)=&\lim \limits_{x \rightarrow+\infty}  \dfrac{2 x^m-3 x+2}{3 x^n+5}\\
    =&\lim \limits_{x  \rightarrow+\infty} \dfrac{\dfrac{2 x^m}{x^n}-\dfrac{3 x}{x^n}+\dfrac{2}{x^n}}{\dfrac{3 x^n}{x^n}+\dfrac{5}{x^n}}\\
    =&\lim \limits_{x \rightarrow+\infty}  \dfrac{\dfrac{2}{x^{n-m}}-\dfrac{3}{x^{n-1}}+\dfrac{2}{x^n}}{3+\dfrac{5}{x^n}}=0.
   \end{align*}
  \end{itemchoice}
  }
\end{ex}


%%%%%%%%%%%%%%%% TRẢ LỜI NGẮN %%%%%%%%%%%%%%%%

% nguồn: Ok_TraLoiNgan_Lop11_Lan2
\begin{ex}%[1D3H2-3]%[ma:d857a7e2]
	Biết $\lim\limits_{x \to 2}\dfrac{x^2+3 x-10}{x^2+2 x-3}=\dfrac{a}{b}$; $a, b \in \mathbb{Z}$; $b \neq 0$;  $\dfrac{a}{b}$ là phân số tối giản. Giá trị của $a \cdot b$ bằng bao nhiêu?
	 \shortans{7}
	\loigiai{ 
		Ta  có
		 \[
		\lim\limits_{x \to 2}\dfrac{x^2+3  x-10}{x^2-3 x+2}
		 =\lim\limits_{x \to 2}\dfrac{(x-2)(x+5)}{(x-2)(x-1)}
		=\lim\limits_{x  \to 2}\dfrac{x+5}{x-1}
		 =\dfrac{7}{1}
		 =\dfrac{14}{2}
		 =7.
		\] 
		Do  đó, giá trị của $a \cdot b$ bằng $7$.
	} 
\end{ex}

% nguồn: Ok_TraLoiNgan_Lop11_Lan1
\begin{ex}%[1D3V2-3]%[ma:7f8a1bae]
	 Tính giới hạn $\lim\limits_{x\to+\infty}\dfrac{(3-2x)^5\cdot\left(3x^2-2x\right)}{\left(2x^2-1\right)^2(7-2x)^3}$
	 \shortans{3}
	\loigiai{ 
		\begin{eqnarray*}
			\lim\limits_{x\to+\infty}\dfrac{(3-2x)^5\cdot\left(3x^2-2x\right)}{\left(2x^2-1\right)^2(7-2x)^3}&=& \lim\limits_{x\to+\infty}\dfrac{x^5\left(\dfrac{3}{x}-2\right)^5\cdot{x^2}\left(3-\dfrac{2}{x}\right)}{x^4\left(2-\dfrac{1}{x^2}\right)^2\cdot{x^3}{\left(\dfrac{7}{x}-2\right)^3}}\\ 
			&=&\lim\limits_{x\to+\infty}\dfrac{\left(\dfrac{3}{x}-2\right)^5\cdot\left(3-\dfrac{2}{x}\right)}{\left(2-\dfrac{1}{x^2}\right)^2\cdot{\left(\dfrac{7}{x}-2\right)^3}}=3. 
		\end{eqnarray*}
	} 
\end{ex}

% nguồn: Ok_TraLoiNgan_Lop11_Lan2
\begin{ex}%[1D3V2-3]%[ma:1363850e]
	Biết rằng $\lim\limits_{x \to +\infty}\left(\dfrac{x^2+1}{x-2}+a  x-b\right)=-5$. Tính tổng $a+b$.
	 \shortans{6}
	 \loigiai{
		 Ta có
		\[ 
		\lim\limits_{x \to +\infty}\left(\dfrac{x^2+1}{x-2}+a x-b\right)=\lim\limits_{x \to +\infty}\left(\dfrac{(a+1)  x^2-(2 a+b) x+2 b+1}{x-2}\right)=-5.
		\] 
		\begin{itemize}
			\item Khi  $a+1>0 \Leftrightarrow a>-1$, ta được
			\[ 
			\lim\limits_{x  \to +\infty}\left(\dfrac{(a+1) x^2-(2 a+b) x+2 b+1}{x-2}\right)=+\infty.
			\] 
			\item Khi $a+1<0 \Leftrightarrow  a<-1$, ta được
			 \[
			\lim\limits_{x \to +\infty}\left(\dfrac{(a+1) x^2-(2 a+b) x+2 b+1}{x-2}\right)=-\infty. 
			 \]
			 \item Khi $a+1=0 \Leftrightarrow a=-1$, ta được
			 \[
			\lim\limits_{x  \to +\infty}\left(\dfrac{(a+1) x^2-(2 a+b) x+2 b+1}{x-2}\right)
			=\lim\limits_{x  \to +\infty}\left(\dfrac{-(-2+b) x+2 b+1}{x-2}\right)
			=2-b. 
			\] 
		\end{itemize}
		Do đó, $2-b=-5 \Leftrightarrow b=7$.  Vậy $a+b=6$.
	} 
\end{ex}

% nguồn: Ok_TraLoiNgan_Lop11_Lan2
\begin{ex}%[1D3V2-3]%[ma:af613ee2]
	 Để $\lim\limits_{x \to -\infty}\dfrac{\sqrt{4 x^2+x+1}+4}{m x-2}=\dfrac{1}{2}$. Tìm giá trị của $m$.
	 \shortans{-4}
	\loigiai{ 
		Xét  $2$ trường hợp
		\begin{itemize}
			\item  Khi $m=0$, ta có
			 \[
			\lim\limits_{x \to -\infty}\dfrac{\sqrt{4 x^2+x+1}+4}{-2} 
			=\lim\limits_{x \to -\infty}x  \cdot \dfrac{\sqrt{4+\dfrac{1}{x}+\dfrac{1}{x^2}}+\dfrac{4}{x}}{2}
			=-\infty 
			 \]
			 nên $m=0$ không thỏa mãn.
			\item  Khi $m \neq 0$, ta có
			\[ 
			\lim\limits_{x \to -\infty}\dfrac{\sqrt{4 x^2+x+1}+4}{m  x-2}
			=\lim\limits_{x \to -\infty}\dfrac{-\sqrt{4+\dfrac{1}{x}+\dfrac{1}{x^2}}+\dfrac{4}{x}}{m-\dfrac{2}{x}} 
			=-\dfrac{2}{m}. 
			\] 
			Theo đề bài,  ta có
			\[ 
			 -\dfrac{2}{m}=\dfrac{1}{2}
			\Leftrightarrow  m=-4.
			\] 
		\end{itemize}
	} 
\end{ex}

% nguồn: Ok_TraLoiNgan_Lop11_Lan2
\begin{ex}%[1D3V2-3]%[ma:25c9df53]
	Cho 
	$\lim\limits_{x \to -\infty}\dfrac{a \sqrt{x^2+2}+2023}{x+2024}=\dfrac{1}{2}$  và
	$\lim\limits_{x \to  +\infty}\left(\sqrt{x^2+b x+3}-x\right)=2$.
	Tính giá trị của $P=4 a+b$. 
	 \shortans{2}
	\loigiai{ 
		Ta  có
		\begin{align*}
			\lim\limits_{x \to  -\infty}\dfrac{a \sqrt{x^2}+2+2023}{x+2024}
			 &=\lim\limits_{x \to -\infty}\dfrac{x\left(-a \sqrt{1+\dfrac{2}{x^2}}+\dfrac{2023}{x}\right)}{x\left(1+\dfrac{2024}{x}\right)} \\
			&=\lim\limits_{x  \to -\infty}\dfrac{-a \sqrt{1+\dfrac{2}{x^2}}+\dfrac{2023}{x}}{1+\dfrac{2024}{x}} \\
			 &=-a.
		\end{align*}
		Theo giả  thiết
		\[ 
		 -a=\dfrac{1}{2}\Leftrightarrow a=-\dfrac{1}{2}.
		\] 
		Ta  có
		\begin{align*}
			 \lim\limits_{x \to +\infty}\left(\sqrt{x^2}+b x+3-x\right)
			&=\lim\limits_{x \to +\infty}\dfrac{\left(\sqrt{x^2}+b x+3-x\right)\left(\sqrt{x^2}+b x+3+x\right)}{\sqrt{x^2}+b  x+3+x} \\
			&=\lim\limits_{x \to +\infty}\dfrac{b x+3}{x\left(\sqrt{1+\dfrac{b}{x}+\dfrac{3}{x^2}}+1\right)}  \\
			&=\lim\limits_{x \to +\infty}\dfrac{x\left(b+\dfrac{3}{x}\right)}{x\left(\sqrt{1+\dfrac{b}{x}+\dfrac{3}{x^2}}+1\right)}  \\
			&=\lim\limits_{x \to  +\infty}\dfrac{b+\dfrac{3}{x}}{\sqrt{1+\dfrac{b}{x}+\dfrac{3}{x^2}}+1} \\
			 &=\dfrac{b}{2}.
		\end{align*}
		Nên  $\dfrac{b}{2}=2 \Leftrightarrow b=4$.
		 \\
		Vậy $P=4\left(-\dfrac{1}{2}\right)+4=2$. 
	} 
\end{ex}

% nguồn: Ok_TraLoiNgan_Lop11_Lan2
\begin{ex}%[1D3V2-3]%[ma:6c46fa5f]
	Cho $C=\lim\limits_{x \to 1}\dfrac{x^2-m x+m-1}{x^2-1}$ với $m$ là tham số thực. 
	Tìm $m$ để $C=2$. 
	 \shortans{-2}
	\loigiai{ 
		\begin{align*}
			C&=\lim\limits_{x \rightarrow 1} \dfrac{x^2-m  x+m-1}{x^2-1} \\
			&=\lim\limits_{x \rightarrow 1}  \dfrac{\left(x^2-1\right)-m(x-1)}{(x-1)(x+1)} \\
			&=\lim\limits_{x  \rightarrow 1} \dfrac{(x-1)(x+1-m)}{(x-1)(x+1)} \\
			&=\lim\limits_{x \rightarrow  1} \dfrac{x+1-m}{x+1} \\
			&=\dfrac{2-m}{2}. 
		\end{align*}
		Khi đó $C=2 \Leftrightarrow \dfrac{2-m}{2}=2 \Leftrightarrow  m=-2$.
	 }
\end{ex}


%%%%%%%%%%%%%%%% TỰ LUẬN %%%%%%%%%%%%%%%%

% nguồn: Ok_TuanLuan_Toan11_Lan1
\begin{bt}%[1D3H2-3]%[1D3V2-5]%[ma:314f439c]
	\textbf{(2.0 điểm)} Tính các giới hạn  sau:
	\begin{multicols}{2}
		\begin{enumerate}
			 \item$\underset{x\to 1}{\mathop{\lim }}\,\dfrac{{{x}^{2}}-3x+2}{{{x}^{2}}-1}$
			\item $\underset{x\to -\infty }{\mathop{\lim }}\,\dfrac{\sqrt{4{{x}^{2}}+5}-2x-1}{x-1}$ 
		\end{enumerate}
	\end{multicols}
	\loigiai{ 
		\begin{enumerate}
			 \item
			$\mathop {\lim }\limits_{x \to 1} \dfrac{{{x^2} - 3x + 2}}{{{x^2} - 1}}= \mathop {\lim }\limits_{x \to  1} \dfrac{{\left( {x - 1} \right)\left( {x - 2} \right)}}{{\left( {x - 1} \right)\left( {x + 1} \right)}}=\mathop {\lim }\limits_{x \to 1} \dfrac{{x - 2}}{{x + 1}} = - \dfrac{1}{2}$.
			\item 
			$\mathop {\lim }\limits_{x \to - \infty } \dfrac{{\sqrt {4{x^2} + 5} - 2x - 1}}{{x - 1}}=\mathop {\lim }\limits_{x \to - \infty } \dfrac{{ - x\sqrt {4 + \dfrac{5}{{{x^2}}}} - 2x - 1}}{{x - 1}}=\mathop {\lim }\limits_{x \to - \infty } \dfrac{{x\left( { - \sqrt {4 + \dfrac{5}{{{x^2}}}} - 2 - \dfrac{1}{x}} \right)}}{{x\left( {1 - \dfrac{1}{x}} \right)}} =\mathop  {\lim }\limits_{x \to - \infty } \dfrac{{ - \sqrt {4 + \dfrac{5}{{{x^2}}}} - 2 - \dfrac{1}{x}}}{{1 - \dfrac{1}{x}}} = - 4$.
		\end{enumerate}
	 }
\end{bt}

% nguồn: Ok_TuanLuan_Toan11_Lan2
\begin{ex}%[1D3H2-3]%[1D3V2-5]%[ma:a43ef871]
	\textbf{(2.0 điểm)}  Tính các giới hạn sau:
	\begin{multicols}{2}
		\begin{enumerate}
			\item$\underset{x\to 1}{\mathop{\lim }}\,\dfrac{{{x}^{2}}-3x+2}{{{x}^{2}}-1}$ 
			 \item $\underset{x\to -\infty }{\mathop{\lim }}\,\dfrac{\sqrt{4{{x}^{2}}+5}-2x-1}{x-1}$
		\end{enumerate}
	\end{multicols}
	\loigiai{ 
		\begin{enumerate}
			\item 
			$\mathop {\lim }\limits_{x \to 1} \dfrac{{{x^2} - 3x + 2}}{{{x^2} - 1}}= \mathop {\lim }\limits_{x \to 1} \dfrac{{\left( {x - 1} \right)\left( {x - 2}  \right)}}{{\left( {x - 1} \right)\left( {x + 1} \right)}}=\mathop {\lim }\limits_{x \to 1} \dfrac{{x - 2}}{{x + 1}} = - \dfrac{1}{2}$.
			 \item
			$\mathop {\lim }\limits_{x \to - \infty  } \dfrac{{\sqrt {4{x^2} + 5} - 2x - 1}}{{x - 1}}=\mathop {\lim }\limits_{x \to - \infty } \dfrac{{ - x\sqrt {4 + \dfrac{5}{{{x^2}}}} - 2x - 1}}{{x - 1}}=\mathop {\lim }\limits_{x \to - \infty } \dfrac{{x\left( { - \sqrt {4 + \dfrac{5}{{{x^2}}}} - 2 - \dfrac{1}{x}} \right)}}{{x\left( {1 - \dfrac{1}{x}} \right)}} =\mathop {\lim }\limits_{x \to - \infty } \dfrac{{ - \sqrt {4 + \dfrac{5}{{{x^2}}}} - 2 - \dfrac{1}{x}}}{{1 - \dfrac{1}{x}}} = - 4$.
		\end{enumerate}
	 }
\end{ex}

% nguồn: Ok_TuanLuan_Toan11_Lan2
\begin{ex}%[1D3H2-3]%[ma:1112cd08]
	Tính giá trị của giới hạn  $ \lim\limits_{x\to +\infty}\dfrac{\sqrt{4x^2-2x+1}+2-x}{\sqrt{9x^2-3x}+2x}$.
	 \loigiai{
	$ \lim\limits_{x\to +\infty  }\dfrac{\sqrt{4x^2-2x+1}+2-x}{\sqrt{9x^2-3x}+2x}= \lim\limits_{x\to +\infty }\dfrac{\sqrt{4-\dfrac{2}{x}+\dfrac{1}{x^2}}+\dfrac{2}{x}-1}{\sqrt{9-\dfrac{3}{x}}+2}=\dfrac{1}{5}$.
	 }
\end{ex}

% nguồn: Ok_TuanLuan_Toan11_Lan2
\begin{ex}%[1D3H2-3]%[ma:b67db890]
	Tính giá  trị của giới hạn $\lim\limits_{x \to -\infty}\dfrac{\sqrt{4x^2-x+1}}{x+1}$.
	 \loigiai{
	 $ \lim\limits_{x\to -\infty }\dfrac{\sqrt{4x^2-x+1}}{x+1}= \lim\limits_{x\to -\infty }\dfrac{-\sqrt{4-\dfrac{1}{x}+\dfrac{1}{x^2}}}{1+\dfrac{1}{x}}=\dfrac{-\sqrt{4}}{1}=-2$.
	} 
\end{ex}

% nguồn: Ok_TuanLuan_Toan11_Lan2
\begin{ex}%[1D3H2-3]%[ma:c67eed4b]
	Tính  giới hạn $\lim\limits_{x\to -\infty}\dfrac{\sqrt{x^2-x+1}}{2x}$.
	\loigiai{ 
		Ta có $\lim\limits_{x\to -\infty}\dfrac{\sqrt{x^2-x+1}}{2x}=\lim\limits_{x\to +\infty}\dfrac{-x\sqrt{1-\dfrac{1}{x}+\dfrac{1}{x^2}}}{2x}=\lim\limits_{x\to +\infty}\dfrac{-\sqrt{1-\dfrac{1}{x}+\dfrac{1}{x^2}}}{2}=-\dfrac{1}{2}$. 
	} 
\end{ex}

% nguồn: Ok_TuanLuan_Toan11_Lan2
\begin{ex}%[1D3H2-3]%[ma:0cb48083]
	 Giá trị của giới hạn $\lim\limits_{x\to-1}\dfrac{x^5+1}{x^3+1}$ là
	\loigiai{ 
	 $\lim\limits_{x\to-1}\dfrac{x^5+1}{x^3+1}=\lim\limits_{x\to-1}\dfrac{\left(x+1\right)\left(x^4-x^3+x^2-x+1\right)}{(x+1)(x^2-x+1)}=\lim\limits_{x\to-1}\dfrac{x^4-x^3+x^2-x+1}{x^2-x+1}=\dfrac{5}{3}$.}
\end{ex}

% nguồn: Ok_TuanLuan_Toan11_Lan2
\begin{ex}%[1D3H2-3]%[ma:54813b2c]
	Tính giới hạn $\lim\limits_{x\to -\infty}\dfrac{5x^2+2x+3}{x^2+1}$. 
	\loigiai{ 
		 Ta có $\lim\limits_{x\to -\infty}\dfrac{5x^2+2x+3}{x^2+1}=\lim\limits_{x\to -\infty}\dfrac{5+\dfrac{2}{x}+\dfrac{3}{x^2}}{1+\dfrac{1}{x^2}}=5$.
	} 
\end{ex}

% nguồn: Ok_TuanLuan_Toan11_Lan2
\begin{ex}%[1D3H2-3]%[ma:082ef5a3]
	Tính giá trị của giới hạn  $ \lim\limits_{x\to -\infty }\dfrac{2x-3}{\sqrt{x^2+1}-x}$.
	\loigiai{ 
	Chia cả tử  và mẫu cho $x$, ta được $ \lim\limits_{x\to -\infty }\dfrac{2x-3}{\sqrt{x^2+1}-x}= \lim\limits_{x\to -\infty }\dfrac{2-\dfrac{3}{x}}{-\sqrt{1+\dfrac{1}{x^2}}-1}=-1$.
	 }
\end{ex}

% nguồn: Ok_TuanLuan_Toan11_Lan2
\begin{ex}%[1D3H2-3]%[ma:cbfcfc60]
	Giá trị của giới hạn $\lim\limits_{x\to  2}\dfrac{x^3-8}{x^2-4}$ là
	\loigiai{ 
	Ta  có $\lim\limits_{x\to 2}\dfrac{x^3-8}{x^2-4}=\lim\limits_{x\to 2}\dfrac{(x-2)(x^2+2x+4)}{(x-2)(x+2)}=\lim\limits_{x\to 2}\dfrac{x^2+2x+4}{x+2}=\dfrac{12}{4}=3$.}
\end{ex}

% nguồn: Ok_TuanLuan_Toan11_Lan2
\begin{ex}%[1D3H2-3]%[ma:0e29c807]
	Biết rằng $\lim\limits_{x\to-\sqrt 3}\dfrac{2x^3+6\sqrt 3}{3-x^2}=a\sqrt  3$. Tính $a$.
	\loigiai{ 
		Ta có $\lim\limits_{x\to-\sqrt 3}\dfrac{2x^3+3\sqrt  3}{3-x^2}=\lim\limits_{x\to-\sqrt 3}\dfrac{2(x+\sqrt 3)(x^2-\sqrt 3 x+3)}{(\sqrt 3-x)(\sqrt 3+x)}=\lim\limits_{x\to-\sqrt 3}\dfrac{2(x^2-\sqrt 3 x+3)}{\sqrt 3-x}\\
	=\dfrac{2\left[\left(-\sqrt 3\right)^2-\sqrt 3 \cdot (-\sqrt 3)+3\right]}{\sqrt 3-(-\sqrt 3)}=\dfrac{18}{2\sqrt 3}=3\sqrt  3$.}
\end{ex}

% nguồn: Ok_TuanLuan_Toan11_Lan2
\begin{ex}%[1D3H2-3]%[ma:5534ae23]
	Tính giá trị của giới hạn $ \lim\limits_{x\to  -\infty }\dfrac{2x^3-7x^2+11}{3x^6+2x^5-5}$.
	 \loigiai{
	Ta có $\lim\limits_{x\to -\infty }\dfrac{2x^3-7x^2+11}{3x^6+2x^5-5}=  \lim\limits_{x\to -\infty }\dfrac{\dfrac{2}{x^3}-\dfrac{7}{x^4}+\dfrac{11}{x^6}}{3+\dfrac{2}{x}-\dfrac{5}{x^6}}=\dfrac{0}{3}=0$.
	} 
\end{ex}

% nguồn: Ok_TuanLuan_Toan11_Lan2
\begin{ex}%[1D3H2-3]%[ma:a554289a]
	 Tìm $\lim\limits_{x\to -\infty}\dfrac{\sqrt{x^2+3x+5}}{4x-1}$.
	\loigiai{ 
		Ta có $\lim\limits_{x\to -\infty}\dfrac{\sqrt{x^2+3x+5}}{4x-1}=\lim\limits_{x\to -\infty}\dfrac{-\sqrt{1+\dfrac{3}{x}+\dfrac{5}{x^2}}}{4-\dfrac{1}{x}}=-\dfrac{1}{4}$. 
	 }
\end{ex}

% nguồn: Ok_TuanLuan_Toan11_Lan2
\begin{ex}%[1D3H2-3]%[ma:2376a8ca]
	Tính giới hạn $K=\lim\limits_{x\to  -\infty}\dfrac{\sqrt{4x^2+1}}{x+1}$.
	 \loigiai{
		 Ta có $K=\lim\limits_{x\to -\infty}\dfrac{\sqrt{4x^2+1}}{x+1}=\lim\limits_{x\to -\infty}\dfrac{-x\sqrt{4+\dfrac{1}{x^2}}}{x+1}=\lim\limits_{x\to -\infty}\dfrac{-\sqrt{4+\dfrac{1}{x^2}}}{1+\dfrac{1}{x}}=-2$.
	 }
\end{ex}

% nguồn: Ok_TuanLuan_Toan11_Lan2
\begin{ex}%[1D3H2-3]%[ma:6c6ab3f2]
	 $\lim\limits_{x\to -\infty}\dfrac{x-\sqrt{x^2+x}}{x+1}$ bằng
	 \loigiai{
		Ta có  $\lim\limits_{x\to -\infty}\dfrac{x-\sqrt{x^2+x}}{x+1}=\lim\limits_{x\to -\infty}\dfrac{x+x\sqrt{1+\dfrac{1}{x}}}{x+1}=\lim\limits_{x\to -\infty}\dfrac{1+\sqrt{1+\dfrac{1}{x}}}{1+\dfrac{1}{x}}=2$.
	 }
\end{ex}

% nguồn: Ok_TuanLuan_Toan11_Lan2
\begin{ex}%[1D3H2-3]%[ma:e65d6e3d]

	Tính các giới hạn sau  đây
	\begin{multicols}{2}
		\begin{enumerate}
			\item $\displaystyle\lim_{x  \to 2} \dfrac{x^2 - 7x + 4}{x^2 - x - 20}$;
			\item $\displaystyle\lim_{x  \to 2} \dfrac{x^2 - 5x + 2}{x^2 - x - 6}$;
			\item $\displaystyle\lim_{x \to -4} \dfrac{x^2 + 2x - 8}{x + 4}$; 
			\item $\displaystyle\lim_{x \to  1} \dfrac{x^2 - 5x + 2}{x - 1}$;
			\item $\displaystyle\lim_{x \to -1} \dfrac{x^3 + 1}{x + 1}$; 
			\item $\displaystyle\lim_{x \to 2} \dfrac{x^2 - 5x + 6}{x^2 +  x - 6}$;
			\item $\displaystyle\lim_{x \to -1} \dfrac{x^3 - 4x^2 - x + 4}{x  + 1}$;
			\item $\displaystyle\lim_{x \to -1} \dfrac{4x^2 - 3x  - 7}{x + 1}$;
			\item $\displaystyle\lim_{x \to 2} \dfrac{(x^2 + x - 6)^2}{x^2  - 5x + 6}$.
		\end{enumerate}
	\end{multicols}
	 \loigiai{
		\begin{enumerate}
			\item Ta  có
			\begin{eqnarray*}
				\displaystyle\lim_{x \to 2} \dfrac{x^2 -  7x + 4}{x^2 - x - 20}
				&=& \dfrac{4 -  14 + 4}{4 - 2 - 20} = \dfrac{1}{3}.
			\end{eqnarray*}
			\item Ta  có
			\begin{eqnarray*}
				\displaystyle\lim_{x \to 2} \dfrac{x^2 -  5x + 2}{x^2 - x - 6}
				&=&  \dfrac{2^2 - 5 \cdot 2 + 2}{2^2 - 2 - 6}\\
				&=& \dfrac{4 - 10 + 2}{4 - 2 - 6} \\ 
				&=&1. 
			\end{eqnarray*}
			\item Ta  có
			\begin{eqnarray*}
				\displaystyle\lim_{x \to -4} \dfrac{x^2 + 2x  - 8}{x + 4}
				&=& \displaystyle\lim_{x \to -4} \dfrac{(x + 4)(x - 2)}{x + 4}\\ 
				&=& \displaystyle\lim_{x \to -4}  (x - 2)\\
				&=& -6. 
			\end{eqnarray*}
			\item Ta có 
			\begin{eqnarray*}
				\displaystyle\lim_{x  \to 1} \dfrac{x^2 - 5x + 2}{x - 1}
				&=& \displaystyle\lim_{x \to 1} \dfrac{(x - 1)(x - 2)}{x - 1}\\ 
				&=& \displaystyle\lim_{x \to 1}  (x - 2)\\
				&=&-1. 
			\end{eqnarray*}
			\item  Ta có
			\begin{eqnarray*}
				\displaystyle\lim_{x \to -1} \dfrac{x^3 + 1}{x + 1} 
				&=& \displaystyle\lim_{x \to -1}\dfrac{(x  + 1)(x^2 - x + 1)}{x + 1}\\
				&=& \displaystyle\lim_{x \to -1} (x^2 - x +  1)\\
				&=&  3.
			\end{eqnarray*}
			\item  Ta có
			\begin{eqnarray*}
				 \displaystyle\lim_{x \to 2} \dfrac{x^2 - 5x + 6}{x^2 + x - 6}
				&=& \dfrac{(x -  2)(x - 3)}{(x - 2)(x + 3)}\\
				 &=& \displaystyle\lim_{x \to 2} \dfrac{x - 3}{x + 3}\\
				&=&  \dfrac{-1}{5}.
			\end{eqnarray*}
			\item Ta có 
			\begin{eqnarray*}
				\displaystyle\lim_{x \to -1} \dfrac{x^3 - 4x^2 -  x + 4}{x + 1}
				&=&\displaystyle\lim_{x \to -1} \dfrac{(x+1)(x^2-5x+4)}{x  + 1} \\
				&=&\displaystyle\lim_{x  \to -1} (x^2-5x+4)\\
				&=&10. 
			\end{eqnarray*}
			\item  Ta có
			\begin{eqnarray*}
				\displaystyle\lim_{x \to -1} \dfrac{4x^2  - 3x - 7}{x + 1}
				&=& \displaystyle\lim_{x \to -1}  \dfrac{(x+1)(4x-7)}{x + 1} \\
				 &=& \displaystyle\lim_{x \to -1} (4x-7)\\
				 &=& -11.
			\end{eqnarray*}
			\item  Ta có
			\begin{eqnarray*}
				\displaystyle\lim_{x  \to 2} \dfrac{(x^2 + x - 6)^2}{x^2 - 5x + 6}
				&=&  \displaystyle\lim_{x\to 2} \dfrac{(x+3)^2\cdot (x-2)^2}{(x-2)(x-3)}\\
				&=&  \displaystyle\lim_{x\to 2} \dfrac{(x+3)^2\cdot (x-2)}{x-3}\\
				 &=& 0.
			\end{eqnarray*}
		\end{enumerate}
	} 
\end{ex}

% nguồn: Ok_TuanLuan_Toan11_Lan2
\begin{ex}%[1D3H2-3]%[ma:f572c256]
	 Tính $\lim\limits_{x\to 3}\dfrac{x^2-9}{x-3}$ bằng
	 \loigiai{
	Ta  có $\lim\limits_{x\to 3}\dfrac{x^2-9}{x-3}=\lim\limits_{x\to 3}\left(x+3\right)=6$.}
\end{ex}

% nguồn: Ok_TuanLuan_Toan11_Lan2
\begin{ex}%[1D3H2-3]%[ma:73ba5c40]
	 Tính giới hạn $I=\lim\limits_{x\to 2}\dfrac{x^2-5x+6}{x-2}$.
	 \loigiai{
	$I=\lim\limits_{x\to 2}\dfrac{x^2-5x+6}{x-2}=  \lim\limits_{x\to 2}\dfrac{\left(x-2\right)\left(x-3\right)}{x-2}= \lim\limits_{x\to 2}\left(x-3\right)=-1$.
	 }
\end{ex}

% nguồn: Ok_TuanLuan_Toan11_Lan2
\begin{ex}%[1D3H2-3]%[ma:aa1afed2]
	Tính $ \lim\limits_{x\to 5}\dfrac{x^2-12x+35}{25-5x}$. 
	\loigiai{ 
	Ta có $ \lim\limits_{x\to  5}\,\dfrac{x^2-12x+35}{25-5x}= \lim\limits_{x\to 5}\,\dfrac{(x-7)(x-5)}{5(5-x)}=\lim\limits_{x\to 5}\,\dfrac{7-x}{5}=\dfrac{2}{5}$.\\
	Vậy  $\lim\limits_{x\to 5}\,\dfrac{x^2-12x+35}{25-5x}=\dfrac{2}{5}$.
	} 
\end{ex}

% nguồn: Ok_TuanLuan_Toan11_Lan2
\begin{ex}%[1D3H2-3]%[1D3V2-6]%[ma:f7b63e89]

	Tính các giới hạn  sau đây
	\begin{multicols}{2}
		\begin{enumerate}
			\item $A=\displaystyle\lim_{x \to +\infty} \dfrac{2x^3-3x^2+4x+1}{x^4-5x^3+2x^2-x+3}$; 
			\item $B=\displaystyle\lim_{x \to  +\infty} \dfrac{x+\sqrt{x^2+2}}{\sqrt[3]{8x^3+x^2+1}}$;
			\item $C=\displaystyle\lim_{x\to  +\infty}\left( \sqrt{x+\sqrt{x}}-\sqrt{x}\right)$;
			 \item $D=\displaystyle\lim_{x\to +\infty} x \left(\sqrt{x^2+1}-x \right)$;
			\item  $E=\displaystyle\lim_{x\to +\infty} \dfrac{3x^2-x+7}{2x^3-1}$;
			\item $F=\displaystyle\lim_{x\to -\infty} \dfrac{2x^3-3x-4}{-x^3-x^2+1}$; 
			\item  $G=\displaystyle\lim_{x\to +\infty} \dfrac{3x^2-x+3}{x-4}$;
			\item $H=\displaystyle\lim_{x\to +\infty}  \dfrac{2x^3-2x+3}{5-x}$;
			 \item $I=\displaystyle\lim_{x\to \infty} \dfrac{(2x-3)^{20}(3x+2)^{30}}{(1+2x)^{50}}$;
			\item $J=\displaystyle\lim_{x\to +\infty}  \dfrac{\sqrt{x^2+x}+2x}{2x+3}$;
			\item $K=\displaystyle\lim_{x\to  -\infty} \dfrac{\sqrt{x^2+x}+2x}{2x+3}$;
			\item $L=\displaystyle\lim_{x\to  -\infty} \dfrac{8x+3}{6x+\sqrt{4x^2+x+3}}$.
		\end{enumerate}
	\end{multicols}
	\loigiai{ 
		\begin{enumerate}
			\item Ta  có
			\begin{eqnarray*}
				A &=& \displaystyle\lim_{x \to +\infty} \dfrac{2x^3 - 3x^2 + 4x + 1}{x^4 - 5x^3 + 2x^2  - x + 3} \\
				&=& \displaystyle\lim_{x \to +\infty} \dfrac{\dfrac{2}{x} - \dfrac{3}{x^2} + \dfrac{4}{x^3} +  \dfrac{1}{x^4}}{1 - \dfrac{5}{x} + \dfrac{2}{x^2} - \dfrac{1}{x^3} + \dfrac{3}{x^4}}\\
				&=& 0. 
			\end{eqnarray*}
			 \item Ta có
			\begin{eqnarray*}
				B &=& \displaystyle\lim_{x \to +\infty} \dfrac{x + \sqrt{x^2 + 2}}{\sqrt[3]{8x^3 + x^2 + 1}} \\ 
				&=& \displaystyle\lim_{x \to +\infty} \dfrac{x\left(1 + \sqrt{1 + \dfrac{2}{x^2}}\right)}{x  \cdot \sqrt[3]{8 + \dfrac{1}{x} + \dfrac{1}{x^3}}}\\
				&=&  1.
			\end{eqnarray*}
			 \item Ta có
			\begin{eqnarray*}
				C  &=& \displaystyle\lim_{x\to +\infty}\left(\sqrt{x+\sqrt{x}}-\sqrt{x}\right) \\
				&=& \displaystyle\lim_{x\to  +\infty}\dfrac{x+\sqrt{x}-x}{\left(\sqrt{x+\sqrt{x}}+\sqrt{x}\right)} \\
				&=& \displaystyle\lim_{x\to +\infty}\dfrac{\sqrt{x}}{\sqrt{x}\cdot\sqrt{1+\dfrac{1}{\sqrt{x}}}+\sqrt{x}}  \\
				&=&  \displaystyle\lim_{x\to +\infty}\dfrac{\sqrt{x}}{\sqrt{x}\left(\sqrt{1+\dfrac{1}{\sqrt{x}}}+1\right)}\\
				&=& \displaystyle\lim_{x\to  +\infty}\dfrac{1}{\sqrt{1+\dfrac{1}{\sqrt{x}}}+1}\\
				&=&1. 
			\end{eqnarray*}
			\item 
			\begin{eqnarray*}
				D &=& \displaystyle\lim_{x \to +\infty} x\left( \sqrt{x^2 +  1} - x \right) \\
				&=& \displaystyle\lim_{x \to +\infty}  \dfrac{x\left(x^2 + 1 - x^2\right)}{\sqrt{x^2 + 1} + x}\\
				&=& \displaystyle\lim_{x \to +\infty}  \dfrac{x}{\sqrt{x^2 + 1} + x} \\
				&=& \dfrac{1}{2}. 
			\end{eqnarray*}
			\item 
			\begin{eqnarray*}
				 E &=& \displaystyle\lim_{x \to +\infty} \dfrac{3x^2-x+7}{2x^3-1}\\
				&=& \displaystyle\lim_{x \to +\infty} \dfrac{\dfrac{3}{x} - \dfrac{1}{x^2} + \dfrac{7}{x^3}}{2 -  \dfrac{1}{x^3}}\\
				 &=& 0.
			\end{eqnarray*}
			 \item
			\begin{eqnarray*}
				F &=& \displaystyle\lim_{x \to -\infty} \dfrac{2x^3 - 3x -  4}{-x^3 - x^2 + 1} \\
				&=& \displaystyle\lim_{x \to -\infty} \dfrac{x^3 \left(2 - \dfrac{3}{x^2} - \dfrac{4}{x^3}\right)}{x^3  \left(-1 - \dfrac{1}{x} + \dfrac{1}{x^3}\right)}\\
				 &=& -2.
			\end{eqnarray*}
			\item 
			\begin{eqnarray*}
				 G &=& \displaystyle\lim_{x \to +\infty} \dfrac{3x^2-x+3}{x-4} \\
				&=& \displaystyle\lim_{x \to +\infty} \dfrac{x^2\left(3-\dfrac{1}{x}+\dfrac{3}{x^2}\right)}{x\left(1-\dfrac{4}{x}\right)}  \\
				&=& \displaystyle\lim_{x \to  +\infty} \dfrac{x\left(3-\dfrac{1}{x}+\dfrac{3}{x^2}\right)}{1-\dfrac{4}{x}} \\
				&=& +\infty. 
			\end{eqnarray*}
			Vì $\displaystyle\lim_{x \to +\infty} x=+\infty$ và $\displaystyle\lim_{x \to  +\infty} \dfrac{3-\dfrac{1}{x}+\dfrac{3}{x^2}}{1-\dfrac{4}{x}}=3>0$.
			\item 
			\begin{eqnarray*}
				H &=& \displaystyle\lim_{x \to +\infty}  \dfrac{2x^3-2x+3}{5-x} \\
				&=& \displaystyle\lim_{x  \to +\infty} \dfrac{x^3\left(2-\dfrac{2}{x^2}+\dfrac{3}{x^3}\right)}{x\left(\dfrac{5}{x}-1\right)} \\
				&=& \displaystyle\lim_{x \to +\infty} \dfrac{x^2\left(2-\dfrac{2}{x^2}+\dfrac{3}{x^3}\right)}{\dfrac{5}{x}-1} \\ 
				&=&  -\infty.
			\end{eqnarray*}
			Vì $\displaystyle\lim_{x \to +\infty} x^2=+\infty$ và  $\displaystyle\lim_{x \to +\infty} \dfrac{2-\dfrac{2}{x^2}+\dfrac{3}{x^3}}{\dfrac{5}{x}-1}=-2<0$.
			 \item
			\begin{eqnarray*}
				 I &=& \displaystyle\lim_{x \to +\infty} \dfrac{(2x-3)^{20}(3x+2)^{30}}{(1+2x)^{50}} \\
				&=&\displaystyle\lim_{x \to +\infty}  \dfrac{x^{20}\left(2-\dfrac{3}{x}\right)^{20}\cdot x^{30}\left(3+\dfrac{2}{x}\right)^{30}}{x^{50}\left(\dfrac{1}{x}+2\right)^{50}} \\
				&=&\displaystyle\lim_{x \to  +\infty} \dfrac{\left(2-\dfrac{3}{x}\right)^{20}\cdot\left(3+\dfrac{2}{x}\right)^{30}}{\left(\dfrac{1}{x}+2\right)^{50}} \\
				 &=& \dfrac{2^{20}\cdot 3^{30}}{2^{50}}\\
				&=&\left(\dfrac{3}{2}\right)^{30}. 
			\end{eqnarray*}
			\item 
			\begin{eqnarray*}
				J &=& \displaystyle\lim_{x\to +\infty} \dfrac{\sqrt{x^2+x}+2x}{2x+3}\\ 
				&=&\displaystyle\lim_{x\to +\infty} \dfrac{x\left( \sqrt{1+\dfrac{1}{x}}+2\right)}{x\left(2+\dfrac{3}{x}\right)}  \\
				&=&\displaystyle\lim_{x\to +\infty} \dfrac{ \sqrt{1+\dfrac{1}{x}}+2}{2+\dfrac{3}{x}}  \\
				 &=& \dfrac{3}{2}.
			\end{eqnarray*}
			\item 
			\begin{eqnarray*}
				K &=& \displaystyle\lim_{x\to -\infty} \dfrac{\sqrt{x^2+x}+2x}{2x+3}\\ 
				&=&\displaystyle\lim_{x\to -\infty} \dfrac{-x\left( \sqrt{1+\dfrac{1}{x}}-2\right)}{x\left(2+\dfrac{3}{x}\right)} \\ 
				&=&\displaystyle\lim_{x\to +\infty} \dfrac{ -\sqrt{1+\dfrac{1}{x}}+2}{2+\dfrac{3}{x}}  \\
				 &=& \dfrac{1}{2}.
			\end{eqnarray*}
			\item 
			\begin{eqnarray*}
				L  &=& \displaystyle\lim_{x\to -\infty} \dfrac{8x+3}{6x+\sqrt{4x^2+x+3}}\\
				&=&\displaystyle\lim_{x\to -\infty}  \dfrac{x\left(8+\dfrac{3}{x}\right)}{x\left(6- \sqrt{4+\dfrac{1}{x}+\dfrac{3}{x^2}}\right)} \\
				 &=&\displaystyle\lim_{x\to -\infty} \dfrac{8+\dfrac{3}{x}}{6- \sqrt{4+\dfrac{1}{x}+\dfrac{3}{x^2}}} \\
				&=&  2.
			\end{eqnarray*}
		\end{enumerate}
	 }
\end{ex}

% nguồn: Ok_TuanLuan_Toan11_Lan2
\begin{ex}%[1D3H2-3]%[ma:e64910c8]
	Kết quả của  giới hạn $\lim\limits_{x\to 0}\left[x\left(1-\dfrac{1}{x}\right)\right]$ là
	 \loigiai{
		Ta có $\lim\limits_{x\to 0}\left[x\left(1-\dfrac{1}{x}\right)\right]=\lim\limits_{x\to  0}\left(x-1\right)=0-1=-1$.
	} 
\end{ex}

% nguồn: Ok_TuanLuan_Toan11_Lan2
\begin{ex}%[1D3H2-3]%[ma:ab944283]
	 $\lim\limits_{x\to -\infty}\dfrac{\left(x-1\right)\left(x+2\right)}{x^2+9}$ bằng
	\loigiai{ 
		Ta có $\lim\limits_{x\to -\infty}\dfrac{\left(x-1\right)\left(x+2\right)}{x^2+9}=\lim\limits_{x\to  -\infty}\dfrac{\left(1-\dfrac{1}{x}\right)\left(1+\dfrac{2}{x}\right)}{1+\dfrac{9}{x^2}}=1$.
	 }
\end{ex}

% nguồn: Ok_TuanLuan_Toan11_Lan2
\begin{ex}%[1D3H2-3]%[ma:6e14bfd0]
	 Tính giá trị của giới hạn $\lim\limits_{x\to -\infty }\dfrac{2x^2+5x-3}{x^2+6x+3}$.
	\loigiai{ 
	Ta có $\lim\limits_{x\to -\infty  }\dfrac{2x^2+5x-3}{x^2+6x+3}=\lim\limits_{x\to +\infty }\dfrac{2+\dfrac{5}{x}-\dfrac{3}{x^2}}{1+\dfrac{6}{x}+\dfrac{3}{x^2}}=2$.
	} 
\end{ex}

% nguồn: Ok_TuanLuan_Toan11_Lan2
\begin{ex}%[1D3H2-3]%[1D3V2-6]%[ma:1e4bc954]

	Tính các giới  hạn sau đây
	\allowdisplaybreaks 
	\begin{multicols}{2}
		\begin{enumerate}
			\item $\displaystyle\lim_{x \to +\infty}\left(\sqrt{x^2+3x}-\sqrt{x^2+4x} \right)$; 
			\item $\displaystyle\lim_{x \to  +\infty}\left(\sqrt[3]{2x-1}-\sqrt[3]{2x+1} \right)$;
			\item  $\displaystyle\lim_{x \to +\infty}\left[x\left(\sqrt{x^2+5x+4}-\sqrt{x^2+5x-2}\right)\right]$;
			\item  $\displaystyle\lim_{x \to +\infty}\left(\sqrt{x^2-2x}-x\sqrt{2} \right)$;
			 \item $\displaystyle\lim_{x \to +\infty}\left(2x-\sqrt{4x^2+1}\right)$;
			\item  $\displaystyle\lim_{x \to -\infty}\left(\sqrt{x^2+4x}+x\right)$;
			\item $\displaystyle\lim_{x  \to +\infty}\left(\sqrt{x^2+x}-2\sqrt{x^2-x}+x\right)$.
		\end{enumerate}
	\end{multicols}
	 \loigiai{
		\begin{enumerate}
			 \item Ta có
			\begin{eqnarray*}
				 &&\displaystyle\lim_{x \to +\infty} \left( \sqrt{x^2+3x}-\sqrt{x^2+4x} \right)\\
				& = &  \lim_{x \to +\infty} \dfrac{\left(x^2+3x \right)- \left(x^2+4x \right)}{\sqrt{x^2+3x} + \sqrt{x^2+4x}} \\
				& = & \lim_{x  \to +\infty} \dfrac{-x}{x\sqrt{1+\dfrac{3}{x}}+x\sqrt{1+\dfrac{4}{x}}} \\
				& = & \lim_{x \to  +\infty} \dfrac{-1}{\sqrt{1+\dfrac{3}{x}}+\sqrt{1+\dfrac{4}{x}}}\\
				 &=& -\dfrac{1}{2}.
			\end{eqnarray*}
			 \item Ta có
			\begin{eqnarray*}
				&&\displaystyle\lim_{x \to +\infty} \left(  \sqrt[3]{2x - 1} - \sqrt[3]{2x + 1} \right)\\
				& = & \lim_{x \to +\infty} \dfrac{2x - 1 - 2x - 1}{\sqrt[3]{\left(2x -  1\right)^2} + \sqrt[3]{\left(2x - 1\right)\left(2x + 1\right)} + \sqrt[3]{\left(2x + 1\right)^2}} \\
				& = & \lim_{x \to +\infty} \dfrac{-2}{\sqrt[3]{\left(2x - 1\right)^2} + \sqrt[3]{\left(2x - 1\right)\left(2x +  1\right)} + \sqrt[3]{\left(2x + 1\right)^2}} \\
				& = & \lim_{x \to +\infty} \dfrac{-2}{\sqrt[3]{x^2 \left(2 - \dfrac{1}{x} \right)^2} + \sqrt[3]{x^2 \left(4 - \dfrac{1}{x^2} \right)} + \sqrt[3]{x^2 \left(2 + \dfrac{1}{x}  \right)^2}} \\
				 & = & \lim_{x \to +\infty} \dfrac{-2}{x^2 \left[\sqrt[3]{\left(2 - \dfrac{1}{x} \right)^2} + \sqrt[3]{4 - \dfrac{1}{x^2}} + \sqrt[3]{\left(2 + \dfrac{1}{x} \right)^2} \right]}\\
				 &=& 0.
			\end{eqnarray*}
			\item Ta  có
			\begin{eqnarray*}
				&&\displaystyle\lim_{x \to +\infty}\left[x\left(  \sqrt{x^2+5x+4}-\sqrt{x^2+5x-2}\right)\right]\\
				& = & \displaystyle\lim_{x \to +\infty} \dfrac{x \left[x^2 + 5x  + 4 - \left( x^2 + 5x - 2 \right) \right]}{\sqrt{x^2 + 5x + 4} + \sqrt{x^2 + 5x - 2}} \\
				& = & \displaystyle\lim_{x \to +\infty} \dfrac{6x}{\sqrt{x^2 + 5x + 4} + \sqrt{x^2 + 5x - 2}} \\ 
				& = & \displaystyle\lim_{x  \to +\infty} \dfrac{6x}{x \sqrt{1 + \dfrac{5}{x} + \dfrac{4}{x^2}} + x \sqrt{1 + \dfrac{5}{x} - \dfrac{2}{x^2}}} \\
				& = & \displaystyle\lim_{x \to +\infty} \dfrac{6}{\sqrt{1  + \dfrac{5}{x} + \dfrac{4}{x^2}} + \sqrt{1 + \dfrac{5}{x} - \dfrac{2}{x^2}}} \\
				& = & 3. 
			\end{eqnarray*}
			 \item Ta có
			\begin{eqnarray*}
				 &&\displaystyle\lim_{x \to +\infty} \left(\sqrt{x^2-2x}-x\sqrt{2} \right)\\
				 &=& \displaystyle\lim_{x \to +\infty} \left(x\sqrt{1-\dfrac{2}{x}} - x \sqrt{2}\right) \\
				&=& \displaystyle\lim_{x \to  +\infty} x\left(\sqrt{1-\dfrac{2}{x}} - \sqrt{2} \right) \\
				&=& -\infty. 
			\end{eqnarray*}
			Vì $\displaystyle\lim_{x  \to +\infty} \left(\sqrt{1-\dfrac{2}{x}}-\sqrt{2}\right)=1-\sqrt{2}<0$ và $\displaystyle\lim_{x \to +\infty} x=+\infty$.
			\item Ta  có
			\begin{eqnarray*}
				 &&\displaystyle\lim_{x \to +\infty} \left(2x-\sqrt{4x^2+1}\right) \\
				&=&\displaystyle\lim_{x  \to +\infty} \dfrac{4x^2-4x^2-1}{2x+\sqrt{4x^2+1}}\\
				&=&\displaystyle\lim_{x \to +\infty} \dfrac{-1}{2x+\sqrt{4x^2+1}}\\ 
				&=&\displaystyle\lim_{x \to +\infty}  \dfrac{-1}{x\left(2+\sqrt{4+\dfrac{1}{x^2}}\right)} \\
				&=&  0.
			\end{eqnarray*}
			\item  Ta có
			\begin{eqnarray*}
				&&\displaystyle\lim_{x\to -\infty} \left(\sqrt{x^2+4x}+x\right)\\ 
				&=&\displaystyle\lim_{x\to -\infty}  \dfrac{x^2+4x-x^2}{\sqrt{x^2+4x}-x}\\
				 &=&\displaystyle\lim_{x\to -\infty} \dfrac{4x}{-x\left( \sqrt{1+\dfrac{4}{x}}+1\right)}\\
				&=&\displaystyle\lim_{x\to  -\infty} \dfrac{4}{-\left( \sqrt{1+\dfrac{4}{x}}+1\right)}\\
				&=&-2. 
			\end{eqnarray*}
			 \item Ta có
			\begin{eqnarray*}
				&&\displaystyle\lim_{x \to +\infty}\left( \sqrt{x^2+x}-2\sqrt{x^2-x}+x\right)\\ 
				&=&\displaystyle\lim_{x \to +\infty}\left(  \sqrt{x^2+x}-\sqrt{x^2-x}+x-\sqrt{x^2-x}\right)\\
				 &=&\displaystyle\lim_{x \to +\infty}\left( \dfrac{2x}{\sqrt{x^2+x}+\sqrt{x^2-x}}+\dfrac{x}{x+\sqrt{x^2-x}}\right)\\
				 &=&\displaystyle\lim_{x \to +\infty}\left( \dfrac{2}{\sqrt{1+\dfrac{1}{x}}+\sqrt{1-\dfrac{1}{x}}}+\dfrac{1}{1+\sqrt{1-\dfrac{1}{x}}}\right)\\
				 &=&\dfrac{5}{2}.
			\end{eqnarray*}
		\end{enumerate}
	 }
\end{ex}

% nguồn: Ok_TuanLuan_Toan11_Lan2
\begin{ex}%[1D3H2-3]%[ma:a488e0eb]
	 Giới hạn $\lim\limits_{x\to +\infty}\dfrac{c{x^2}+a}{x^2+b}$ bằng
	 \loigiai{
		 Ta có $\lim\limits_{x\to +\infty}\dfrac{c{x^2}+a}{x^2+b}=\lim\limits_{x\to +\infty}\dfrac{c+\dfrac{a}{x^2}}{1+\dfrac{b}{x^2}}=\dfrac{c+0}{1+0}=c$.
	} 
\end{ex}

% nguồn: Ok_TuanLuan_Toan11_Lan2
\begin{ex}%[1D3H2-3]%[ma:a1abc144]
	Kết quả của giới hạn $\lim\limits_{x\to 0}x^2\left(\sin\pi x-\dfrac{1}{x^2}\right)$ là 
	 \loigiai{
		Ta có $\lim\limits_{x\to  0}x^2\left(\sin\pi x-\dfrac{1}{x^2}\right)=\lim\limits_{x\to 0}\left(x^2\sin\pi x-1\right)=-1$.
	 }
\end{ex}

% nguồn: Ok_TuanLuan_Toan11_Lan2
\begin{ex}%[1D3H2-3]%[ma:3e451a1a]
	Giá trị của $\lim\limits_{x\to  -\infty}\dfrac{2x-1}{\sqrt{x^2+1}-1}$ bằng
	 \loigiai{
		Ta có $\lim\limits_{x\to -\infty}\dfrac{2x-1}{\sqrt{x^2+1}-1}=\lim\limits_{x\to -\infty}\dfrac{2x-1}{-x\sqrt{1+\dfrac{1}{x^2}}-1}=\lim\limits_{x\to -\infty}\dfrac{2-\dfrac{1}{x}}{-\sqrt{1+\dfrac{1}{x^2}}-\dfrac{1}{x}}=-2$. 
	} 
\end{ex}

% nguồn: Ok_TuanLuan_Toan11_Lan2
\begin{ex}%[1D3H2-3]%[ma:e645bc5d]
	Tính $\lim\limits_ {x\to \sqrt{3}}\dfrac{2x^2-6}{x-\sqrt{3}}=a\sqrt{b}$ ($a$, $b$  nguyên). Khi đó giá trị của $P=a+b$ bằng bao nhiêu?
	\loigiai{ 
	Ta  có $\lim\limits_{x\to \sqrt{3}}\dfrac{2x^2-6}{x-\sqrt{3}}=\lim\limits_{x\to \sqrt{3}}\dfrac{2\left(x^2-3\right)}{x-\sqrt{3}}=\lim\limits_{x\to \sqrt{3}} 2\left(x+\sqrt{3} \right)=4\sqrt{3}$.\\
	Suy ra  $a=4$, $b=3$.\\
	Vậy $P=a+b=7$. 
	} 
\end{ex}

% nguồn: Ok_TuanLuan_Toan11_Lan2
\begin{ex}%[1D3V2-3]%[ma:4e93e06e]
	Tính giá trị của cặp số $(a,b)$ thỏa mãn $ \lim\limits_{x\to 3}\,\dfrac{x^2+ax+b}{x-3}=3$. 
	\loigiai{ 
	\begin{itemize}
		\item Cách 1:\\ 
		Để $ \lim\limits_{x\to  3}\,\dfrac{x^2+ax+b}{x-3}=3$ thì ta phải có $x^2+ax+b=(x-3)(x-m)$.\\
		Khi đó  $3-m=3\Leftrightarrow m=0$.\\
		Vậy $x^2+ax+b=(x-3)x=x^2-3x$.\\ 
		Suy ra  $a=-3$ và $b=0$.
		\item Cách  2:\\
		Ta  có $\dfrac{x^2+ax+b}{x-3}=x+a+3+\dfrac{3a+b+9}{x-3}$.\\
		Vậy để có $ \lim\limits_{x\to 3}\,\dfrac{x^2+ax+b}{x-3}=3$  thì ta phải có $\heva{&3a+b+9=0 \\ &a+6=3}\Leftrightarrow\heva{&a=-3 \\ &b=0.}$
	\end{itemize}
	 }
\end{ex}

% nguồn: Ok_TuanLuan_Toan11_Lan2
\begin{ex}%[1D3V2-3]%[ma:b064a26e]
	Cho $a$, $b$ là các số thực  khác $0$. Để giới hạn $\lim\limits_{x\to -\infty}\dfrac{\sqrt{x^2-3x}+ax}{bx-1}=3$ thì
	 \loigiai{
		 Ta có
		 \allowdisplaybreaks
		\begin{eqnarray*}
			\lim\limits_{x\to -\infty}\dfrac{\sqrt{x^2-3x}+ax}{bx-1}&=&\lim\limits_{x\to -\infty}\dfrac{-x\sqrt{1-\dfrac{3}{x}}+ax}{bx-1}\\ 
			&=&\lim\limits_{x\to -\infty}\dfrac{x\left(-\sqrt{1-\dfrac{3}{x}}+a\right)}{x\left(b-\dfrac{1}{x}\right)}\\ 
			 &=&\lim\limits_{x\to -\infty}\dfrac{-\sqrt{1-\dfrac{3}{x}}+a}{b-\dfrac{1}{x}}\\
			&=&\dfrac{a-1}{b}. 
		\end{eqnarray*}
		Theo đề bài ta có $\dfrac{a-1}{b}=3\Leftrightarrow a-1=3b  \Leftrightarrow a=3b+1$.
	} 
\end{ex}

% nguồn: Ok_TuanLuan_Toan11_Lan2
\begin{ex}%[1D3V2-3]%[ma:fb7be103]
	Cho số thực $a$ thỏa mãn $\lim\limits_{x\to +\infty}\dfrac{a\sqrt{2x^2+3}+2017}{2x+2018}=\dfrac{1}{2}$. Khi đó giá trị  của $a$ là
	 \loigiai{
		Ta có $\lim\limits_{x\to +\infty}\dfrac{a\sqrt{2x^2+3}+2017}{2x+2018}=\dfrac{1}{2}$ $\Leftrightarrow\lim\limits_{x\to +\infty}\dfrac{a\sqrt{2+\dfrac{3}{x^2}}+\dfrac{2017}{x}}{2+\dfrac{2018}{x}}=\dfrac{1}{2}\Leftrightarrow\dfrac{a\sqrt{2}}{2}=\dfrac{1}{2}\Leftrightarrow a=\dfrac{\sqrt{2}}{2}$. 
	} 
\end{ex}

% nguồn: Ok_TuanLuan_Toan11_Lan2
\begin{ex}%[1D3V2-3]%[ma:6e527eea]
	Biết $\lim\limits_{x\to +\infty}\dfrac{\left(2-a\right)x-3}{x-\sqrt{x^2+1}}=+\infty$. Giá trị nhỏ nhất của $P=a^2-2a+4$ là 
	\loigiai{ 
		Ta có $\lim\limits_{x\to +\infty}\dfrac{\left(2-a\right)x-3}{x-\sqrt{x^2+1}}=\lim\limits_{x\to +\infty}\left[-\left(\left(2-a\right)x-3\right)\left(x+\sqrt{x^2+1}\right)\right]=+\infty$.\\ 
		Suy ra  $-\left(2-a\right)\ge 0\Leftrightarrow a\ge 2$.\\
		Với $a\ge 2\Rightarrow a\left(a-2\right)\ge 0$ suy ra $P=a\left(a-2\right)+4\ge 4$. 
	} 
\end{ex}

% nguồn: Ok_TuanLuan_Toan11_Lan2
\begin{ex}%[1D3V2-3]%[ma:bdee4b1d]
	Tính giá trị của giới hạn $\lim\limits_{x\to 2^{2018}}\dfrac{x^2-4^{2018}}{x-2^{2018}}$. 
	 \loigiai{
	Ta có $ \lim\limits_{x\to 2^{2018}}\,\dfrac{x^2-4^{2018}}{x-2^{2018}}=\lim\limits_{x\to 2^{2018}}\,\dfrac{\left( x-2^{2018} \right)\left(  x+2^{2018} \right)}{x-2^{2018}}=\lim\limits_{x\to 2^{2018}}\,\left( x+2^{2018} \right)=2^{2018}+2^{2018}=2^{2019}$.
	} 
\end{ex}

% nguồn: Ok_TuanLuan_Toan11_Lan2
\begin{ex}%[1D3V2-3]%[ma:daf47ad7]
	Để  $\lim\limits_{x\to -\infty}\dfrac{\sqrt{4x^2+x+1}+4}{mx-2}=\dfrac{1}{2}$. Tìm giá trị của $m$.
	\loigiai{ 
		Ta có $\lim\limits_{x\to -\infty}\dfrac{\sqrt{4x^2+x+1}+4}{mx-2}=\lim\limits_{x\to  -\infty}\dfrac{-\sqrt{4+\dfrac{1}{x}+\dfrac{1}{x^2}}+\dfrac{4}{x}}{m-\dfrac{2}{x}}=-\dfrac{2}{m}$.\\
		Theo bài  ra ta có $-\dfrac{2}{m}=\dfrac{1}{2}\Leftrightarrow m=-4$.
	} 
\end{ex}

% nguồn: Ok_TuanLuan_Toan11_Lan2
\begin{ex}%[1D3V2-3]%[ma:c4e19232]
	Biết $ \lim\limits_{x\to 1}\dfrac{x^3-ax+a-1}{x-1}=2$. Tính $M=a^2+2a$. 
	\loigiai{ 
	 Ta có
	\begin{align*}
		 \lim\limits_{x\to 1}\dfrac{x^3-ax+a-1}{x-1}&=\lim\limits_{x\to 1}\dfrac{(x-1)\left(x^2+x+1 \right)-a(x-1)}{x-1}\\
		 &= \lim\limits_{x\to 1}\,\left(x^2+x+1-a \right)\\
		&=3-a. 
	\end{align*}
	Suy ra  $3-a=2 \Leftrightarrow a=1$.\\
	Vậy $M=a^2+2a=3$. 
	} 
\end{ex}

% nguồn: Ok_TuanLuan_Toan11_Lan2
\begin{ex}%[1D3V2-3]%[ma:2cf87d92]
	Cho $a$, $b$ là số nguyên và $  \lim\limits_{x\to 1}\dfrac{ax^2+bx-5}{x-1}=7$. Tính $a^2+b^2+a+b$.
	\loigiai{ 
	 Vì $ \lim\limits_{x\to 1}\,\dfrac{ax^2+bx-5}{x-1}=7$ hữu hạn nên $x=1$ phải là nghiệm của phương trình $ax^2+bx-5=0$ suy ra $a+b-5=0\Rightarrow b=5-a$.\\
	Khi đó $ \lim\limits_{x\to  1}\,\dfrac{ax^2+(5-a)x-5}{x-1}=\lim\limits_{x\to 1}\dfrac{(x-1)(ax+5)}{x-1}=a+5=7\Rightarrow a=2$ nên $b=3$.\\
	 Suy ra $a^2+b^2+a+b=18$.
	 }
\end{ex}

% nguồn: Ok_TuanLuan_Toan11_Lan1
\begin{bt}%[1D3Y2-3]%[ma:21dda78a]
	Cho dãy số $( u_n)$  với $u_n=\sin n$. Xét tính tăng giảm của dãy số.
	 \loigiai{
	Ta có $\sin 1 < \sin 2$ nhưng $\sin 2 > \sin 3$ nên $u_n$ không tăng không giảm.  }
\end{bt}

% nguồn: Ok_TuanLuan_Toan11_Lan1
\begin{bt}%[1D3Y2-3]%[ma:781965a3]
	Cho dãy  số $( u_n)$ với $u_n=\dfrac{1}{n+1}$. Xét tính tăng giảm của dãy số.
	\loigiai{ 
	Ta có ${u}_{n+1}-{u}_{n}=\dfrac{1}{n+2}-\dfrac{1}{n+1}=\dfrac{n+1-n-2}{\left( n+2 \right)\left( n+1 \right)}=\dfrac{-1}{\left(  n+2 \right)\left( n+1 \right)}<0,\forall n\in {{\mathbb{N}}^{*}}$\\ Vậy ${u}_{n}$ là dãy giảm}
\end{bt}

% nguồn: Ok_TuanLuan_Toan11_Lan1
\begin{bt}%[1D3Y2-3]%[ma:a8175aca]
	Cho dãy số $( u_n)$ với $u_n={{n}^2}+1$.  Xét tính tăng giảm của dãy số.
	\loigiai{ 
	Ta có $u_{n+1}-u_{n}=\left(n+1\right)^2+1-n^2-1=2n+1>0, \forall n\in {{\mathbb{N}}^{*}}$.\\  Vậy $u_n$ là dãy tăng.}
\end{bt}

% nguồn: Ok_TuanLuan_Toan11_Lan2
\begin{ex}%[1D3Y2-3]%[ma:f60be9b1]
	Cho dãy số $( u_n)$ với $u_n=\sin n$. Xét  tính tăng giảm của dãy số.
	 \loigiai{
	Ta có $\sin 1 < \sin 2$ nhưng $\sin 2  > \sin 3$ nên $u_n$ không tăng không giảm. }
\end{ex}

% nguồn: Ok_TuanLuan_Toan11_Lan2
\begin{ex}%[1D3Y2-3]%[ma:dd9df834]
	Cho dãy số $( u_n)$ với $u_n=\dfrac{1}{n+1}$. Xét tính tăng  giảm của dãy số.
	 \loigiai{
	Ta có ${u}_{n+1}-{u}_{n}=\dfrac{1}{n+2}-\dfrac{1}{n+1}=\dfrac{n+1-n-2}{\left( n+2 \right)\left( n+1 \right)}=\dfrac{-1}{\left( n+2 \right)\left( n+1  \right)}<0,\forall n\in {{\mathbb{N}}^{*}}$\\ Vậy ${u}_{n}$ là dãy giảm}
\end{ex}

% nguồn: Ok_TuanLuan_Toan11_Lan2
\begin{ex}%[1D3Y2-3]%[ma:f1cb87e9]
	Cho dãy số $( u_n)$ với $u_n={{n}^2}+1$. Xét tính tăng giảm  của dãy số.
	\loigiai{ 
	Ta có  $u_{n+1}-u_{n}=\left(n+1\right)^2+1-n^2-1=2n+1>0, \forall n\in {{\mathbb{N}}^{*}}$.\\ Vậy $u_n$ là dãy tăng.}
\end{ex}
